% Đặc tả chi tiết controller Portal WujiaTea — build:
%   cd docs && lualatex controller-spec-detail.tex && lualatex controller-spec-detail.tex
% (KHÔNG qua scripts/build-doc.sh — đây là tài liệu độc lập, không thuộc master 300 trang)
\documentclass[11pt,a4paper]{report}

\usepackage{fontspec}
\setmainfont{Noto Sans}
\setsansfont{Noto Sans}
\setmonofont{DejaVu Sans Mono}[Scale=0.82]

\usepackage[margin=2cm]{geometry}
\usepackage{parskip}
\setlength{\parindent}{0pt}

\usepackage{booktabs}
\usepackage{longtable}
\usepackage{array}
\usepackage{tabularx}
\usepackage{enumitem}
\usepackage{xcolor}
\usepackage{colortbl}
\usepackage{mdframed}
\newcolumntype{L}[1]{>{\raggedright\arraybackslash}p{#1}}

\definecolor{hdrbg}{RGB}{238,242,248}
\definecolor{askbg}{RGB}{255,247,230}
\definecolor{askln}{RGB}{217,119,6}
\definecolor{okc}{RGB}{21,128,61}
\definecolor{partc}{RGB}{180,110,10}
\definecolor{noc}{RGB}{185,28,28}
\definecolor{greyc}{RGB}{110,110,110}

% --- ký hiệu dùng chung -------------------------------------------------------
% \rt{} dùng \nolinkurl: gõ thẳng is_public_portal / <int:id> không cần escape,
% và đường dẫn dài tự ngắt dòng ở dấu / _ thay vì tràn ra lề (bẫy build đầu tiên).
\newcommand{\rt}[1]{{\footnotesize\nolinkurl{#1}}}
\newcommand{\uat}[1]{\textcolor{okc}{\textbf{#1}}}         % số đo thật trên UAT
\newcommand{\ok}{\textcolor{okc}{\textbf{ĐÃ CÓ}}}
\newcommand{\pt}{\textcolor{partc}{\textbf{MỘT PHẦN}}}
\newcommand{\no}{\textcolor{noc}{\textbf{CHƯA CÓ}}}
\newcommand{\note}[1]{\textcolor{greyc}{\small #1}}

% --- khuôn "phiếu màn": 7 mục cố định, không màn nào được thiếu ---------------
\newcommand{\screen}[2]{\section{#1}\nopagebreak\par\rt{#2}\par\medskip}
\newcommand{\mvao}{\subsection*{1. Vào màn bằng đâu}}
\newcommand{\mai}{\subsection*{2. Ai xem được gì}}
\newcommand{\mhien}{\subsection*{3. Màn hiện gì --- lấy ở đâu, lọc điều kiện gì}}
\newcommand{\mloc}{\subsection*{4. Người dùng lọc / nhập được gì}}
\newcommand{\mkiem}[1]{\subsection*{5. Bấm \emph{#1} thì hệ thống kiểm gì}}
\newcommand{\mghi}{\subsection*{6. Ghi gì vào hệ thống}}
\newcommand{\mhoi}{\subsection*{7. Điểm cần BA xác nhận}}

% khối câu hỏi cho BA
\newenvironment{ask}
  {\begin{mdframed}[backgroundcolor=askbg,linecolor=askln,linewidth=0.8pt,
      innertopmargin=6pt,innerbottommargin=6pt,skipabove=6pt,skipbelow=6pt]}
  {\end{mdframed}}

% bảng "màn hiện gì"
\newenvironment{hienbang}
  {\par\footnotesize\begin{longtable}{L{2.4cm}L{2.7cm}L{5.6cm}L{1.9cm}L{1.9cm}}
   \rowcolor{hdrbg}\textbf{Khối trên màn} & \textbf{Lấy từ bảng} &
   \textbf{Điều kiện lọc} & \textbf{Sắp xếp / số dòng} & \textbf{Thực tế UAT}\\ \midrule
   \endfirsthead
   \rowcolor{hdrbg}\textbf{Khối trên màn} & \textbf{Lấy từ bảng} &
   \textbf{Điều kiện lọc} & \textbf{Sắp xếp / số dòng} & \textbf{Thực tế UAT}\\ \midrule
   \endhead}
  {\end{longtable}\par}

% bảng "kiểm gì theo thứ tự"
\newenvironment{kiembang}
  {\par\footnotesize\begin{longtable}{L{0.5cm}L{7.0cm}L{7.2cm}}
   \rowcolor{hdrbg}\textbf{\#} & \textbf{Điều kiện phải đạt} &
   \textbf{Không đạt thì người dùng thấy gì}\\ \midrule
   \endfirsthead
   \rowcolor{hdrbg}\textbf{\#} & \textbf{Điều kiện phải đạt} &
   \textbf{Không đạt thì người dùng thấy gì}\\ \midrule
   \endhead}
  {\end{longtable}\par}

\usepackage{url}
\usepackage{hyperref}
\hypersetup{
    colorlinks=true, linkcolor=blue!50!black, urlcolor=blue!60!black,
    pdftitle={WujiaTea — Đặc tả chi tiết controller Portal},
    pdfauthor={WujiaTea Dev},
}

\title{\textbf{WujiaTea --- Đặc tả chi tiết Controller Portal}\\[4pt]
\large Mỗi màn: hiện gì theo điều kiện nào · ai xem được gì\\[2pt]
\large bấm nút thì kiểm gì · ghi gì vào hệ thống · thực tế trên UAT đang ra sao}
\author{Dev · Odoo 19}
\date{20/09/2026}

% số trang 3 chữ số trong mục lục cần ô rộng hơn mặc định
\makeatletter
\renewcommand{\@pnumwidth}{2.2em}
\renewcommand{\@tocrmarg}{3.2em}
\makeatother

\sloppy
\begin{document}
\maketitle
\tableofcontents

\chapter{Đọc tài liệu này thế nào}

\section{Tài liệu này khác bản trước ở chỗ nào}

Bản kiểm kê 20 trang giao ngày 19/09 trả lời câu \emph{``portal đã có những controller nào''}.
Tài liệu này trả lời câu tiếp theo, cụ thể hơn: \textbf{mỗi màn đang chạy luật gì}.

Với mỗi màn, tài liệu nói đủ bảy điều --- luôn đúng thứ tự này, không màn nào thiếu:

\begin{enumerate}[leftmargin=*,itemsep=1pt]
  \item \textbf{Vào màn bằng đâu} --- đường dẫn, vào từ menu nào, kèm những đường dẫn gọi ngầm.
  \item \textbf{Ai xem được gì} --- cần đăng nhập/chọn cửa hàng chưa, vai trò nào vào được, thấy
        dữ liệu của cửa hàng nào.
  \item \textbf{Màn hiện gì} --- từng khối lấy từ bảng nào, \textbf{lọc theo điều kiện gì}, sắp
        xếp ra sao, mỗi trang bao nhiêu dòng, và \textbf{thực tế trên UAT đang có bao nhiêu bản
        ghi} thoả điều kiện đó.
  \item \textbf{Người dùng lọc / nhập được gì} --- ô tìm kiếm tìm theo trường nào, giá trị lạ thì
        hệ thống xử lý sao.
  \item \textbf{Bấm nút thì kiểm gì} --- liệt kê \textbf{đúng thứ tự} các điều kiện hệ thống
        kiểm, và không đạt thì người dùng thấy câu gì.
  \item \textbf{Ghi gì vào hệ thống} --- tạo/sửa bản ghi nào, trạng thái nào, ai đứng tên.
  \item \textbf{Điểm cần BA xác nhận} --- chỗ hệ thống đang làm theo một giả định, chỗ dữ liệu
        thật đang lệch cấu hình, chỗ cần anh chốt trước khi lên task.
\end{enumerate}

Mục 7 in trên \textbf{nền vàng} để dễ tìm; toàn bộ mục 7 của mọi màn được gom lại ở chương cuối.

\section{Quy ước ký hiệu}

\begin{itemize}[leftmargin=*]
  \item \rt{chữ máy} --- tên bảng dữ liệu, tên trường, đường dẫn, đúng như trong hệ thống. Viết
        ra để anh đối chiếu được với tab \emph{1. Model/ Field} và để Dev tìm đúng chỗ.
  \item \uat{Số xanh} --- số bản ghi \textbf{thật} đang có trên máy chủ UAT tại thời điểm chụp.
  \item ``Cửa hàng'' = một bản ghi \rt{wujia.franchise.management}; ``thành viên'' =
        \rt{wujia.franchise.member} (một tài khoản gắn vào một cửa hàng với một vai trò).
  \item Vai trò trong portal chỉ có ba mức: \textbf{Chủ cửa hàng} > \textbf{Quản lý} >
        \textbf{Nhân viên}.
\end{itemize}

\section{Số liệu lấy ở đâu, lúc nào}

\begin{itemize}[leftmargin=*]
  \item \textbf{Luật nghiệp vụ}: đọc trực tiếp mã nguồn nhánh \rt{main} --- 23 tệp controller
        (\textbf{7.347 dòng}) cộng phần mã nghiệp vụ mà chúng gọi tới. Phần liệt kê đường dẫn,
        điều kiện lọc, chuỗi kiểm được \textbf{máy trích bằng cách đọc cây cú pháp}
        (\rt{scripts/qa/controller_spec.py}), không dùng tìm chuỗi --- tìm chuỗi đếm nhầm cả
        chú thích. Phần diễn giải nghiệp vụ do người đọc mã viết; máy không đoán được ý nghiệp vụ.
  \item \textbf{Số thực tế}: đọc \textbf{chỉ-đọc} máy chủ UAT \rt{113.161.187.126:8019}, cơ sở
        dữ liệu \rt{wujia_tea_19}, bằng \rt{scripts/qa/uat_domain_probe.py}. Script chặn cứng
        mọi lệnh ghi --- phiên làm tài liệu không đụng một bản ghi nào của anh.
  \item \textbf{Ngày chụp}: 20/09/2026.
\end{itemize}

\begin{ask}
\textbf{Lưu ý quan trọng khi đọc các con số:} UAT đang là môi trường \textbf{dữ liệu mẫu mỏng}
--- 8 sản phẩm, 4 cửa hàng, 7 tài khoản, 47 đơn hàng. Vì vậy con số trong tài liệu dùng để
\textbf{soi cấu hình} (\emph{bao nhiêu sản phẩm đang bật cờ hiện trên portal}, \emph{khu vực nào
chưa cấu hình khung giờ}), \textbf{không} dùng để suy ra tải thật khi lên production.
\end{ask}

\section{Bản đồ chương --- tìm nhanh}

Tài liệu mô tả đủ \textbf{104/104 đường dẫn} của \textbf{16 mô-đun}, gom thành \textbf{67 phiếu
màn}. Nếu chỉ cần tra một phân hệ, vào thẳng chương của nó:

{\footnotesize
\begin{center}
\begin{tabular}{@{}L{1.3cm}L{6.4cm}L{7.3cm}@{}}
\toprule
\textbf{Chương} & \textbf{Phân hệ} & \textbf{Gồm}\\
\midrule
2 & Bản đồ 104 đường dẫn & Bảng tra: đường dẫn nào thuộc mô-đun nào, làm gì. \\
3 & Đặt hàng & Sản phẩm, danh mục, giỏ hàng, gửi đơn, khung giờ đặt hàng. \\
4 & Khung portal · Trang chủ & Đăng nhập, chọn cửa hàng, hồ sơ, Trang chủ, thông tin cửa hàng, \textbf{quy ước dùng chung}. \\
5 & Lịch sử mua · Giao hàng · Báo cáo & Đơn đã đặt, chuyến giao, báo cáo đặt hàng và tệp Excel. \\
6 & Công nợ · Đổi trả · Bù hàng & Công nợ theo tuần, lịch sử thanh toán, chuyển khoản, yêu cầu đổi trả. \\
7 & Hỗ trợ · Yêu cầu thông tin · Kiến thức · Thông báo & Bốn phân hệ nhẹ, dùng chung nhiều quy ước. \\
8 & Đăng ký thi & Lịch thi, khung giờ, gửi phiếu, kết quả. \\
9 & Khảo sát · Metabase & Phía cửa hàng và phía người đi khảo sát; nhúng bảng số liệu. \\
10 & \textbf{Tổng hợp} & Đối chiếu CT-001…CT-071, \textbf{157 điểm cần anh chốt}, việc BA không cần lên task. \\
\bottomrule
\end{tabular}
\end{center}}

\textbf{Nếu chỉ có 15 phút}: đọc chương 10 trước --- ba bảng đầu của chương đó đã đủ để lên
task. Các chương 3--9 là phần tra cứu khi cần biết chi tiết một màn.

\section{Ba điều nên biết trước khi lên task}

\begin{enumerate}[leftmargin=*]
  \item \textbf{Đường dẫn sẽ không đổi.} Các đợt chuẩn hoá sắp tới có dời mã nguồn của một số
        màn sang mô-đun khác, nhưng \textbf{đường dẫn và hành vi giữ nguyên}. Anh không cần lên
        task ``viết lại'' những controller đó.
  \item \textbf{Có màn chưa có dòng mã nào.} Chương cuối liệt kê các mục \rt{CT-0xx} trong tab
        \emph{3. Controller} chưa có code --- đó mới là chỗ cần task mới.
  \item \textbf{Có mã đang chạy mà tab CT chưa mô tả.} Cũng liệt kê ở chương cuối, để anh bổ
        sung spec cho khớp thực tế.
\end{enumerate}

\chapter{Bản đồ đường dẫn}

Portal hiện có \textbf{104 đường dẫn} trong \textbf{16 mô-đun}. Bảng này là mục lục tra cứu: mỗi đường dẫn nằm ở chương nào trong tài liệu. Cột \emph{Kiểu} phân biệt \emph{Trang} (người dùng nhìn thấy, gõ được trên thanh địa chỉ) với \emph{Gọi ngầm} (trình duyệt gọi khi bấm nút, không phải màn riêng).

\section{Chương 3 --- Đặt hàng}
\footnotesize
\begin{longtable}{L{7.0cm}L{1.8cm}L{6.0cm}}
\rowcolor{hdrbg}\textbf{Đường dẫn} & \textbf{Kiểu} & \textbf{Làm gì}\\ \midrule
\endfirsthead
\rowcolor{hdrbg}\textbf{Đường dẫn} & \textbf{Kiểu} & \textbf{Làm gì}\\ \midrule
\endhead
\rt{/portal/order} & Trang & Danh sách sản phẩm đặt hàng\\
\rt{/portal/order/cart} & Trang & Màn giỏ hàng của cửa hàng\\
\rt{/portal/order/cart/add} & Gọi ngầm & Thêm sản phẩm vào giỏ\\
\rt{/portal/order/cart/count} & Gọi ngầm & Số hiển thị trên biểu tượng giỏ\\
\rt{/portal/order/cart/fragment} & Gọi ngầm & Lấy lại toàn bộ giỏ khi người cùng cửa hàng vừa sửa\\
\rt{/portal/order/cart/note} & Gọi ngầm & Lưu ghi chú đơn hàng (dùng chung cả cửa hàng)\\
\rt{/portal/order/cart/remove} & Gọi ngầm & Xoá một dòng khỏi giỏ\\
\rt{/portal/order/cart/step} & Gọi ngầm & Bấm + / − tăng giảm đúng một bước\\
\rt{/portal/order/cart/update} & Gọi ngầm & Gõ thẳng số lượng cho một dòng\\
\rt{/portal/order/product/<int:product_id>} & Trang & Chi tiết một sản phẩm\\
\rt{/portal/order/rejected} & Trang & Màn Không thể gửi đơn — ngoài khung giờ\\
\rt{/portal/order/results} & Trang & Phần kết quả khi lọc / tìm kiếm (không tải lại trang)\\
\rt{/portal/order/submit} & Trang & Gửi đơn — tạo đơn hàng nháp trong ERP\\
\rt{/portal/order/submitted/<int:order_id>} & Trang & Màn Đặt hàng thành công (điện thoại)\\
\end{longtable}
\normalsize

\section{Chương 4 --- Khung portal + Trang chủ}
\footnotesize
\begin{longtable}{L{7.0cm}L{1.8cm}L{6.0cm}}
\rowcolor{hdrbg}\textbf{Đường dẫn} & \textbf{Kiểu} & \textbf{Làm gì}\\ \midrule
\endfirsthead
\rowcolor{hdrbg}\textbf{Đường dẫn} & \textbf{Kiểu} & \textbf{Làm gì}\\ \midrule
\endhead
\rt{/my/branches} & Trang & ---\\
\rt{/my/branches/<int:branch_id>} & Trang & ---\\
\rt{/my/franchises} & Trang & ---\\
\rt{/my/franchises/<int:franchise_id>} & Trang & ---\\
\rt{/my/franchises/<int:franchise_id>/members} & jsonrpc & ---\\
\rt{/portal} & Trang & ---\\
\rt{/portal/_pc-preview} & Trang & ---\\
\rt{/portal/branches} & Trang & ---\\
\rt{/portal/branches/<int:branch_id>} & Trang & ---\\
\rt{/portal/change-password} & Trang & ---\\
\rt{/portal/forgot-pass} & Trang & ---\\
\rt{/portal/franchise-information} & Trang & ---\\
\rt{/portal/franchise/clear} & Trang & ---\\
\rt{/portal/franchise/switch} & Trang & ---\\
\rt{/portal/franchises} & Trang & ---\\
\rt{/portal/franchises/<int:franchise_id>} & Trang & ---\\
\rt{/portal/franchises/<int:franchise_id>/profile} & Trang & ---\\
\rt{/portal/login} & Trang & ---\\
\rt{/portal/logout} & Trang & ---\\
\rt{/portal/profile} & Trang & ---\\
\rt{/portal/profile/avatar/<int:user_id>} & Trang & ---\\
\rt{/portal/profile/update} & Trang & ---\\
\rt{/portal/redirect} & Trang & ---\\
\rt{/portal/reset_password} & Trang & ---\\
\rt{/portal/set-lang/<string:lang_code>} & Trang & ---\\
\end{longtable}
\normalsize

\section{Chương 5 --- Lịch sử mua · Giao hàng · Báo cáo}
\footnotesize
\begin{longtable}{L{7.0cm}L{1.8cm}L{6.0cm}}
\rowcolor{hdrbg}\textbf{Đường dẫn} & \textbf{Kiểu} & \textbf{Làm gì}\\ \midrule
\endfirsthead
\rowcolor{hdrbg}\textbf{Đường dẫn} & \textbf{Kiểu} & \textbf{Làm gì}\\ \midrule
\endhead
\rt{/portal/delivery} & Trang & ---\\
\rt{/portal/delivery/<int:batch_id>} & Trang & ---\\
\rt{/portal/delivery/<int:batch_id>.ics} & Trang & ---\\
\rt{/portal/delivery/results} & Trang & ---\\
\rt{/portal/purchase-history} & Trang & ---\\
\rt{/portal/purchase-history/<int:order_id>} & Trang & ---\\
\rt{/portal/purchase-history/<int:order_id>.pdf} & Trang & ---\\
\rt{/portal/purchase-history/results} & Trang & ---\\
\rt{/portal/purchase_history} & Trang & ---\\
\rt{/portal/reports/orders} & Trang & ---\\
\rt{/portal/reports/orders/export.xlsx} & Trang & ---\\
\end{longtable}
\normalsize

\section{Chương 6 --- Công nợ · Đổi trả}
\footnotesize
\begin{longtable}{L{7.0cm}L{1.8cm}L{6.0cm}}
\rowcolor{hdrbg}\textbf{Đường dẫn} & \textbf{Kiểu} & \textbf{Làm gì}\\ \midrule
\endfirsthead
\rowcolor{hdrbg}\textbf{Đường dẫn} & \textbf{Kiểu} & \textbf{Làm gì}\\ \midrule
\endhead
\rt{/portal/debt} & Trang & ---\\
\rt{/portal/debt/pay} & Trang & ---\\
\rt{/portal/debt/payment-history} & Trang & ---\\
\rt{/portal/return} & Trang & ---\\
\rt{/portal/return-request-list} & Trang & ---\\
\rt{/portal/return/<int:request_id>} & Trang & ---\\
\rt{/portal/return/<int:request_id>/attachment/<int:att_id>} & Trang & ---\\
\rt{/portal/return/new} & Trang & ---\\
\end{longtable}
\normalsize

\section{Chương 7 --- Hỗ trợ · Yêu cầu thông tin · Kiến thức · Thông báo}
\footnotesize
\begin{longtable}{L{7.0cm}L{1.8cm}L{6.0cm}}
\rowcolor{hdrbg}\textbf{Đường dẫn} & \textbf{Kiểu} & \textbf{Làm gì}\\ \midrule
\endfirsthead
\rowcolor{hdrbg}\textbf{Đường dẫn} & \textbf{Kiểu} & \textbf{Làm gì}\\ \midrule
\endhead
\rt{/portal/info-request} & Trang & ---\\
\rt{/portal/info-request/<int:req_id>} & Trang & ---\\
\rt{/portal/info-request/<int:req_id>/cancel} & Trang & ---\\
\rt{/portal/info-request/franchise/<int:fid>/values} & Gọi ngầm & ---\\
\rt{/portal/info-request/new} & Trang & ---\\
\rt{/portal/knowledge} & Trang & ---\\
\rt{/portal/knowledge/<string:slug>} & Trang & ---\\
\rt{/portal/knowledge/<string:slug>/attachment/<int:att_id>} & Trang & ---\\
\rt{/portal/knowledge/search} & Gọi ngầm & ---\\
\rt{/portal/notification} & Trang & ---\\
\rt{/portal/notification/<int:notification_id>} & Trang & ---\\
\rt{/portal/notification/<int:notification_id>/attachment/<int:attachment_id>} & Trang & ---\\
\rt{/portal/notification/mark-all-read} & Gọi ngầm & ---\\
\rt{/portal/notification/mark-read} & Gọi ngầm & ---\\
\rt{/portal/notification/recent} & Gọi ngầm & ---\\
\rt{/portal/notification/results} & Trang & ---\\
\rt{/portal/notification/unread-count} & Gọi ngầm & ---\\
\rt{/portal/support} & Trang & ---\\
\rt{/portal/support/<int:ticket_id>} & Trang & ---\\
\rt{/portal/support/<int:ticket_id>/attachment/<int:att_id>} & Trang & ---\\
\rt{/portal/support/<int:ticket_id>/reply} & Trang & ---\\
\rt{/portal/support/new} & Trang & ---\\
\end{longtable}
\normalsize

\section{Chương 8 --- Đăng ký thi}
\footnotesize
\begin{longtable}{L{7.0cm}L{1.8cm}L{6.0cm}}
\rowcolor{hdrbg}\textbf{Đường dẫn} & \textbf{Kiểu} & \textbf{Làm gì}\\ \midrule
\endfirsthead
\rowcolor{hdrbg}\textbf{Đường dẫn} & \textbf{Kiểu} & \textbf{Làm gì}\\ \midrule
\endhead
\rt{/portal/exam} & Trang & ---\\
\rt{/portal/exam-registration} & Trang & ---\\
\rt{/portal/exam/calendar} & Gọi ngầm & ---\\
\rt{/portal/exam/line/<int:line_id>/photo} & Trang & ---\\
\rt{/portal/exam/register} & Trang & ---\\
\rt{/portal/exam/registration/<int:reg_id>} & Trang & ---\\
\rt{/portal/exam/slots} & Gọi ngầm & ---\\
\end{longtable}
\normalsize

\section{Chương 9 --- Khảo sát cửa hàng · Metabase}
\footnotesize
\begin{longtable}{L{7.0cm}L{1.8cm}L{6.0cm}}
\rowcolor{hdrbg}\textbf{Đường dẫn} & \textbf{Kiểu} & \textbf{Làm gì}\\ \midrule
\endfirsthead
\rowcolor{hdrbg}\textbf{Đường dẫn} & \textbf{Kiểu} & \textbf{Làm gì}\\ \midrule
\endhead
\rt{/franchise/inspection/do/<int:inspection_id>} & Trang & ---\\
\rt{/franchise/inspection/do/<int:inspection_id>/attendance/add} & Gọi ngầm & ---\\
\rt{/franchise/inspection/do/<int:inspection_id>/attendance/deactivate_member} & Gọi ngầm & ---\\
\rt{/franchise/inspection/do/<int:inspection_id>/attendance/delete_line} & Gọi ngầm & ---\\
\rt{/franchise/inspection/do/<int:inspection_id>/attendance/save_member} & Gọi ngầm & ---\\
\rt{/franchise/inspection/do/<int:inspection_id>/save} & Gọi ngầm & ---\\
\rt{/franchise/inspection/do/<int:inspection_id>/sync_posapp} & Gọi ngầm & ---\\
\rt{/metabase/embed/<int:dashboard_id>} & Trang & ---\\
\rt{/metabase/embed_url/<int:dashboard_id>} & Trang & ---\\
\rt{/portal/inspection} & Trang & ---\\
\rt{/portal/inspection/ajax} & Trang & ---\\
\rt{/portal/inspection/detail/<int:inspection_id>} & Trang & ---\\
\rt{/portal/inspection/remediation/<int:line_id>} & Trang & ---\\
\rt{/portal/inspection/remediation/submit} & Trang & ---\\
\rt{/portal/remediation} & Trang & ---\\
\rt{/portal/remediation/<int:line_id>} & Trang & ---\\
\rt{/portal/remediation/submit} & Trang & ---\\
\end{longtable}
\normalsize


\chapter{Phân hệ Đặt hàng}

Phân hệ này gồm \textbf{14 đường dẫn} gom thành \textbf{6 màn} người dùng nhìn thấy, nằm trong
mô-đun \rt{wujia_portal_sale} (tệp \rt{controllers/portal.py}). Đây là phân hệ phức tạp nhất của
portal: nó vừa hiển thị hàng hoá, vừa giữ giỏ, vừa sinh ra đơn hàng thật trong ERP.

\textbf{Bốn nguyên tắc đang chạy}, cần nắm trước khi đọc từng màn:

\begin{enumerate}[leftmargin=*]
  \item \textbf{Giỏ hàng là của CỬA HÀNG, không phải của người dùng.} Mỗi cửa hàng có đúng một
        giỏ (\rt{wujia.portal.cart}, ràng buộc duy nhất trên cửa hàng). Chủ, quản lý và nhân
        viên cùng cửa hàng thao tác trên cùng một giỏ; ai sửa sau thì giá trị của người đó
        thắng. Người khác đang mở màn giỏ sẽ thấy thay đổi \textbf{ngay lập tức} (bản tin
        realtime theo kênh của cửa hàng).
  \item \textbf{Giá luôn do máy chủ tính}, không tin số do trình duyệt gửi lên. Giá lấy theo
        \emph{bảng giá} gắn với khách hàng của cửa hàng, rồi qua bộ tính thuế của Odoo. Cùng một
        hàm dựng giá được dùng cho lưới sản phẩm, trang chi tiết, giỏ và màn kết quả --- nên
        năm chỗ này không bao giờ lệch số.
  \item \textbf{Đơn gửi từ portal luôn ở trạng thái Nháp}, hệ thống không tự xác nhận. Ngô Gia
        xử lý tiếp trong ERP.
  \item \textbf{Một cửa hàng chỉ có một đơn nháp từ portal.} Gửi đơn mới thì đơn nháp/đã gửi cũ
        của chính cửa hàng đó bị huỷ trong cùng một lượt.
\end{enumerate}

\section{Màn Danh sách sản phẩm}
\rt{/portal/order} · \rt{/portal/order/results}
\medskip

\mvao

Vào từ mục \emph{Đặt hàng} trên thanh điều hướng (cột trái ở máy tính, thanh dưới ở điện thoại).
Đường dẫn thứ hai \rt{/portal/order/results} \textbf{không phải màn riêng}: khi người dùng gõ
tìm kiếm hoặc bấm danh mục, trình duyệt gọi ngầm đường dẫn này để lấy về đúng phần kết quả và
thay tại chỗ, không tải lại cả trang. Hai đường dẫn dùng chung một bộ dữ liệu, nên mọi luật mô
tả dưới đây đúng cho cả hai.

\mai

\begin{itemize}[leftmargin=*]
  \item Phải \textbf{đăng nhập}. Khách vãng lai bị đẩy về trang đăng nhập.
  \item Nếu tài khoản \textbf{không thuộc cửa hàng nào} còn hiệu lực, màn hiện trạng thái rỗng
        ``chưa có cửa hàng'' --- không có sản phẩm, không có giỏ.
  \item \textbf{Cửa hàng đang chọn} được lấy từ \emph{cookie} \rt{wujia_active_franchise_id}
        (lưu 30 ngày trên máy người dùng), và chỉ được chấp nhận nếu cửa hàng đó nằm trong danh
        sách cửa hàng mà tài khoản có thành viên còn hiệu lực. Ai chỉ thuộc \textbf{một} cửa
        hàng thì hệ thống tự chọn, không hỏi.
  \item \textbf{Vai trò không ảnh hưởng màn này}: Chủ, Quản lý hay Nhân viên đều xem và thêm vào
        giỏ như nhau. Phân biệt vai trò chỉ xuất hiện ở các phân hệ khác (công nợ, thành viên).
  \item Nếu \textbf{chưa chọn được cửa hàng} (tài khoản có nhiều cửa hàng mà chưa bấm chọn),
        màn \textbf{vẫn hiện đủ sản phẩm} nhưng: giá rơi về \emph{giá niêm yết} của sản phẩm
        thay vì giá theo bảng giá cửa hàng, giỏ hàng trống, và có dải cảnh báo
        ``Vui lòng chọn cửa hàng trước khi đặt hàng''. Đây là điểm dễ hiểu nhầm --- xem mục 7.
  \item Dữ liệu sản phẩm được đọc \textbf{bằng quyền hệ thống} (bỏ qua phân quyền của người
        dùng portal). Việc giới hạn nằm hoàn toàn ở điều kiện lọc bên dưới, không dựa vào
        quy tắc quyền của Odoo.
\end{itemize}

\mhien

\begin{hienbang}
Lưới sản phẩm &
  \rt{product.product} &
  Bật cờ \emph{Hiện trên portal} (\rt{is_public_portal = true}) \textbf{và} chưa lưu trữ
  (\rt{active = true}). Có từ khoá thì thêm: tên \textbf{hoặc} mã sản phẩm chứa từ khoá. Có
  chọn danh mục thì thêm: đúng danh mục đó. &
  Theo tên A$\rightarrow$Z; 24 sản phẩm/trang &
  \uat{5}/8 sản phẩm đang bật cờ \\ \addlinespace
Bộ lọc danh mục &
  \rt{wujia.product.category} &
  Chỉ những danh mục \textbf{đang có ít nhất một sản phẩm hiện trên portal}. Danh mục đã lưu
  trữ tự động không hiện. &
  Theo \emph{Thứ tự} rồi tên &
  \uat{2} danh mục (Trà, Topping) \\ \addlinespace
Giá từng sản phẩm &
  Bảng giá của khách hàng gắn với cửa hàng &
  Lấy bảng giá ở trường \emph{Bảng giá} của khách hàng cửa hàng; chưa có bảng giá (hoặc chưa
  chọn cửa hàng) thì dùng \emph{giá niêm yết} của sản phẩm. Số hiển thị là giá \textbf{đã gồm
  thuế}, tính theo vị thế tài khoán của chính khách hàng đó. &
  --- &
  4/4 cửa hàng đều đã gắn khách hàng \\ \addlinespace
Số trong giỏ trên mỗi thẻ &
  \rt{wujia.portal.cart.line} &
  Dòng thuộc giỏ của cửa hàng đang chọn. &
  --- &
  \uat{6} dòng trong \uat{2} giỏ \\ \addlinespace
Dải khung giờ đặt hàng &
  \rt{wujia.order.window} + tham số chung &
  Theo khu vực của cửa hàng; khu vực chưa cấu hình thì rơi về khung giờ chung. Chi tiết ở
  \textbf{mục 3.7}. &
  --- &
  \uat{1} khung theo khu vực / 2 khu vực \\
\end{hienbang}

\textbf{Trường nào hiện lên thẻ sản phẩm:} ảnh sản phẩm, tên (\rt{name}), tên tiếng Trung
(\rt{name_chinese}, để trống thì không hiện dòng đó), mã (\rt{default_code}), quy cách đóng gói
(\rt{wujia_packaging}, trống thì hiện dấu ``---''), đơn vị tính, giá đã gồm thuế, số lượng tối
thiểu (\rt{min_qty}) dùng làm \textbf{bước tăng giảm}.

\mloc

\begin{itemize}[leftmargin=*]
  \item \textbf{Ô tìm kiếm} --- tìm theo \emph{tên} hoặc \emph{mã sản phẩm}, không phân biệt hoa
        thường, khớp một phần. \textbf{Không} tìm theo tên tiếng Trung và \textbf{không} tìm
        theo quy cách.
  \item \textbf{Danh mục} --- một danh mục tại một thời điểm. Giá trị lạ (không phải số) bị bỏ
        qua, coi như không lọc.
  \item \textbf{Số sản phẩm mỗi trang} --- chỉ nhận \textbf{24, 48, 96}; gửi số khác thì hệ
        thống lặng lẽ dùng 24.
  \item \textbf{Trang} --- số trang vượt quá trang cuối thì tự kéo về trang cuối, không báo lỗi.
        Đổi trang giữ nguyên từ khoá và danh mục đang chọn.
\end{itemize}

\mkiem{Thêm vào giỏ}

Nút này không tải lại trang mà gọi ngầm \rt{/portal/order/cart/add}. Thứ tự kiểm đúng như code
đang chạy:

\begin{kiembang}
1 & Đã chọn được cửa hàng hợp lệ &
    ``Vui lòng chọn cửa hàng trước khi đặt hàng.'' --- hoặc ``Bạn không có quyền thao tác tại
    cửa hàng này'' / ``Tài khoản của bạn hiện không còn hiệu lực tại cửa hàng này'' nếu cookie
    còn trỏ tới cửa hàng cũ \\
2 & Mã sản phẩm gửi lên là số & ``Dữ liệu gửi lên không hợp lệ.'' \\
3 & Sản phẩm còn tồn tại, chưa lưu trữ, còn bật cờ hiện trên portal &
    ``Sản phẩm này hiện không còn được phép đặt hàng.'' \\
4 & Sản phẩm đã cấu hình số lượng tối thiểu (\rt{min_qty > 0}) &
    ``Sản phẩm chưa được cấu hình số lượng đặt tối thiểu. Vui lòng liên hệ Ngô Gia.'' \\
5 & Số lượng gửi lên (nếu có) là \textbf{số nguyên} &
    ``Số lượng của \emph{tên sản phẩm} phải là số nguyên.'' --- 1,5 bị từ chối, \textbf{không}
    làm tròn \\
6 & Số lượng $\geq$ tối thiểu & ``Số lượng tối thiểu của \emph{tên sản phẩm} là \emph{n}.'' \\
7 & Số lượng là \textbf{bội số} của tối thiểu &
    ``Số lượng của \emph{tên sản phẩm} phải tăng theo bước \emph{n}.'' \\
8 & Số lượng $\leq$ tối đa (khi \rt{max_qty > 0}) &
    ``Số lượng tối đa của \emph{tên sản phẩm} là \emph{n}.'' \\
\end{kiembang}

Bấm từ lưới sản phẩm thì trình duyệt \textbf{không gửi số lượng}, hệ thống tự cộng đúng một
bước (bằng số lượng tối thiểu). Các bước 5--8 chỉ áp dụng khi người dùng gõ số tay ở trang chi
tiết. Nếu cộng dồn vượt trần tối đa, hệ thống \textbf{cắt xuống đúng trần} và trả về cảnh báo
chứ không báo lỗi.

\mghi

\begin{itemize}[leftmargin=*]
  \item Tạo giỏ cho cửa hàng nếu chưa có, rồi \textbf{thêm hoặc cộng dồn} dòng
        \rt{wujia.portal.cart.line} (mỗi sản phẩm một dòng duy nhất trong giỏ).
  \item Thao tác cộng dồn chạy thẳng bằng một câu lệnh cơ sở dữ liệu nên hai người cùng bấm một
        lúc \textbf{không mất số của nhau}.
  \item Phát một bản tin realtime tới kênh của cửa hàng để mọi thiết bị khác đang mở giỏ tự cập
        nhật.
  \item Không ghi gì vào đơn hàng --- giỏ và đơn là hai thứ tách rời.
\end{itemize}

\mhoi

\begin{ask}
\begin{enumerate}[leftmargin=*]
  \item \textbf{Sản phẩm chưa gán danh mục vẫn hiện.} Trên UAT có \uat{2}/5 sản phẩm bật cờ
        portal nhưng bỏ trống \emph{Danh mục portal} (\emph{Hồng Trà Sữa size M},
        \emph{Matcha Latte size M}). Chúng hiện ở danh sách chung nhưng \textbf{không bao giờ
        xuất hiện} khi người dùng bấm lọc theo danh mục. Anh muốn (a) bắt buộc phải có danh mục
        mới cho hiện, hay (b) thêm nhóm ``Khác''?
  \item \textbf{Chưa chọn cửa hàng thì giá đang là giá niêm yết.} Người dùng nhiều cửa hàng, khi
        chưa bấm chọn, vẫn thấy giá --- nhưng là giá gốc, không phải giá theo bảng giá của cửa
        hàng. Có nên ẩn hẳn giá cho tới khi chọn cửa hàng không?
  \item \textbf{Tìm kiếm chưa bao gồm tên tiếng Trung.} Màn có hiện \rt{name_chinese} nhưng ô
        tìm kiếm không tìm theo trường đó. Có cần bổ sung không?
  \item \textbf{Trên UAT chưa sản phẩm nào đặt trần tối đa} (\rt{max_qty} đều bằng 0 = không
        giới hạn). Toàn bộ luật chặn ``số lượng tối đa'' vì vậy chưa từng chạy thật khi BA test.
\end{enumerate}
\end{ask}

\section{Màn Chi tiết sản phẩm}
\rt{/portal/order/product/<int:product_id>}
\medskip

\mvao

Bấm vào một thẻ trong lưới sản phẩm.

\mai

Như màn danh sách. Riêng màn này, nếu sản phẩm \textbf{không tồn tại, đã lưu trữ, hoặc đã tắt cờ
hiện trên portal}, người dùng bị đưa về màn danh sách kèm thông báo ``Sản phẩm này hiện không
còn được phép đặt hàng'' --- hệ thống \textbf{cố ý không nói rõ} là sản phẩm không tồn tại hay
người dùng không được xem, để không lộ dữ liệu.

\mhien

\begin{hienbang}
Khối thông tin sản phẩm &
  \rt{product.product} &
  Chính sản phẩm được mở, phải đang \rt{is_public_portal = true} và \rt{active = true} &
  --- & --- \\ \addlinespace
Gợi ý ``sản phẩm liên quan'' &
  \rt{product.product} &
  Cùng \emph{danh mục portal} với sản phẩm đang xem, trừ chính nó, và cũng phải đang hiện trên
  portal. Sản phẩm \textbf{không có danh mục} thì khối này rỗng. &
  Theo tên A$\rightarrow$Z; tối đa 6 &
  2/5 sản phẩm không có danh mục $\Rightarrow$ khối rỗng \\ \addlinespace
Giá &
  Bảng giá của cửa hàng &
  Như màn danh sách --- cùng một hàm tính, nên hai màn luôn cùng số. &
  --- & --- \\
\end{hienbang}

\mloc

Không có bộ lọc. Người dùng chỉ gõ được \textbf{số lượng} trước khi bấm thêm vào giỏ.

\mkiem{Thêm vào giỏ}

Cùng đúng 8 bước kiểm của màn danh sách. Khác một điểm: ở đây trình duyệt \textbf{gửi số lượng
tường minh}, nên các bước 5--8 (số nguyên, tối thiểu, bội số, tối đa) thực sự chạy.

\mghi

Giống màn danh sách.

\mhoi

\begin{ask}
Khối ``sản phẩm liên quan'' đang dựa hoàn toàn vào danh mục portal. Với dữ liệu hiện tại, sản
phẩm không có danh mục sẽ có khối này trống trơn. Anh muốn thay bằng ``sản phẩm hay đặt cùng''
hay để nguyên?
\end{ask}

\section{Màn Giỏ hàng}
\rt{/portal/order/cart} --- kèm 6 đường dẫn thao tác ngầm
\medskip

\mvao

Bấm biểu tượng giỏ trên thanh trên cùng, hoặc nút ``Xem giỏ hàng''. Sáu đường dẫn còn lại không
phải màn, mà là các thao tác trình duyệt gọi ngầm khi người dùng bấm nút:

\begin{itemize}[leftmargin=*,itemsep=1pt]
  \item \rt{/portal/order/cart/add} --- thêm sản phẩm (đã mô tả ở màn danh sách)
  \item \rt{/portal/order/cart/step} --- bấm nút \emph{+} / \emph{--} (tăng giảm đúng một bước)
  \item \rt{/portal/order/cart/update} --- gõ thẳng số lượng vào ô
  \item \rt{/portal/order/cart/remove} --- bấm xoá dòng
  \item \rt{/portal/order/cart/count} --- lấy số hiển thị trên biểu tượng giỏ
  \item \rt{/portal/order/cart/note} --- lưu ghi chú đơn hàng
  \item \rt{/portal/order/cart/fragment} --- lấy lại toàn bộ giỏ khi có người khác vừa sửa
\end{itemize}

\mai

Người dùng \textbf{chưa thuộc cửa hàng nào} bị đưa thẳng về màn danh sách sản phẩm. Ngoài ra như
màn danh sách: giỏ là giỏ chung của cửa hàng đang chọn, ai trong cửa hàng cũng sửa được, không
phân biệt vai trò.

\mhien

\begin{hienbang}
Danh sách dòng &
  \rt{wujia.portal.cart.line} &
  Thuộc giỏ của cửa hàng đang chọn (mỗi cửa hàng đúng một giỏ) &
  Không phân trang &
  \uat{6} dòng / \uat{2} giỏ \\ \addlinespace
Đơn giá từng dòng &
  Bảng giá của cửa hàng &
  Tính lại \textbf{theo số lượng của chính dòng đó} --- bảng giá có bậc số lượng thì giá đổi
  theo. Giỏ \textbf{không lưu giá}, mỗi lần mở là tính lại. &
  --- & --- \\ \addlinespace
Tạm tính · Thuế · Tổng thanh toán &
  Tính tại máy chủ &
  Cộng từ các dòng, làm tròn theo tiền tệ của bảng giá. Trình duyệt không tự cộng. &
  --- & --- \\ \addlinespace
Cảnh báo trên dòng &
  Tính tại máy chủ &
  Mỗi dòng được soi lại ba điều: sản phẩm còn được phép đặt không · đã cấu hình tối thiểu chưa ·
  số lượng có đúng tối thiểu/bội số/tối đa không. &
  --- & --- \\ \addlinespace
Ghi chú đơn hàng &
  \rt{wujia.portal.cart} trường \rt{note} &
  Ghi chú \textbf{dùng chung cả cửa hàng}, tối đa 1000 ký tự. &
  --- & --- \\
\end{hienbang}

\mloc

Không có bộ lọc. Người dùng sửa được số lượng từng dòng và ghi chú chung.

\mkiem{+ / -- / gõ số lượng / xoá}

\begin{kiembang}
1 & Đã chọn được cửa hàng hợp lệ & Thông báo như bảng ở màn danh sách \\
2 & Dòng còn tồn tại và \textbf{thuộc đúng giỏ của cửa hàng đang chọn} &
    \textbf{Không báo lỗi}: coi như dòng đã bị người cùng cửa hàng xoá, trả về giỏ mới nhất và
    màn tự cập nhật \\
3 & Sản phẩm đã cấu hình số lượng tối thiểu &
    ``Sản phẩm chưa cấu hình số lượng tối thiểu.'' \\
4 & Số lượng mới $\geq$ tối thiểu, đúng bội số, $\leq$ tối đa &
    Thông báo tương ứng (tối thiểu / bước / tối đa) \\
\end{kiembang}

Hai hành vi đáng chú ý: \textbf{gõ số lượng bằng 0} hoặc \textbf{bấm ``--'' xuống dưới mức tối
thiểu} đều \textbf{xoá dòng khỏi giỏ}, không báo lỗi. Và nút \emph{+} không cho trình duyệt gửi
số lượng tuỳ ý --- nó chỉ gửi \emph{hướng} tăng hay giảm, máy chủ tự lấy bước --- nên hai người
bấm cùng lúc không ghi đè nhau.

\mghi

\begin{itemize}[leftmargin=*]
  \item Sửa hoặc xoá dòng trong giỏ của cửa hàng; lưu ghi chú vào giỏ.
  \item Mỗi thay đổi phát một bản tin realtime tới kênh của cửa hàng.
  \item Không đụng tới đơn hàng.
\end{itemize}

\mhoi

\begin{ask}
\begin{enumerate}[leftmargin=*]
  \item \textbf{Ghi chú là của cửa hàng, không của người gửi.} Nhân viên A gõ ghi chú, nhân
        viên B gửi đơn thì ghi chú của A đi theo đơn. Anh xác nhận đúng ý nghiệp vụ chứ?
  \item \textbf{Giỏ không có lịch sử.} Khi hai người sửa chồng nhau, hệ thống lấy giá trị của
        người sửa sau và không lưu lại ai đã sửa gì. Có cần nhật ký thao tác giỏ không?
\end{enumerate}
\end{ask}

\section{Gửi đơn hàng}
\rt{/portal/order/submit}
\medskip

\mvao

Nút \emph{Gửi đơn} ở màn giỏ. Đây là đường dẫn \textbf{ghi dữ liệu thật} --- nó tạo đơn hàng
trong ERP, nên có chống giả mạo biểu mẫu (CSRF) và chỉ nhận phương thức POST.

\mai

Chỉ người có thành viên còn hiệu lực tại cửa hàng đang chọn. Vai trò \textbf{không} giới hạn:
nhân viên cũng gửi được đơn.

\mhien

Đường dẫn này \textbf{không hiện màn nào} --- nó chỉ nhận lệnh, ghi dữ liệu rồi chuyển hướng.
Gửi thành công: trên điện thoại sang màn \emph{Đặt hàng thành công} (mục 3.5), trên máy tính sang
thẳng trang chi tiết đơn ở phân hệ \emph{Lịch sử mua hàng}. Gửi hỏng: quay về màn giỏ kèm dải đỏ,
riêng trường hợp ngoài khung giờ trên điện thoại thì sang màn \emph{Không thể gửi đơn} (mục 3.6).

\mloc

\textbf{Không có ô nhập nào.} Toàn bộ nội dung đơn lấy từ giỏ của cửa hàng đang chọn --- trình
duyệt \textbf{không} gửi lên danh sách sản phẩm, số lượng hay giá. Đây là chủ đích: người dùng
không thể can thiệp vào nội dung đơn bằng cách sửa dữ liệu gửi lên.

\mkiem{Gửi đơn}

Đây là chuỗi kiểm dài nhất của portal --- \textbf{11 lớp}, đúng thứ tự chạy:

\begin{kiembang}
1 & Chọn được cửa hàng hợp lệ & Về màn giỏ, dải đỏ theo từng trường hợp (chưa chọn / không có
    quyền / hết hiệu lực) \\
2 & Cửa hàng còn tồn tại & ``Bạn không có quyền thao tác tại cửa hàng này.'' \\
3 & Cửa hàng \textbf{không bị khoá đặt hàng} (\rt{portal_locked}) &
    ``Cửa hàng đang tạm khóa đặt hàng. Vui lòng liên hệ Ngô Gia.'' \\
4 & Giỏ có ít nhất một dòng & ``Giỏ hàng chưa có sản phẩm.'' \\
5 & \textbf{Không có người khác đang gửi đơn cùng lúc} (giỏ bị khoá trong lúc xử lý) &
    ``Giỏ hàng đang được một người dùng khác gửi đơn. Vui lòng thử lại sau.'' \\
6 & \textbf{Đang trong khung giờ đặt hàng} &
    Điện thoại: chuyển sang màn riêng ``Không thể gửi đơn'' có ghi giờ mở lại. Máy tính: về giỏ
    kèm dải ``Hiện ngoài khung giờ đặt hàng.'' \\
7 & Mọi dòng trong giỏ đều hợp lệ (sản phẩm còn bán, số lượng đúng bước) &
    ``Một số sản phẩm trong giỏ không còn được phép đặt'' hoặc ``Số lượng một số sản phẩm chưa
    hợp lệ'' \\
8 & Cửa hàng đã gắn \textbf{khách hàng} &
    ``Cửa hàng chưa được cấu hình khách hàng đặt hàng. Vui lòng liên hệ Ngô Gia.'' \\
9 & Đơn nháp cũ của cửa hàng \textbf{không bị khoá} &
    ``Không thể hoàn tất đơn hàng do đơn nháp cũ chưa được xử lý.'' \\
10 & Tạo được đơn hàng (ERP không từ chối) &
    ``Không thể tạo đơn hàng. Vui lòng thử lại hoặc liên hệ Ngô Gia.'' \\
11 & Huỷ được đơn nháp cũ \textbf{và} dọn được giỏ &
    Hoàn tác toàn bộ --- kể cả đơn vừa tạo --- rồi báo lỗi. Không bao giờ để tồn tại hai đơn
    nháp. \\
\end{kiembang}

Lớp 6 được kiểm \textbf{hai lần}: một lần ở màn, một lần nữa ngay trong ERP lúc tạo đơn. Kể cả
khi ai đó gọi thẳng đường dẫn, đơn ngoài khung giờ vẫn bị chặn.

\mghi

\begin{itemize}[leftmargin=*]
  \item \textbf{Tạo đơn} \rt{sale.order} ở trạng thái \textbf{Nháp}, đánh dấu \rt{is_portal_order},
        gắn cửa hàng, khách hàng của cửa hàng, bảng giá tương ứng, nguồn ``Wujia Portal'',
        người gửi và bản ghi thành viên tại thời điểm gửi.
  \item Ghi chú của giỏ được chép sang đơn ở \textbf{hai trường}: một trường văn bản thuần cho
        trang Lịch sử của portal, một trường ghi chú chuẩn của Odoo.
  \item Mỗi dòng giỏ thành một dòng đơn (sản phẩm, số lượng, đơn vị tính). \textbf{Giá không
        chép sang} --- Odoo tự tính lại theo bảng giá đã gắn, nên giá trên đơn khớp giá người
        dùng vừa thấy.
  \item \textbf{Huỷ} mọi đơn portal cũ của cửa hàng đang ở Nháp/Đã gửi.
  \item \textbf{Dọn sạch} dòng và ghi chú của giỏ, rồi phát bản tin realtime.
  \item Người đứng tên đơn là \textbf{chính tài khoản portal} đã bấm gửi.
\end{itemize}

\mhoi

\begin{ask}
\begin{enumerate}[leftmargin=*]
  \item \textbf{Nhân viên gửi được đơn.} Hiện không phân biệt vai trò ở bước gửi. Nếu nghiệp vụ
        muốn ``chỉ Chủ/Quản lý được gửi đơn'' thì đây là chỗ cần thêm luật.
  \item \textbf{Huỷ đơn nháp cũ là im lặng.} Cửa hàng gửi đơn lần hai trong ngày sẽ mất đơn lần
        một mà không có cảnh báo trước. Có cần hộp xác nhận ``đơn cũ sẽ bị huỷ'' không?
  \item \textbf{Đơn không tự xác nhận.} Hiện Ngô Gia phải vào ERP xác nhận. Đúng như thiết kế
        chứ, hay cần tự xác nhận trong khung giờ?
\end{enumerate}
\end{ask}

\section{Màn Đặt hàng thành công}
\rt{/portal/order/submitted/<int:order_id>}
\medskip

\mvao

Tự chuyển tới sau khi gửi đơn thành công \textbf{trên điện thoại}. Trên máy tính, người dùng
được đưa thẳng sang trang chi tiết đơn ở phân hệ Lịch sử mua hàng --- hai luồng khác nhau có chủ ý.

\mai

Chỉ xem được đơn \textbf{của chính cửa hàng đang chọn} và phải là đơn đặt qua portal. Sai điều
kiện --- kể cả khi gõ tay số đơn của cửa hàng khác --- thì bị đưa về màn danh sách sản phẩm,
không báo gì thêm. Đơn đã huỷ cũng bị chặn ở đây, vì màn này chỉ nói về đơn vừa gửi thành công.

\mhien

\begin{hienbang}
Mã đơn, ngày giờ gửi &
  \rt{sale.order} &
  Đúng đơn vừa tạo; ngày giờ quy về múi giờ người dùng &
  --- &
  \uat{23} đơn đặt từ portal \\ \addlinespace
Tạm tính · Thuế · Tổng &
  \rt{sale.order} &
  Lấy thẳng số tiền của đơn, ký hiệu tiền tệ theo chính đơn đó (không viết cứng ``đ'') &
  --- & --- \\ \addlinespace
Trạng thái đơn &
  \rt{sale.order} trường \rt{state} &
  Dùng \textbf{chung bảng nhãn} với trang Lịch sử mua hàng --- nháp hiển thị là
  ``Chờ xác nhận''. &
  --- &
  \uat{3} đơn portal đang chờ xác nhận \\
\end{hienbang}

\mloc\ Không có ô nhập nào. \quad
\mkiem{Về trang chủ / Xem chi tiết đơn}\ Hai nút trên màn chỉ là liên kết sang màn khác,
không kiểm gì. \quad \mghi\ Không ghi gì --- màn chỉ đọc.

\mhoi

\begin{ask}
Máy tính và điện thoại kết thúc ở \textbf{hai màn khác nhau} sau khi gửi đơn (máy tính sang
trang chi tiết đơn, điện thoại ở màn xác nhận này). Anh xác nhận giữ nguyên chứ?
\end{ask}

\section{Màn Không thể gửi đơn}
\rt{/portal/order/rejected}
\medskip

\mvao

Chỉ dùng cho \textbf{một trường hợp duy nhất}: gửi đơn ngoài khung giờ, trên điện thoại. Mọi lý
do khác đều quay về màn giỏ kèm dải cảnh báo, vì ở đó người dùng sửa được ngay. Nếu ai gõ tay
đường dẫn với lý do lạ, hệ thống đưa về giỏ --- \textbf{không hiển thị lại chuỗi người dùng tự
gõ}, tránh bị lợi dụng để hiện nội dung giả.

\mai\ Như màn giỏ.

\mhien

\begin{hienbang}
Khung giờ mở lại &
  \rt{wujia.order.window} hoặc tham số chung &
  Lấy khung gần nhất sắp mở; nếu bây giờ đã qua giờ mở của mọi khung thì lấy khung đầu tiên của
  \textbf{ngày mai}. Có ghi rõ ``Hôm nay'' hay ngày cụ thể. &
  --- &
  Khu vực Hà Nội chưa cấu hình $\Rightarrow$ dùng khung chung 10:00--04:00 \\
\end{hienbang}

\mloc\ Không có ô nhập nào. \quad
\mkiem{Quay lại giỏ hàng}\ Nút duy nhất trên màn chỉ là liên kết về màn giỏ, không kiểm gì.
\quad \mghi\ Không ghi gì.

\mhoi\ Không có điểm treo.

\section{Khung giờ đặt hàng --- luật dùng chung}
\label{sec:khung-gio}

Đây là luật bị hỏi nhiều nhất, nên tách riêng. Hệ thống quyết định ``có đang trong giờ đặt hàng
không'' theo \textbf{ba bậc}:

\begin{enumerate}[leftmargin=*]
  \item \textbf{Công tắc tổng}. Tham số \rt{portal_order_time_limit_enabled}. Tắt thì đặt hàng
        \textbf{mọi lúc}, không xét gì thêm. \note{(UAT: đang BẬT)}
  \item \textbf{Khung giờ theo khu vực}. Nếu khu vực của cửa hàng có ít nhất một khung
        (\rt{wujia.order.window}) đang hoạt động, hệ thống chỉ xét các khung đó và cho đặt khi
        \textbf{bất kỳ khung nào} đang mở. Một khu vực được phép có nhiều khung (sáng + chiều).
  \item \textbf{Khung giờ chung}. Khu vực không có khung riêng thì dùng cặp tham số chung
        \rt{portal_order_time_from} / \rt{portal_order_time_to}. \note{(UAT: 10:00--04:00)}
\end{enumerate}

Khung có giờ kết thúc \textbf{nhỏ hơn} giờ bắt đầu được hiểu là \textbf{vắt qua nửa đêm} (ví dụ
10:00 hôm nay đến 04:00 ngày mai); màn hiển thị ghi rõ ``hôm nay'' / ``ngày mai'' kèm ``UTC+7''
để không hiểu nhầm.

\textbf{Hiện trạng trên UAT:}

\begin{center}\footnotesize
\begin{tabular}{L{3.9cm}L{3.3cm}L{3.9cm}L{3.5cm}}
\rowcolor{hdrbg}\textbf{Cửa hàng} & \textbf{Khu vực} & \textbf{Khung giờ áp dụng} &
  \textbf{Nguồn}\\ \midrule
TP HCM Quận 1 (HCM-01) & TP HCM Quận 1--3 & 00:00 -- 23:50 & Khung của khu vực \\
Cầu Giấy (HN-01) & Hà Nội Trung tâm & 10:00 -- 04:00 & Khung chung \\
Đống Đa (HN-02) & Hà Nội Trung tâm & 10:00 -- 04:00 & Khung chung \\
Cửa hàng TEST Issue 134 & Hà Nội Trung tâm & 10:00 -- 04:00 & Khung chung \\
\end{tabular}
\end{center}

\mhoi

\begin{ask}
\begin{enumerate}[leftmargin=*]
  \item \textbf{Tên khung giờ không khớp giờ thật.} Khung duy nhất trên UAT \textbf{tên là}
        ``00:00 - 05:00'' nhưng giờ thật đang là \textbf{00:00 -- 23:50}. Tên là chữ người dùng
        tự gõ, hệ thống không kiểm --- và chính chuỗi tên này hiện lên portal. Có cần bắt hệ
        thống tự sinh tên theo giờ không?
  \item \textbf{Giờ được tính theo múi giờ của TÀI KHOẢN đăng nhập}, không theo khu vực của cửa
        hàng, trong khi dải thông báo luôn ghi ``UTC+7''. Tài khoản chưa đặt múi giờ sẽ bị tính
        theo giờ UTC --- lệch 7 tiếng so với chữ hiển thị. Đề nghị chốt: lấy múi giờ theo
        \textbf{khu vực cửa hàng}, hay ép mọi tài khoản portal về giờ Việt Nam?
  \item \textbf{Khu vực Hà Nội (3/4 cửa hàng) chưa có khung riêng} nên đang chạy khung chung. BA
        xác nhận đây là ý muốn, hay cần cấu hình khung riêng cho Hà Nội?
\end{enumerate}
\end{ask}

\chapter{Khung portal --- Đăng nhập, Trang chủ, Cửa hàng}

Chương này gồm \textbf{25 đường dẫn} gom thành \textbf{13 màn}, nằm trong hai mô-đun:

\begin{itemize}[leftmargin=*,itemsep=1pt]
  \item \rt{wujia_portal_layout} (11 đường dẫn) --- \emph{cái khung}: đăng nhập, quên và đổi
        mật khẩu, hồ sơ cá nhân, đổi ngôn ngữ. Phần này cố ý \textbf{không biết gì} về cửa
        hàng, đơn hàng hay công nợ, để màn đăng nhập vẫn chạy được kể cả khi các phân hệ
        nghiệp vụ chưa bật.
  \item \rt{wujia_portal_base} (14 đường dẫn) --- \emph{cái nền ghép}: Trang chủ tổng hợp số
        liệu từ các phân hệ, việc chọn cửa hàng đang làm việc, và các màn thông tin cửa hàng.
\end{itemize}

\textbf{Năm nguyên tắc chi phối toàn bộ portal}, mô tả một lần ở đây, các chương sau chỉ trỏ về:

\begin{enumerate}[leftmargin=*]
  \item \textbf{Cửa hàng đang làm việc nằm trong cookie, không nằm trong tài khoản.} Trình duyệt
        giữ \rt{wujia_active_franchise_id} trong \textbf{30 ngày}. Mỗi lần mở trang, hệ thống
        kiểm lại: cửa hàng trong cookie có nằm trong danh sách cửa hàng tài khoản còn hiệu lực
        không --- không thì coi như chưa chọn. Ai chỉ thuộc \textbf{một} cửa hàng thì hệ thống
        tự chọn, không hỏi. Hệ quả: đổi máy hoặc xoá dữ liệu trình duyệt là phải chọn lại.
  \item \textbf{Chưa chọn cửa hàng thì dữ liệu là TỔNG của mọi cửa hàng.} Đây là điểm dễ hiểu
        nhầm nhất. Với người thuộc nhiều cửa hàng mà chưa bấm chọn, hệ thống \emph{không} chặn,
        mà lấy dữ liệu của \textbf{tất cả} cửa hàng họ thuộc. Trang chủ vì thế có thể hiện con
        số cộng gộp của 3 cửa hàng.
  \item \textbf{Vai trò có ba bậc: Nhân viên $<$ Quản lý $<$ Chủ tiệm.} Bậc dùng để so sánh
        ``ít nhất phải là\ldots''. Một tài khoản thuộc nhiều cửa hàng thì vai trò được lấy là
        \textbf{bậc cao nhất} trong các cửa hàng đang xét. Rất nhiều màn hiện nay \emph{không}
        phân biệt vai trò --- chỗ nào có phân biệt sẽ được nói rõ ở mục 2 của màn đó.
  \item \textbf{Hầu hết dữ liệu được đọc bằng quyền hệ thống.} Nghĩa là phân quyền chuẩn của
        Odoo bị bỏ qua, và việc ``ai thấy gì'' \textbf{hoàn toàn} do điều kiện lọc trong từng
        màn quyết định. Nếu một điều kiện lọc viết thiếu, dữ liệu cửa hàng khác sẽ lọt ra --- đó
        là lý do tài liệu này liệt kê từng điều kiện một.
  \item \textbf{Giờ hiển thị đổi theo múi giờ của TÀI KHOẢN.} Hệ thống lưu giờ chuẩn quốc tế,
        khi hiện lên mới quy đổi theo ô \emph{Múi giờ} trong tài khoản người dùng; ô trống thì
        dùng \emph{Asia/Ho\_Chi\_Minh}. Không phải theo cửa hàng, cũng không phải theo khu vực.
\end{enumerate}

\section{Màn Đăng nhập}
\rt{/portal/login} · \rt{/portal/logout} · \rt{/portal/redirect}
\medskip

\mvao

Là màn đầu tiên. Mọi đường dẫn portal khi chưa đăng nhập đều bị đẩy về đây.
\rt{/portal/logout} là nút Đăng xuất; \rt{/portal/redirect} \textbf{không phải màn}, mà là
trạm trung chuyển ngay sau khi đăng nhập thành công.

\mai

\begin{itemize}[leftmargin=*]
  \item Màn này \textbf{không yêu cầu đăng nhập} (đương nhiên). Đã đăng nhập rồi mà mở lại
        \rt{/portal/login} thì bị đẩy thẳng qua trạm trung chuyển.
  \item Trạm trung chuyển phân luồng theo \textbf{loại tài khoản}: tài khoản \emph{nội bộ}
        (nhân sự Ngô Gia) đi vào giao diện quản trị \rt{/odoo}; tài khoản \emph{cửa hàng} đi
        vào \rt{/portal}. Trên UAT: \uat{1} tài khoản nội bộ, \uat{6} tài khoản cửa hàng.
  \item Nếu địa chỉ gọi tới có kèm tham số ``sau khi vào thì đi đâu'', hệ thống đi theo tham số
        đó \textbf{trước}, không qua trạm phân luồng.
\end{itemize}

\mhien

\begin{hienbang}
Ô tài khoản, ô mật khẩu &
  \rt{res.users} &
  Không truy vấn gì lúc hiện màn. Ô tài khoản được điền sẵn nếu phiên trước có ghi nhận. &
  --- & \uat{7} tài khoản đang có \\ \addlinespace
Nút đổi ngôn ngữ &
  \rt{res.lang} &
  Chỉ ngôn ngữ \textbf{đang bật} (\rt{active = true}) &
  --- & \uat{3} ngôn ngữ: Tiếng Việt, English, Thai \\
\end{hienbang}

\mloc

Hai ô: \emph{tài khoản} (chính là email đăng nhập) và \emph{mật khẩu}. Người chưa đăng nhập
vẫn đổi được ngôn ngữ ngay tại màn này --- lựa chọn đó được nhớ tạm và \textbf{ghi vào tài
khoản} ngay sau khi đăng nhập thành công.

\mkiem{Đăng nhập}

\begin{kiembang}
1 & Xác định đúng cơ sở dữ liệu cần đăng nhập &
    Máy chủ chỉ có một cơ sở dữ liệu nên bước này luôn qua \\
2 & Tài khoản và mật khẩu khớp &
    ``Wrong login/password''. Hệ thống \textbf{không} nói rõ sai tài khoản hay sai mật khẩu \\
3 & Tài khoản còn hoạt động &
    Thông báo do Odoo sinh ra (ví dụ tài khoản đã bị tắt) \\
\end{kiembang}

\textbf{Đáng lưu ý:} màn đăng nhập \textbf{không} có cơ chế giới hạn số lần thử của riêng
portal --- chỉ màn \emph{Quên mật khẩu} mới có (10 lần/giờ/địa chỉ mạng). Xem mục 7.

\mghi

\begin{itemize}[leftmargin=*]
  \item Mở một \emph{phiên làm việc} cho tài khoản. Không tạo bản ghi nghiệp vụ nào.
  \item Nếu người dùng đã đổi ngôn ngữ ở màn đăng nhập, hệ thống \textbf{ghi ngôn ngữ đó vào
        tài khoản} (\rt{res.users.lang}) --- đây là lần duy nhất màn đăng nhập ghi dữ liệu.
  \item Đăng xuất: đóng phiên, quay về màn đăng nhập. \textbf{Không} xoá cookie cửa hàng đang
        chọn --- lần sau đăng nhập lại vẫn còn nguyên cửa hàng cũ, kể cả khi người đăng nhập
        là tài khoản khác trên cùng máy (hệ thống vẫn kiểm lại quyền nên không lộ dữ liệu, chỉ
        là cửa hàng hiển thị có thể không như mong đợi).
\end{itemize}

\mhoi

\begin{ask}
\begin{enumerate}[leftmargin=*]
  \item \textbf{Chưa chặn thử mật khẩu hàng loạt ở màn đăng nhập.} Màn \emph{Quên mật khẩu} đã
        có chặn 10 lần/giờ/địa chỉ mạng, nhưng màn đăng nhập thì chưa. Anh có muốn đặt luật
        tương tự (ví dụ khoá tạm 5 phút sau 10 lần sai) không, và thông báo hiển thị thế nào?
  \item \textbf{Đăng xuất không xoá cửa hàng đang chọn.} Máy dùng chung ở cửa hàng, người sau
        đăng nhập sẽ thấy tên cửa hàng của người trước cho tới khi bấm đổi. Có cần xoá khi
        đăng xuất không?
\end{enumerate}
\end{ask}

\section{Màn Quên mật khẩu}
\rt{/portal/forgot-pass}
\medskip

\mvao

Liên kết ``Quên mật khẩu?'' ngay dưới ô mật khẩu ở màn đăng nhập.

\mai

Không cần đăng nhập. Ai cũng mở được.

\mhien

Một ô nhập \emph{email}, một nút gửi. Không hiển thị dữ liệu nào từ hệ thống.

\mloc

Chỉ ô email.

\mkiem{Gửi yêu cầu}

\begin{kiembang}
1 & Địa chỉ mạng chưa gửi quá \textbf{10 lần trong 1 giờ} &
    ``Quá nhiều yêu cầu. Vui lòng thử lại sau.'' \\
2 & Có nhập email &
    ``Vui lòng nhập email.'' \\
3 & \emph{(không kiểm email có tồn tại hay không --- xem giải thích dưới)} &
    --- \\
\end{kiembang}

\textbf{Cố ý không báo email có tồn tại hay không.} Dù email đúng hay sai, người dùng luôn nhận
đúng một câu: \emph{``Nếu email tồn tại trong hệ thống, hướng dẫn đặt lại mật khẩu đã được gửi
tới hộp thư của bạn.''} Đây là chủ đích --- nếu báo ``email không tồn tại'' thì kẻ xấu có thể
dò ra danh sách email đang dùng hệ thống. Lỗi thật (email không có, hoặc máy chủ thư chưa cấu
hình) chỉ được ghi vào nhật ký hệ thống cho kỹ thuật xem.

\mghi

Sinh một \textbf{mã dùng một lần} gắn vào tài khoản và gửi email chứa liên kết đặt lại. Nếu
máy chủ gửi thư chưa được cấu hình thì \textbf{không có gì được gửi đi} nhưng người dùng vẫn
thấy câu thông báo thành công.

\mhoi

\begin{ask}
Vì màn luôn báo thành công, người dùng nhập nhầm email sẽ ngồi chờ một email không bao giờ tới,
và Ngô Gia cũng không biết. Anh muốn giữ nguyên (an toàn hơn), hay thêm dòng hướng dẫn kiểu
``không nhận được sau 5 phút, vui lòng liên hệ hotline\ldots''?
\end{ask}

\section{Màn Đặt lại mật khẩu}
\rt{/portal/reset_password}
\medskip

\mvao

Bấm liên kết trong email đặt lại mật khẩu. Liên kết có kèm \textbf{mã dùng một lần}.

\mai

Không cần đăng nhập, nhưng \textbf{bắt buộc có mã hợp lệ}. Vào thẳng đường dẫn mà không có mã
thì bị đẩy về màn đăng nhập. Mã hết hạn hoặc đã dùng cũng bị đẩy về.

\mhien

Tên tài khoản (lấy ra từ chính mã), ô \emph{mật khẩu mới}, ô \emph{nhập lại mật khẩu}.

\mloc

Hai ô mật khẩu.

\mkiem{Đặt lại mật khẩu}

\begin{kiembang}
1 & Liên kết có mã & Đẩy về màn đăng nhập \\
2 & Hai ô mật khẩu khớp nhau & ``Mật khẩu xác nhận không khớp.'' \\
3 & Mật khẩu \textbf{tối thiểu 8 ký tự} & ``Mật khẩu tối thiểu 8 ký tự.'' \\
\end{kiembang}

\textbf{Chỉ có luật độ dài 8 ký tự} --- không bắt buộc chữ hoa, chữ số hay ký tự đặc biệt.

\mghi

Ghi mật khẩu mới vào tài khoản, huỷ hiệu lực của mã, rồi đưa người dùng về màn đăng nhập kèm
thông báo thành công. \textbf{Không} tự đăng nhập hộ.

\mhoi

\begin{ask}
Luật mật khẩu hiện chỉ là ``từ 8 ký tự''. Anh có yêu cầu thêm (hoa/số/ký tự đặc biệt, không
trùng mật khẩu cũ, buộc đổi sau N tháng) không? Lưu ý luật này phải áp \textbf{đồng thời} cho
màn Đổi mật khẩu ở mục 4.7.
\end{ask}

\section{Màn Trang chủ}
\rt{/portal}
\medskip

\mvao

Vào thẳng sau khi đăng nhập, hoặc bấm biểu tượng nhà trên thanh điều hướng.

\mai

\begin{itemize}[leftmargin=*]
  \item Phải đăng nhập.
  \item \textbf{Phạm vi dữ liệu}: nếu đã chọn cửa hàng thì \textbf{chỉ} cửa hàng đó; chưa chọn
        thì \textbf{cộng gộp mọi cửa hàng} tài khoản thuộc. Trên UAT có \uat{2}/7 tài khoản
        thuộc nhiều hơn một cửa hàng (một tài khoản 3 cửa hàng, một tài khoản 2 cửa hàng), nên
        hiện tượng cộng gộp này BA test bằng hai tài khoản đó.
  \item Tài khoản \textbf{không thuộc cửa hàng nào}: mọi con số về 0, mọi danh sách rỗng ---
        màn vẫn mở được, không báo lỗi.
  \item \textbf{Hộp chọn cửa hàng} tự bật lên khi tài khoản thuộc \textbf{từ 2 cửa hàng trở
        lên} \emph{và} cookie đang trống.
  \item \textbf{Vai trò không giới hạn gì ở màn này} --- Nhân viên nhìn thấy đúng bằng Chủ tiệm,
        kể cả các con số về đơn hàng và đổi trả.
  \item Toàn bộ số liệu đọc bằng quyền hệ thống; giới hạn duy nhất là điều kiện
        \rt{franchise_id} trong bảng dưới.
\end{itemize}

\mhien

\textbf{Bốn thẻ số liệu} ở đầu màn:

\begin{hienbang}
Thẻ \emph{Thông báo chưa đọc} &
  \rt{wujia.notification} trừ đi \rt{wujia.notification.read} &
  Số thông báo đã công bố (\rt{is_published_portal = true}) gửi cho \textbf{mọi cửa hàng}
  (\rt{franchise_ids} trống) \textbf{hoặc} đúng cửa hàng đang xét; rồi \textbf{trừ} đi số lượt
  người này đã đánh dấu đọc. Không âm. &
  --- &
  \uat{2} thông báo đã công bố, cả \uat{2} đều gửi toàn hệ \\ \addlinespace
Thẻ \emph{Đơn hàng 30 ngày} &
  \rt{sale.order} &
  Đơn của cửa hàng đang xét, ngày đặt trong \textbf{30 ngày} gần nhất, trạng thái khác
  \emph{Đã huỷ}. &
  --- &
  \uat{31} đơn chưa huỷ trên toàn hệ \\ \addlinespace
Thẻ \emph{Đơn chờ xác nhận} &
  \rt{sale.order} &
  Đơn của cửa hàng đang xét, trạng thái \emph{Nháp} hoặc \emph{Đã gửi}
  (\rt{state in ('draft','sent')}). Không giới hạn thời gian. &
  --- &
  \uat{17} đơn \\ \addlinespace
Thẻ \emph{Yêu cầu đổi trả đang mở} &
  \rt{wujia.return.request} &
  Yêu cầu của cửa hàng đang xét, trạng thái \emph{Đã gửi}, \emph{Đang xử lý} hoặc
  \emph{Đã duyệt} (\rt{state in ('submitted','processing','approved')}). &
  --- &
  \uat{7} yêu cầu \\
\end{hienbang}

\textbf{Các khối danh sách} bên dưới --- mọi khối đều cắt ở \textbf{2 dòng}, xem đầy đủ phải
bấm ``Xem tất cả'':

\begin{hienbang}
Thông báo mới nhất &
  \rt{wujia.notification} &
  Đã công bố \textbf{và} (gửi toàn hệ \textbf{hoặc} đúng cửa hàng đang xét) &
  Ngày công bố mới nhất trước; 2 dòng &
  \uat{2} thông báo \\ \addlinespace
Đơn hàng gần đây &
  \rt{sale.order} &
  Đơn của cửa hàng đang xét, trạng thái khác \emph{Đã huỷ} &
  Ngày đặt mới nhất trước; 2 dòng &
  \uat{31} đơn \\ \addlinespace
Yêu cầu đổi trả gần đây &
  \rt{wujia.return.request} &
  Của cửa hàng đang xét, \textbf{loại bỏ} \emph{Từ chối} và \emph{Đã huỷ}
  (\rt{state not in ('rejected','cancelled')}) --- lưu ý bộ trạng thái này
  \textbf{khác} bộ của thẻ số ở trên. &
  Ngày gửi mới nhất trước; 2 dòng &
  \uat{12} yêu cầu lọt bộ lọc này \\ \addlinespace
Sản phẩm đặt nhiều nhất &
  \rt{sale.order.line} &
  Dòng hàng thuộc đơn của cửa hàng đang xét, đơn đã \emph{Xác nhận} hoặc \emph{Hoàn tất}
  (\rt{state in ('sale','done')}), ngày đặt trong \textbf{90 ngày}. Gộp theo sản phẩm, cộng số
  lượng và thành tiền. &
  Số lượng nhiều nhất trước; 5 dòng &
  --- \\ \addlinespace
Giao hàng sắp tới \note{(bản điện thoại)} &
  \rt{stock.picking.batch} &
  Chuyến có phiếu giao thuộc cửa hàng đang xét, trạng thái \textbf{chưa hoàn thành}, và (khởi
  hành từ đầu hôm nay trở đi \textbf{hoặc} chưa đặt lịch \textbf{hoặc} đang bốc/đang chạy). &
  Giờ khởi hành sớm nhất trước; 2 dòng &
  \uat{1} chuyến chưa hoàn thành \\ \addlinespace
Bài kiến thức \note{(bản điện thoại)} &
  \rt{wujia.knowledge.article} &
  Đã công bố lên portal. \textbf{Không} lọc theo cửa hàng --- mọi cửa hàng thấy như nhau. &
  Ngày công bố mới nhất trước; 2 dòng &
  \uat{21} bài \\ \addlinespace
Dải chào ở đầu màn \note{(bản điện thoại)} &
  \rt{wujia.franchise.management} + \rt{wujia.franchise.member} &
  Tên cửa hàng đang chọn và vai trò của chính người đang xem tại cửa hàng đó. Chưa chọn cửa
  hàng thì để trống. &
  --- &
  \uat{4} cửa hàng, \uat{10} thành viên còn hiệu lực \\ \addlinespace
Dải khung giờ đặt hàng \note{(bản điện thoại)} &
  \rt{wujia.order.window} + tham số chung &
  Theo \textbf{khu vực của cửa hàng đang chọn}. Ba trạng thái: \emph{Đặt hàng 24/7} (khung giờ
  tắt), \emph{Đang mở} kèm thời gian còn lại, \emph{Đã đóng} kèm giờ mở lại. Luật đầy đủ ở
  \textbf{mục 3.7}. &
  --- &
  \uat{1} khung riêng / 2 khu vực \\
\end{hienbang}

\textbf{Ký hiệu tiền trên Trang chủ.} Số tiền gộp từ các đơn (bảng \emph{Sản phẩm đặt nhiều
nhất}, tổng tiền chuyến giao) mang ký hiệu tiền tệ \textbf{của chính các đơn đó} --- lấy từ bảng
giá gắn với khách hàng của cửa hàng, không phải tiền tệ của công ty. Nếu các đơn được gộp lại
đang dùng \textbf{nhiều loại tiền khác nhau} thì con số cộng gộp vốn đã vô nghĩa, hệ thống rơi
về tiền tệ công ty.

\mloc

Trang chủ \textbf{không có bộ lọc}. Thứ duy nhất người dùng đổi được là \emph{cửa hàng đang làm
việc} (mục 4.5) --- đổi cửa hàng thì toàn bộ số liệu trên màn đổi theo.

\mkiem{một thẻ số hoặc ``Xem tất cả''}

Không kiểm gì --- đây chỉ là liên kết sang màn danh sách của phân hệ tương ứng.

\mghi

\textbf{Trang chủ không ghi gì vào hệ thống.} Mở màn bao nhiêu lần cũng không tạo, không sửa
bản ghi nào; kể cả khối Thông báo cũng không tự đánh dấu là đã đọc.

\mhoi

\begin{ask}
\begin{enumerate}[leftmargin=*]
  \item \textbf{Số ``Thông báo chưa đọc'' ở Trang chủ có thể lệch với số trên màn Thông báo.}
        Ba khác biệt trong cách đếm: (a) Trang chủ \textbf{không} loại thông báo \emph{đã hết
        hiệu lực}, màn Thông báo thì có; (b) Trang chủ \textbf{không} loại thông báo có ngày
        công bố ở tương lai; (c) Trang chủ trừ \textbf{mọi} lượt đánh dấu đọc của người đó,
        trong khi trạng thái đọc được lưu \textbf{theo từng cửa hàng} --- người làm ở 2 cửa hàng
        đọc cùng một thông báo ở cả hai nơi sẽ bị trừ 2. Trên UAT tài khoản
        \emph{Administrator} thuộc 2 cửa hàng và đang có \uat{3} lượt đọc trên \uat{2} thông
        báo, nên con số bị kẹp về 0. Anh chốt giúp: lấy cách đếm của màn Thông báo làm chuẩn
        cho cả Trang chủ, đúng không?
  \item \textbf{Thẻ ``Đổi trả đang mở'' bỏ sót trạng thái \emph{Đang xét}.} Trạng thái
        \rt{reviewing} có thật trong hệ thống nhưng không nằm trong bộ đếm
        (\rt{submitted}, \rt{processing}, \rt{approved}). Trên UAT hiện chưa có yêu cầu nào ở
        trạng thái này nên chưa ai thấy. Có phải \emph{Đang xét} cũng là ``đang mở'' không?
  \item \textbf{Thẻ số và danh sách bên dưới dùng hai bộ trạng thái khác nhau.} Thẻ đếm
        \emph{Đã gửi + Đang xử lý + Đã duyệt}; danh sách hiện \emph{mọi trạng thái trừ Từ chối
        và Đã huỷ} --- nghĩa là danh sách \textbf{có} cả phiếu \emph{Nháp} và \emph{Hoàn thành}
        mà thẻ không đếm. Trên UAT: thẻ ra \uat{7}, danh sách lọc ra \uat{12}. Anh muốn thống
        nhất theo bộ nào?
  \item \textbf{Mọi khối danh sách chỉ hiện 2 dòng.} Con số 2 đang cố định trong code, không
        cấu hình được. Trên màn hình máy tính rộng, 2 dòng trông khá trống. Anh muốn giữ 2, hay
        đổi (ví dụ 5 dòng cho máy tính, 2 dòng cho điện thoại)?
  \item \textbf{Chưa chọn cửa hàng thì các con số là tổng của mọi cửa hàng.} Người quản lý nhiều
        cửa hàng mở Trang chủ lần đầu sẽ thấy số cộng gộp mà không có nhãn nào nói điều đó. Anh
        muốn (a) giữ nguyên nhưng thêm dòng ``Đang xem: tất cả 3 cửa hàng'', hay (b) bắt buộc
        chọn cửa hàng mới cho xem số?
  \item \textbf{Khối ``Bài kiến thức'' không lọc theo cửa hàng.} Mọi cửa hàng thấy chung
        \uat{21} bài. Nếu sau này có tài liệu riêng theo khu vực/cấp cửa hàng thì cần bổ sung
        điều kiện lọc --- anh xác nhận trước mắt là dùng chung?
\end{enumerate}
\end{ask}

\section{Chọn cửa hàng đang làm việc}
\rt{/portal/franchise/switch} · \rt{/portal/franchise/clear}
\medskip

\mvao

Bấm tên cửa hàng trên thanh trên cùng, hoặc chọn trong hộp tự bật ở Trang chủ. Hai đường dẫn
này \textbf{không phải màn} --- chúng chỉ nhận lệnh rồi đưa người dùng quay lại trang cũ.

\mai

Phải đăng nhập. Chỉ chọn được cửa hàng mà tài khoản đang có thành viên \textbf{còn hiệu lực}.

\mhien

Danh sách cửa hàng trong hộp chọn lấy từ \rt{wujia.franchise.management}, giới hạn đúng những
cửa hàng tài khoản thuộc về.

\mloc

Người dùng chọn đúng một cửa hàng, hoặc bấm ``Bỏ chọn''.

\mkiem{Chọn cửa hàng}

\begin{kiembang}
1 & Mã cửa hàng gửi lên là số & Đưa về Trang chủ, không đổi gì \\
2 & Mã đó nằm trong danh sách cửa hàng tài khoản còn hiệu lực &
    Đưa về Trang chủ, \textbf{không có thông báo lỗi} --- người dùng chỉ thấy cửa hàng không
    đổi \\
3 & Địa chỉ quay lại phải là \textbf{đường dẫn nội bộ} &
    Địa chỉ trỏ ra ngoài (trang web khác) bị thay bằng \rt{/portal}; đây là lớp chặn chuyển
    hướng sang trang giả mạo \\
\end{kiembang}

\mghi

\begin{itemize}[leftmargin=*]
  \item Ghi \textbf{cookie} \rt{wujia_active_franchise_id} trên máy người dùng, hạn \textbf{30
        ngày}. \emph{Không} ghi gì vào cơ sở dữ liệu --- đây là lý do đổi máy là mất lựa chọn.
  \item ``Bỏ chọn'' xoá cookie và đưa về Trang chủ.
  \item Cookie \textbf{đọc được bằng mã trình duyệt} (có chủ đích, để hiện tên cửa hàng ở thanh
        trên mà không phải hỏi máy chủ).
\end{itemize}

\mhoi

\begin{ask}
\begin{enumerate}[leftmargin=*]
  \item \textbf{Chọn cửa hàng không được phép thì im lặng.} Người dùng bấm vào cửa hàng vừa bị
        gỡ quyền sẽ chỉ thấy trang tải lại mà không đổi gì, không có câu nào giải thích. Anh
        muốn thêm thông báo ``Bạn không còn quyền tại cửa hàng này'' không?
  \item \textbf{Cửa hàng đang chọn lưu trên trình duyệt.} Dùng máy khác, trình duyệt ẩn danh
        hoặc xoá dữ liệu duyệt web là phải chọn lại từ đầu. Anh xác nhận chấp nhận cách này,
        hay muốn nhớ theo tài khoản (lưu vào hồ sơ người dùng)?
\end{enumerate}
\end{ask}

\section{Màn Hồ sơ cá nhân}
\rt{/portal/profile} · \rt{/portal/profile/update} · \rt{/portal/profile/avatar/<int:user_id>}
\medskip

\mvao

Bấm ảnh đại diện ở góc trên bên phải $\rightarrow$ \emph{Hồ sơ}.
\rt{/portal/profile/update} là hành động Lưu; \rt{/portal/profile/avatar/<id>} là đường dẫn
\textbf{trả về ảnh đại diện} --- mọi chỗ trong portal hiện ảnh người dùng đều gọi qua đây.

\mai

\begin{itemize}[leftmargin=*]
  \item Phải đăng nhập. Mỗi người chỉ xem và sửa \textbf{hồ sơ của chính mình} --- không có
        đường dẫn nào cho xem hồ sơ người khác.
  \item \textbf{Ảnh đại diện thì khác}: xem được ảnh của \textbf{chính mình}, và ảnh của bất kỳ
        ai \textbf{cùng thuộc một cửa hàng} với mình. Người ngoài cửa hàng bị từ chối. Với dữ
        liệu UAT, tài khoản thuộc cửa hàng Cầu Giấy thấy được ảnh của \uat{5} người cùng cửa hàng
        đó (kể cả chính mình).
  \item Khối ``Cửa hàng / Vai trò'' trên màn chỉ hiện khi \textbf{đã chọn được cửa hàng}; người
        nhiều cửa hàng chưa chọn thì hai dòng này hiện dấu ``---''.
\end{itemize}

\mhien

\begin{hienbang}
Họ tên, email liên hệ, điện thoại, địa chỉ &
  \rt{res.partner} (liên hệ gắn với tài khoản) &
  Chính liên hệ của người đang đăng nhập &
  --- & --- \\ \addlinespace
Ảnh đại diện &
  \rt{res.users} &
  Bản ảnh cỡ 256 điểm ảnh. Trình duyệt được phép giữ lại \textbf{5 phút} nên đổi ảnh xong có
  thể phải tải lại trang mới thấy. &
  --- & --- \\ \addlinespace
Cửa hàng + Vai trò &
  \rt{wujia.franchise.member} &
  Thành viên \textbf{còn hiệu lực} của người này tại cửa hàng đang chọn. Chưa chọn và thuộc
  nhiều cửa hàng $\rightarrow$ để trống. &
  --- &
  \uat{10} thành viên còn hiệu lực (\uat{3} Chủ tiệm, \uat{3} Quản lý, \uat{4} Nhân viên) \\
\end{hienbang}

\mloc

Bốn ô nhập: \emph{họ tên}, \emph{email}, \emph{điện thoại}, \emph{địa chỉ}; và một ô chọn ảnh.

\mkiem{Lưu}

\begin{kiembang}
1 & Họ tên không được trống & ``Họ tên không được trống'' \\
2 & Điện thoại (nếu có nhập) đúng khuôn: 7--20 ký tự, chỉ số và các dấu \rt{+ ( ) . -} và
    khoảng trắng & ``Số điện thoại không hợp lệ'' \\
3 & \emph{(ảnh, nếu có chọn)} kích thước không quá \textbf{2MB} & ``Ảnh đại diện vượt quá 2MB'' \\
4 & \emph{(ảnh)} định dạng thuộc \textbf{PNG, JPEG, WEBP} & ``Định dạng ảnh không hỗ trợ'' \\
\end{kiembang}

\textbf{Email không bị kiểm khuôn} ở tầng này --- việc kiểm do Odoo làm khi ghi; sai thì hiện
thông báo do Odoo sinh ra.

\mghi

\begin{itemize}[leftmargin=*]
  \item Ghi \emph{họ tên}, \emph{điện thoại}, \emph{địa chỉ} vào liên hệ của tài khoản.
  \item \textbf{Email chỉ ghi vào liên hệ, KHÔNG đổi email đăng nhập.} Người dùng sửa email ở
        đây rồi lần sau vẫn phải đăng nhập bằng email cũ. Mỗi lần sửa đều được ghi vào nhật ký
        hệ thống.
  \item Ảnh (nếu có) ghi vào tài khoản ở cỡ gốc; Odoo tự sinh các cỡ nhỏ.
  \item Lưu xong hệ thống \textbf{chuyển hướng} về lại màn hồ sơ kèm câu ``Cập nhật thành công''
        --- cách này để bấm F5 không lưu lại lần nữa.
\end{itemize}

\mhoi

\begin{ask}
\begin{enumerate}[leftmargin=*]
  \item \textbf{Sửa email ở Hồ sơ không đổi được email đăng nhập} --- đây là điểm gần như chắc
        chắn gây hiểu nhầm. Anh muốn (a) đổi nhãn ô thành ``Email liên hệ'' và ghi chú rõ, hay
        (b) cho đổi luôn email đăng nhập (phải có bước xác nhận qua email mới)?
  \item \textbf{Ai cùng cửa hàng đều xem được ảnh đại diện của nhau.} Xác nhận đúng ý anh chứ?
  \item \textbf{Hồ sơ chưa cho sửa gì khác} --- không có ngày sinh, không có chức danh, không
        có liên kết tới hồ sơ nhân sự. Anh cần thêm trường nào không?
\end{enumerate}
\end{ask}

\section{Màn Đổi mật khẩu}
\rt{/portal/change-password}
\medskip

\mvao

Bấm ảnh đại diện góc trên bên phải $\rightarrow$ \emph{Đổi mật khẩu}.

\mai

Phải đăng nhập. Chỉ đổi được mật khẩu của chính mình. Không có đường dẫn nào cho Chủ tiệm đặt
lại mật khẩu cho nhân viên --- việc đó phải làm trong giao diện quản trị.

\mhien

Ba ô mật khẩu, kèm khối ``Cửa hàng / Vai trò'' giống màn Hồ sơ.

\mloc

\emph{Mật khẩu hiện tại}, \emph{mật khẩu mới}, \emph{nhập lại mật khẩu mới}.

\mkiem{Đổi mật khẩu}

\begin{kiembang}
1 & Có nhập cả mật khẩu cũ và mật khẩu mới & ``Vui lòng nhập đầy đủ.'' \\
2 & Hai ô mật khẩu mới khớp nhau & ``Mật khẩu xác nhận không khớp với mật khẩu mới.'' \\
3 & Mật khẩu mới \textbf{tối thiểu 8 ký tự} & ``Mật khẩu mới tối thiểu 8 ký tự.'' \\
4 & Mật khẩu mới \textbf{khác} mật khẩu cũ & ``Mật khẩu mới phải khác mật khẩu cũ.'' \\
5 & Mật khẩu hiện tại \textbf{đúng} &
    ``Mật khẩu hiện tại không đúng. Vui lòng kiểm tra lại.'' \\
\end{kiembang}

Lưu ý thứ tự: hệ thống chỉ kiểm mật khẩu cũ ở \textbf{bước cuối}, sau khi mật khẩu mới đã qua
hết các luật khuôn dạng.

\mghi

Ghi mật khẩu mới vào tài khoản. \textbf{Phiên đang mở vẫn giữ nguyên} --- người dùng không bị
đăng xuất, và các thiết bị khác đang đăng nhập cũng không bị đẩy ra.

\mhoi

\begin{ask}
\begin{enumerate}[leftmargin=*]
  \item \textbf{Đổi mật khẩu không đăng xuất các thiết bị khác.} Nếu ai đó đổi mật khẩu vì nghi
        bị lộ, thiết bị của kẻ kia \textbf{vẫn đang đăng nhập}. Anh có muốn đóng mọi phiên khác
        khi đổi mật khẩu không?
  \item Luật mật khẩu ở đây (8 ký tự, khác mật khẩu cũ) \textbf{không giống hoàn toàn} màn Đặt
        lại mật khẩu qua email (chỉ 8 ký tự, không kiểm khác cũ). Cần thống nhất --- anh chốt
        bộ luật chung.
\end{enumerate}
\end{ask}

\section{Đổi ngôn ngữ}
\rt{/portal/set-lang/<string:lang_code>}
\medskip

\mvao

Nút cờ/ngôn ngữ trên thanh trên cùng và ở màn đăng nhập. Không phải màn --- bấm xong quay lại
đúng trang đang đứng.

\mai

\textbf{Không cần đăng nhập} --- có chủ đích, để khách đổi được ngôn ngữ ngay ở màn đăng nhập.

\mhien

Danh sách ngôn ngữ lấy từ \rt{res.lang} với điều kiện \rt{active = true}. Trên UAT đang bật
\uat{3} ngôn ngữ: Tiếng Việt, English (US), Thai.

\mloc

Chọn một ngôn ngữ.

\mkiem{một ngôn ngữ}

\begin{kiembang}
1 & Mã ngôn ngữ có trong danh sách \textbf{đang bật} &
    Không đổi gì, vẫn quay lại trang cũ (im lặng bỏ qua) \\
2 & Địa chỉ quay lại là \textbf{đường dẫn nội bộ cùng máy chủ} &
    Bị thay bằng \rt{/portal} --- lớp chặn chuyển hướng sang trang giả mạo \\
\end{kiembang}

\mghi

\begin{itemize}[leftmargin=*]
  \item Cập nhật ngôn ngữ cho \textbf{phiên làm việc hiện tại} (có hiệu lực ngay ở trang kế).
  \item \textbf{Đã đăng nhập} $\rightarrow$ ghi luôn vào tài khoản (\rt{res.users.lang}), nên
        lần sau đăng nhập vẫn giữ ngôn ngữ đó.
  \item \textbf{Chưa đăng nhập} $\rightarrow$ chỉ nhớ tạm trong phiên, và được ghi vào tài
        khoản ngay sau khi đăng nhập thành công.
\end{itemize}

\mhoi

\begin{ask}
Portal đang bật \uat{3} ngôn ngữ, nghĩa là nút đổi ngôn ngữ hiện đủ 3 lựa chọn. Nếu bản dịch
tiếng Anh/tiếng Thái chưa đầy đủ thì người dùng chọn vào sẽ thấy màn lẫn lộn hai thứ tiếng. Anh
muốn tạm tắt bớt ngôn ngữ chưa dịch xong không?
\end{ask}

\section{Màn Danh sách cửa hàng của tôi}
\rt{/portal/franchises} · \rt{/my/franchises}
\medskip

\mvao

Menu tài khoản $\rightarrow$ \emph{Cửa hàng của tôi}. Hai đường dẫn cho \textbf{cùng một nội
dung}: bản \rt{/portal/\ldots} dùng giao diện Vuexy của dự án, bản \rt{/my/\ldots} dùng giao
diện portal mặc định của Odoo. Xem mục 7.

\mai

Phải đăng nhập. Chỉ hiện những cửa hàng mà tài khoản có thành viên \textbf{còn hiệu lực}.
Vai trò không giới hạn gì --- Nhân viên cũng thấy đủ.

\mhien

\begin{hienbang}
Danh sách cửa hàng &
  \rt{wujia.franchise.member} $\rightarrow$ \rt{wujia.franchise.management} &
  Thành viên của chính người đang đăng nhập, \rt{is_currently_valid = true}. ``Còn hiệu lực''
  nghĩa là: bản ghi chưa lưu trữ \textbf{và} đang làm việc \textbf{và} hôm nay nằm trong khoảng
  \emph{từ ngày} -- \emph{đến ngày}. &
  Theo thứ tự mặc định của bảng thành viên (cửa hàng, chủ chính trước, ngày bắt đầu mới trước)
  &
  \uat{10} bản ghi thành viên, \uat{0} bản hết hiệu lực \\ \addlinespace
Nhãn vai trò trên từng dòng &
  \rt{wujia.franchise.member.role} &
  \emph{Chủ tiệm} / \emph{Quản lý} / \emph{Nhân viên} &
  --- &
  \uat{3} / \uat{3} / \uat{4} \\
\end{hienbang}

\mloc

Không có bộ lọc, không phân trang. Một người thuộc bao nhiêu cửa hàng thì hiện bấy nhiêu dòng
(nhiều nhất trên UAT là \uat{3}).

\mkiem{một cửa hàng}

Không kiểm gì --- sang màn chi tiết.

\mghi

Không ghi gì.

\mhoi

\begin{ask}
\textbf{Đang có hai bộ màn song song cho cùng một nội dung}: \rt{/portal/franchises} (giao diện
Vuexy) và \rt{/my/franchises} (giao diện Odoo mặc định), tương tự ở màn chi tiết. Hai bộ này lấy
cùng dữ liệu nhưng khác giao diện, và \textbf{không} có liên kết nào trong portal trỏ tới bộ
\rt{/my/\ldots}. Anh cho biết: giữ cả hai (bộ \rt{/my/} để nhân sự Ngô Gia đối chiếu), hay gỡ
hẳn bộ \rt{/my/}? \note{Gỡ thì cần giữ chuyển hướng để liên kết cũ không chết.}
\end{ask}

\section{Màn Chi tiết cửa hàng}
\rt{/portal/franchises/<int:franchise_id>} · \rt{/my/franchises/<int:franchise_id>} ·
\rt{/portal/franchises/<int:franchise_id>/profile} ·
\rt{/my/franchises/<int:franchise_id>/members}
\medskip

\mvao

Bấm một dòng ở màn Danh sách cửa hàng. Đường dẫn thứ ba (\rt{/profile}) là bản \emph{hồ sơ đầy
đủ} của cửa hàng; đường dẫn thứ tư \textbf{không phải màn} mà là nguồn dữ liệu danh sách thành
viên cho giao diện gọi ngầm.

\mai

\begin{itemize}[leftmargin=*]
  \item Phải đăng nhập \textbf{và} có thành viên còn hiệu lực tại \textbf{đúng} cửa hàng đó.
        Không đạt $\rightarrow$ bị đưa về danh sách cửa hàng, \textbf{không báo lỗi} (cố ý
        không cho biết cửa hàng đó có tồn tại hay không).
  \item \textbf{Đây là điểm quan trọng về quyền riêng tư}: bước kiểm ở trên chỉ hỏi ``bạn có
        thuộc cửa hàng này không''. Sau khi qua được, danh sách \textbf{tất cả thành viên} của
        cửa hàng được đọc bằng quyền hệ thống. Nghĩa là \textbf{Nhân viên cũng xem được đầy đủ
        danh sách đồng nghiệp}, gồm tên, vai trò, số điện thoại, ngày bắt đầu và ngày kết thúc.
        Xem mục 7.
  \item Vai trò \textbf{không} giới hạn gì trong màn này.
\end{itemize}

\mhien

\begin{hienbang}
Khối thông tin cửa hàng &
  \rt{wujia.franchise.management} &
  Đúng cửa hàng được mở &
  --- &
  \uat{4} cửa hàng, cả \uat{4} đều \emph{Đang hoạt động} \\ \addlinespace
Khối thành viên &
  \rt{wujia.franchise.member} &
  Mọi thành viên của cửa hàng đó, \rt{is_currently_valid = true}. \textbf{Không} lọc theo người
  đang xem. &
  Theo thứ tự mặc định của bảng &
  Cầu Giấy \uat{5} người · TP HCM Quận 1 \uat{3} · Đống Đa \uat{1} · TEST \uat{1} \\ \addlinespace
Trạng thái cửa hàng &
  \rt{wujia.franchise.management.status} &
  Nhãn tiếng Việt cố định trong portal: \emph{Nháp} / \emph{Đang hoạt động} / \emph{Khoá} /
  \emph{Đã đóng} / \emph{Hết hạn} &
  --- &
  \uat{4}/4 \emph{Đang hoạt động} \\ \addlinespace
Hồ sơ đầy đủ \note{(đường dẫn /profile)} &
  \rt{wujia.franchise.management} + \rt{wujia.franchise.member} &
  Thêm các thông tin hợp đồng, khu vực, ngày bắt đầu/kết thúc nhượng quyền --- vẫn chỉ của cửa
  hàng người dùng thuộc về &
  --- & --- \\
\end{hienbang}

\mloc

Không có bộ lọc, không phân trang trên các màn này (khác với màn \emph{Thông tin cửa hàng} ở
mục 4.11 --- màn đó \textbf{có} phân trang).

\mkiem{một liên kết}

Không kiểm gì.

\mghi

Không ghi gì.

\mhoi

\begin{ask}
\begin{enumerate}[leftmargin=*]
  \item \textbf{Nhân viên xem được đầy đủ danh sách đồng nghiệp}, kèm số điện thoại và ngày
        hiệu lực. Anh xác nhận đúng mong muốn, hay muốn giới hạn: chỉ \emph{Chủ tiệm} và
        \emph{Quản lý} thấy đủ, còn \emph{Nhân viên} chỉ thấy tên và vai trò?
  \item \textbf{Bốn đường dẫn cho ba màn gần giống nhau.} \rt{/portal/franchises/<id>},
        \rt{/my/franchises/<id>} và \rt{/portal/franchises/<id>/profile} hiển thị phần lớn cùng
        thông tin, cộng thêm màn \emph{Thông tin cửa hàng} ở mục 4.11 nữa là \textbf{bốn} chỗ
        xem thông tin cửa hàng. Anh chốt giúp giữ lại \textbf{một} màn chính là màn nào; các
        đường dẫn còn lại sẽ chuyển hướng về đó.
\end{enumerate}
\end{ask}

\section{Màn Thông tin cửa hàng}
\rt{/portal/franchise-information}
\medskip

\mvao

Menu \emph{Thông tin cửa hàng} trên thanh điều hướng. Màn này luôn nói về \textbf{cửa hàng đang
chọn}, không nhận mã cửa hàng trên đường dẫn.

\mai

\begin{itemize}[leftmargin=*]
  \item Phải đăng nhập \textbf{và} đã chọn được cửa hàng. Chưa chọn $\rightarrow$ bị đưa về
        Trang chủ.
  \item Phải có thành viên còn hiệu lực tại cửa hàng đó.
  \item \textbf{Cổng chặn cứng}: cửa hàng \emph{bị khoá portal} (\rt{portal_locked = true})
        hoặc \textbf{không} ở trạng thái \emph{Đang hoạt động} (\rt{status != 'active'})
        $\rightarrow$ hiện màn ``bị khoá'' thay cho nội dung. Trên UAT hiện \uat{0} cửa hàng bị
        khoá và \uat{0} cửa hàng ngoài trạng thái \emph{Đang hoạt động}, nên \textbf{nhánh này
        chưa từng được kiểm thử bằng dữ liệu thật}.
  \item Vai trò không giới hạn gì.
\end{itemize}

\mhien

\begin{hienbang}
Khối hồ sơ cửa hàng &
  \rt{wujia.franchise.management} &
  Cửa hàng đang chọn &
  --- & --- \\ \addlinespace
Khối thành viên &
  \rt{wujia.franchise.member} &
  Thành viên của cửa hàng đang chọn, \rt{is_currently_valid = true} &
  Theo \textbf{vai trò} rồi mã bản ghi; \textbf{10} người mỗi trang &
  Cầu Giấy \uat{5} người --- chưa cửa hàng nào chạm ngưỡng phân trang \\ \addlinespace
Khối thành viên của chính tôi &
  \rt{wujia.franchise.member} &
  Bản ghi thành viên của người đang xem tại cửa hàng này &
  --- & --- \\
\end{hienbang}

\mloc

\begin{itemize}[leftmargin=*]
  \item \textbf{Số người mỗi trang} --- chỉ nhận \textbf{10, 20, 50}; gửi số khác thì lặng lẽ
        dùng 10.
  \item \textbf{Trang} --- vượt quá trang cuối thì tự kéo về trang cuối.
\end{itemize}

\mkiem{---}

Màn \textbf{chỉ xem}, không có nút ghi nào. Việc đề nghị sửa thông tin cửa hàng nằm ở phân hệ
\emph{Yêu cầu thông tin} (chương 7).

\mghi

Không ghi gì.

\mhoi

\begin{ask}
\begin{enumerate}[leftmargin=*]
  \item \textbf{Sắp xếp thành viên theo vai trò đang là thứ tự chữ cái của mã trạng thái}
        (\rt{manager} $\rightarrow$ \rt{owner} $\rightarrow$ \rt{staff}), tức ra \emph{Quản lý
        $\rightarrow$ Chủ tiệm $\rightarrow$ Nhân viên} --- không phải thứ tự cấp bậc. Anh muốn
        đổi thành \emph{Chủ tiệm $\rightarrow$ Quản lý $\rightarrow$ Nhân viên} không?
  \item \textbf{Nhánh ``cửa hàng bị khoá'' chưa có dữ liệu để test.} Nếu anh muốn nghiệm thu
        nhánh này, cần Ngô Gia khoá thử một cửa hàng trên UAT --- em không tự đổi dữ liệu.
\end{enumerate}
\end{ask}

\section{Bốn đường dẫn cũ đã đổi tên}
\rt{/my/branches} · \rt{/my/branches/<int:branch_id>} · \rt{/portal/branches} ·
\rt{/portal/branches/<int:branch_id>}
\medskip

Trước đây cửa hàng nhượng quyền được gọi là ``branch''. Khi đổi sang ``franchise'', bốn đường
dẫn cũ được \textbf{giữ lại làm đường chuyển hướng vĩnh viễn} sang đường dẫn mới tương ứng, để
liên kết cũ (trong email, tin nhắn, dấu trang của người dùng) không bị chết.

Bốn đường dẫn này \textbf{không hiện màn nào}, không đọc và không ghi dữ liệu. BA không cần lên
task cho chúng; chỉ cần biết là chúng tồn tại khi nhìn danh sách đường dẫn.

\section{Trang xem trước bộ giao diện}
\rt{/portal/_pc-preview}
\medskip

Trang \textbf{dành riêng cho đội phát triển}: bày sẵn mọi thành phần giao diện (nút, thẻ, bảng,
phân trang, nhãn trạng thái) để đối chiếu với bản thiết kế. \textbf{Chỉ nhân sự nội bộ mở
được} --- tài khoản cửa hàng vào sẽ bị đưa về Trang chủ. Không có liên kết nào trong portal trỏ
tới trang này, và nó không đọc dữ liệu nghiệp vụ nào.

\section{Quy ước dùng chung toàn portal}
\label{sec:quy-uoc-chung}

Những quy tắc dưới đây được viết \textbf{một lần duy nhất} trong mã nguồn và mọi màn của mọi
phân hệ đều dùng chung. Các chương sau sẽ trỏ về mục này thay vì lặp lại.

\subsection*{Phân trang}

\begin{itemize}[leftmargin=*]
  \item Số bản ghi mỗi trang \textbf{chỉ nhận đúng các giá trị có trong ô chọn} của màn đó
        (thường là 10 / 20 / 50, riêng màn Đặt hàng là 24 / 48 / 96). Gửi số khác --- kể cả sửa
        tay trên thanh địa chỉ --- thì hệ thống lặng lẽ dùng giá trị mặc định, \textbf{không
        báo lỗi}. Trần cứng là 100 bản ghi/trang.
  \item Số trang vượt quá trang cuối thì \textbf{tự kéo về trang cuối}.
  \item \textbf{Đổi trang luôn giữ nguyên mọi bộ lọc}: từ khoá, khoảng ngày, trạng thái đang
        chọn đều được mang theo. Đây là lý do các màn không còn hiện tượng ``bấm trang 2 là mất
        bộ lọc''.
  \item Thanh phân trang hiện dạng \emph{1 \ldots 4 5 6 \ldots 20} --- luôn có trang đầu, trang
        cuối, và một trang liền kề mỗi bên.
\end{itemize}

\subsection*{Ngày giờ và múi giờ}

\begin{itemize}[leftmargin=*]
  \item Hệ thống lưu giờ chuẩn quốc tế; khi hiện lên mới quy đổi theo ô \emph{Múi giờ} của
        \textbf{tài khoản người dùng}. Ô trống thì dùng \emph{Asia/Ho\_Chi\_Minh}. Múi giờ ghi
        sai (giá trị lạ) cũng rơi về mặc định chứ \textbf{không} làm lỗi trang.
  \item Trên UAT: \uat{1} tài khoản để trống múi giờ, \uat{6} tài khoản đặt giá trị khác
        \emph{Asia/Ho\_Chi\_Minh} --- cụ thể \emph{Asia/Saigon} (4 tài khoản) và
        \emph{Asia/Bangkok} (2 tài khoản). Cả ba giá trị này hiện đều là UTC+7 nên chưa gây
        lệch giờ; nhưng nếu sau này có người dùng ở múi giờ khác thì \textbf{mọi màn lọc theo
        ngày sẽ lệch theo tài khoản đó}, không theo cửa hàng.
  \item \textbf{Lọc theo khoảng ngày} được quy đổi đúng biên: ``từ ngày X'' tính từ 00:00 giờ
        địa phương, ``đến ngày Y'' tính hết 23:59:59 giờ địa phương.
  \item Khoảng ngày ngược (\emph{Từ ngày} sau \emph{Đến ngày}) dùng chung một câu báo trên mọi
        màn: \emph{``Từ ngày không được lớn hơn Đến ngày''}. Ngày sai định dạng thì tuỳ màn:
        có màn coi như bỏ lọc, có màn báo riêng.
\end{itemize}

\subsection*{Định dạng tiền}

\begin{itemize}[leftmargin=*]
  \item Một chỗ duy nhất sinh ra chuỗi tiền cho cả portal. Tiền Việt hiện dạng
        \emph{48.000 đ} (không phần lẻ, nhóm nghìn bằng dấu chấm); tiền có phần lẻ hiện dạng
        \emph{10.990,99 \$} --- số chữ số lẻ và ký hiệu đều lấy theo \textbf{loại tiền của
        chứng từ}, không cố định trong giao diện.
  \item Với các con số \textbf{cộng gộp nhiều đơn}, ký hiệu tiền lấy theo loại tiền của chính
        các đơn đó; trộn nhiều loại tiền thì rơi về tiền tệ công ty (và con số gộp khi đó vốn
        đã không có ý nghĩa).
\end{itemize}

\subsection*{Nhãn trạng thái (màu badge)}

\begin{itemize}[leftmargin=*]
  \item Màu của nhãn trạng thái \textbf{do chữ trên nhãn quyết định}, không do từng màn tự
        chọn --- nên cùng một chữ (ví dụ \emph{Đã xác nhận}) ra cùng một màu ở mọi màn, trên cả
        máy tính lẫn điện thoại. Đây là kết quả xử lý lỗi \emph{UI-STATUSBADGE-001}.
  \item Bảy bậc màu: \emph{trung tính} (chưa bắt đầu / đã khép lại) · \emph{thông tin} (mới, đã
        xác nhận) · \emph{chờ} (đang chờ bên khác) · \emph{đang chạy} · \emph{thành công} ·
        \emph{hỏng/dừng} · \emph{cần bạn làm gì đó}.
  \item Nhãn lạ chưa nằm trong bảng thì ra màu trung tính, \textbf{không bao giờ mất nhãn}.
  \item \textbf{Bản điện thoại có một số nhãn khác bản máy tính} theo đúng thiết kế đã duyệt ---
        ví dụ ticket hỗ trợ ở trạng thái chờ khách trả lời: điện thoại ghi \emph{``Có phản
        hồi''}, máy tính ghi \emph{``Chờ phản hồi''}. Đây là khác biệt \textbf{có chủ đích}
        theo bản thiết kế, xem mục 7 chương 7.
\end{itemize}

\subsection*{Tệp đính kèm}

\begin{itemize}[leftmargin=*]
  \item Mọi màn cho tải tệp lên (đổi trả, hỗ trợ, yêu cầu thông tin) dùng chung một bộ luật:
        tối đa \textbf{10 tệp}, mỗi tệp tối đa \textbf{5MB}, và \textbf{định dạng được kiểm ở
        máy chủ} chứ không tin khai báo của trình duyệt. Màn ảnh chỉ nhận PNG/JPEG/WEBP; màn tài
        liệu nhận thêm PDF.
  \item Tên tệp được làm sạch trước khi lưu (chống tên tệp độc hại).
  \item \textbf{Xem lại tệp đã tải lên}: hệ thống kiểm người xem có thuộc cửa hàng của chứng từ
        gốc không; không thuộc thì từ chối.
\end{itemize}

\subsection*{Vai trò}

Thang bậc \emph{Nhân viên (1) $<$ Quản lý (2) $<$ Chủ tiệm (3)}. Chỗ nào yêu cầu ``ít nhất
Quản lý'' thì Chủ tiệm đương nhiên qua. Người thuộc nhiều cửa hàng thì lấy \textbf{bậc cao
nhất}. Không đủ bậc $\rightarrow$ hệ thống từ chối thẳng với mã lỗi ``không có quyền'', không
phải chuyển hướng êm.

\subsection*{Giới hạn số lần gọi}

Chỉ \textbf{một} chỗ trong portal đang đặt giới hạn: màn \emph{Quên mật khẩu} --- 10 lần/giờ cho
mỗi địa chỉ mạng. Bộ đếm nằm trong bộ nhớ của tiến trình máy chủ, nghĩa là khởi động lại máy chủ
là bộ đếm về 0.

\chapter{Lịch sử mua hàng --- Giao hàng --- Báo cáo}

Chương này gồm \textbf{11 đường dẫn} gom thành \textbf{8 màn}, thuộc ba mô-đun:
\rt{wujia_portal_purchase_history} (5), \rt{wujia_portal_delivery} (4),
\rt{wujia_portal_report} (2). Ba phân hệ này đều \textbf{chỉ đọc} --- không màn nào tạo hay sửa
dữ liệu nghiệp vụ, trừ việc sinh tệp tải về.

\textbf{Ba điểm cần biết trước khi đọc:}

\begin{enumerate}[leftmargin=*]
  \item \textbf{Ba phân hệ xác định ``cửa hàng đang xem'' theo hai cách khác nhau.} Lịch sử mua
        hàng \textbf{bắt buộc} phải chọn được cửa hàng, chưa chọn thì báo lỗi và không hiện gì.
        Giao hàng và Báo cáo thì \textbf{cộng gộp mọi cửa hàng} khi chưa chọn. Đây là khác biệt
        có thật trong code, không phải mô tả nhầm --- xem mục 7 của màn Lịch sử.
  \item \textbf{Trạng thái đơn mà cửa hàng nhìn thấy là trạng thái GỘP.} Odoo chỉ có bốn trạng
        thái đơn (\emph{Nháp}, \emph{Đã gửi}, \emph{Đã xác nhận}, \emph{Đã huỷ}). Portal thêm
        hai nhãn \emph{Đang giao} và \emph{Hoàn tất} \textbf{suy ra từ chuyến giao} và cho đè
        lên nhãn \emph{Đã xác nhận}. Đây là phương án (a) chủ dự án chốt 03/08: cửa hàng chỉ
        nhìn một cột trạng thái duy nhất.
  \item \textbf{Đơn Đã huỷ bị giấu khỏi phân hệ Lịch sử.} Trên UAT hiện có \uat{16} đơn đã huỷ
        trên tổng \uat{47} --- tức \textbf{một phần ba số đơn không xuất hiện} ở màn Lịch sử.
        Nhưng chúng \textbf{vẫn có} trong tệp Excel của phân hệ Báo cáo. Xem mục 7 màn Báo cáo.
\end{enumerate}

\section{Màn Lịch sử mua hàng}
\rt{/portal/purchase-history} · \rt{/portal/purchase-history/results}
\medskip

\mvao

Mục \emph{Lịch sử mua hàng} trên thanh điều hướng. Đường dẫn \rt{/results} không phải màn
riêng --- đó là phần kết quả được gọi ngầm khi người dùng đổi bộ lọc, để không phải tải lại cả
trang. Hai đường dẫn dùng chung một bộ dữ liệu nên mọi luật dưới đây đúng cho cả hai.

Ngoài ra có \textbf{một đường dẫn cũ} \rt{/portal/purchase_history} (gạch dưới, từ bản v14) được
giữ lại để chuyển hướng vĩnh viễn sang đường dẫn mới --- liên kết cũ trong email hay dấu trang
của cửa hàng không bị chết.

\mai

\begin{itemize}[leftmargin=*]
  \item Phải đăng nhập \textbf{và phải chọn được cửa hàng}. Chưa chọn (tài khoản nhiều cửa hàng
        chưa bấm chọn) $\rightarrow$ màn hiện đúng một dòng:
        \emph{``Không xác định được cửa hàng đang thao tác. Vui lòng chọn lại cửa hàng.''} và
        \textbf{không chạy truy vấn nào}.
  \item Chỉ thấy đơn của \textbf{đúng một cửa hàng} đang chọn --- khác với Trang chủ (cộng gộp)
        và khác với Giao hàng / Báo cáo.
  \item \textbf{Vai trò không giới hạn gì}: Nhân viên xem được toàn bộ lịch sử đơn và số tiền
        của cửa hàng, ngang Chủ tiệm.
  \item Đọc bằng quyền hệ thống; giới hạn duy nhất là điều kiện \rt{franchise_id} trong bảng
        dưới.
\end{itemize}

\mhien

\begin{hienbang}
Bảng đơn hàng &
  \rt{sale.order} &
  Đơn của \textbf{đúng cửa hàng đang chọn} và \textbf{không phải đơn huỷ}
  (\rt{state != 'cancel'}). Lấy \textbf{cả} đơn đặt qua portal lẫn đơn Ngô Gia tạo hộ trong
  ERP. &
  \emph{Ngày tạo} mới nhất trước; 10 đơn/trang &
  Cầu Giấy \uat{8} · Đống Đa \uat{5} · TP HCM Q1 \uat{18} · TEST \uat{0} \\ \addlinespace
Cột \emph{Trạng thái} &
  \rt{sale.order.state} + \rt{stock.picking.batch.delivery_batch_status} &
  Năm nhãn: \emph{Chờ xác nhận} (nháp), \emph{Đã gửi}, \emph{Đã xác nhận}, và hai nhãn đè từ
  chuyến giao: \emph{Đang giao}, \emph{Hoàn tất}. Đơn đã xác nhận mà chuyến đang chạy thì hiện
  \emph{Đang giao}; chuyến đã xong thì hiện \emph{Hoàn tất}. &
  --- &
  \uat{17} Chờ xác nhận · \uat{0} Đã gửi · \uat{14} Đã xác nhận \\ \addlinespace
Cột \emph{Số dòng sản phẩm} &
  \rt{sale.order.line} &
  Đếm dòng hàng thật, \textbf{bỏ} dòng tiêu đề/ghi chú &
  --- & --- \\ \addlinespace
Cột \emph{Người đặt} &
  \rt{sale.order.portal_requester_user_id} &
  Đơn đặt qua portal $\rightarrow$ tên người bấm gửi. Đơn Ngô Gia tạo trong ERP $\rightarrow$
  nhãn cố định \emph{``Ngô Gia tạo đơn''}; \textbf{cố ý không} lộ tên nhân sự nội bộ. &
  --- &
  \uat{23} đơn portal / \uat{24} đơn tạo từ ERP \\ \addlinespace
Cột \emph{Chuyến giao} &
  \rt{stock.picking.batch} &
  Mã chuyến và trạng thái chuyến (nhãn tiếng Việt cố định trong portal, không đổi theo ngôn
  ngữ người dùng) &
  --- &
  \uat{3}/47 đơn đã được xếp chuyến \\ \addlinespace
Cột \emph{Tổng tiền} &
  \rt{sale.order.amount_total} &
  Đọc thẳng số Odoo đã tính. Ký hiệu tiền lấy theo \textbf{loại tiền của chính đơn}. &
  --- &
  \uat{47}/47 đơn đang ghi bằng \textbf{USD} \\
\end{hienbang}

\mloc

\begin{itemize}[leftmargin=*]
  \item \textbf{Từ khoá} --- chỉ tìm theo \textbf{mã đơn}, khớp một phần, không phân biệt hoa
        thường. Không tìm theo tên sản phẩm, không tìm theo người đặt.
  \item \textbf{Khoảng ngày} --- lọc theo \emph{ngày tạo đơn}, quy đổi đúng theo múi giờ tài
        khoản (``từ ngày X'' tính từ 00:00 giờ địa phương). Khoảng ngày \textbf{ngược} thì
        \textbf{không chạy truy vấn}, giữ nguyên hai ô đã nhập và báo ngay tại thanh lọc
        \emph{``Từ ngày không được lớn hơn Đến ngày''} --- cố ý không để màn rơi vào trạng thái
        ``Chưa có đơn hàng'' làm người dùng tưởng mất dữ liệu.
  \item \textbf{Nút nhanh} --- \emph{Tháng này} / \emph{Tháng trước} tự điền hai ô ngày và
        \textbf{ghi đè} ngày người dùng vừa gõ.
  \item \textbf{Trạng thái} --- một nhãn tại một thời điểm, chọn trong đúng năm nhãn kể trên.
        Giá trị lạ bị bỏ qua (coi như không lọc).
  \item \textbf{Số đơn mỗi trang} --- chỉ nhận \textbf{10, 20, 50}; số khác thì lặng lẽ dùng 10.
\end{itemize}

\mkiem{một dòng đơn}

Không kiểm gì ở màn danh sách --- sang màn chi tiết, và màn chi tiết \textbf{kiểm lại quyền từ
đầu} (xem màn sau).

\mghi

\textbf{Không ghi gì.} Màn chỉ đọc.

\mhoi

\begin{ask}
\begin{enumerate}[leftmargin=*]
  \item \textbf{Ba phân hệ hiểu ``cửa hàng đang xem'' khác nhau.} Lịch sử mua hàng bắt buộc
        chọn cửa hàng; Giao hàng và Báo cáo thì cộng gộp mọi cửa hàng khi chưa chọn; Trang chủ
        cũng cộng gộp. Người quản lý 3 cửa hàng vì vậy thấy Trang chủ báo 31 đơn nhưng màn Lịch
        sử lại báo lỗi ``chưa chọn cửa hàng''. Anh chốt giúp một cách thống nhất cho cả portal.
  \item \textbf{Đơn Đã huỷ bị giấu hoàn toàn.} Trên UAT là \uat{16}/47 đơn. Cửa hàng không có
        cách nào tra lại đơn mình đã bị huỷ, kể cả khi biết mã đơn. Anh muốn (a) giữ nguyên,
        (b) thêm một nhãn lọc \emph{Đã huỷ} để tra khi cần?
  \item \textbf{Tìm kiếm chỉ theo mã đơn.} Cửa hàng thường nhớ ``đơn có món Hồng Trà'' chứ ít
        khi nhớ mã. Có cần tìm theo tên sản phẩm trong đơn không?
  \item \textbf{Nhân viên xem được toàn bộ số tiền các đơn.} Xác nhận đúng ý anh chứ, hay muốn
        giới hạn cột \emph{Tổng tiền} cho Chủ tiệm/Quản lý?
\end{enumerate}
\end{ask}

\section{Màn Chi tiết đơn hàng}
\rt{/portal/purchase-history/<int:order_id>}
\medskip

\mvao

Bấm một dòng ở màn Lịch sử mua hàng. Trên máy tính, đây cũng là nơi người dùng được đưa tới
ngay sau khi gửi đơn thành công (xem mục 3.4).

\mai

Phải đăng nhập, đã chọn cửa hàng, và đơn phải \textbf{thuộc đúng cửa hàng đó} và
\textbf{không phải đơn huỷ}. \textbf{Quyền được kiểm lại từ đầu ở mỗi lần mở}, không tin việc
người dùng vừa bấm từ danh sách --- chống trường hợp sửa số đơn trên thanh địa chỉ để xem đơn
cửa hàng khác. Sai điều kiện thì hiện đúng một câu: \emph{``Không tìm thấy đơn hàng hoặc bạn
không có quyền xem đơn hàng này.''} --- cùng một câu cho cả hai trường hợp, \textbf{cố ý không}
cho biết đơn đó có tồn tại hay không.

\mhien

\begin{hienbang}
Khối đầu đơn &
  \rt{sale.order} &
  Mã đơn, ngày tạo, ngày đặt, trạng thái (gộp như màn danh sách), tổng tiền, mã + tên cửa hàng,
  người đặt, số dòng sản phẩm &
  --- & --- \\ \addlinespace
Khối chuyến giao &
  \rt{stock.picking.batch} &
  Chỉ hiện khi đơn \textbf{đã được xếp chuyến}: mã chuyến, trạng thái chuyến, giờ khởi hành
  (\textbf{thực tế} nếu đã đi, chưa đi thì \textbf{dự kiến}), ghi chú giao hàng &
  --- &
  \uat{3}/47 đơn có chuyến $\Rightarrow$ đa số đơn không có khối này \\ \addlinespace
Bảng dòng sản phẩm &
  \rt{sale.order.line} &
  Dòng hàng thật của đơn, \textbf{bỏ} dòng tiêu đề/ghi chú. Mỗi dòng: tên sản phẩm, quy cách
  đóng gói, đơn vị tính, số lượng, chiết khấu, \textbf{đơn giá đã gồm thuế}, thành tiền. &
  Theo thứ tự dòng trong đơn &
  --- \\ \addlinespace
Ghi chú đơn &
  \rt{sale.order.portal_note} &
  Ghi chú người dùng đã gõ ở giỏ hàng lúc gửi đơn &
  --- & --- \\
\end{hienbang}

\textbf{Về con số đơn giá:} đơn giá hiển thị được tính lại qua bộ tính thuế của Odoo cho
\textbf{đúng một đơn vị}, \textbf{không} lấy thành tiền chia cho số lượng. Phép chia đó chỉ
đúng khi thuế là phần trăm thuần --- sai với thuế cố định, sai khi một dòng chịu nhiều thuế, và
sai khi làm tròn ở loại tiền có phần lẻ. Thành tiền thì đọc thẳng số của Odoo.

\mloc

Không có bộ lọc.

\mkiem{Tải PDF}

Xem màn tiếp theo.

\mghi

Không ghi gì.

\mhoi

\begin{ask}
Màn chi tiết \textbf{không hiện lịch sử thay đổi} của đơn --- cửa hàng không thấy được đơn bị
sửa số lượng hay bị đổi giá lúc nào, do ai. Anh có cần một khối ``Nhật ký đơn'' không?
\end{ask}

\section{Tải đơn hàng dạng PDF}
\rt{/portal/purchase-history/<int:order_id>.pdf}
\medskip

\mvao

Nút \emph{Tải PDF} ở màn chi tiết đơn.

\mai

Cùng điều kiện với màn chi tiết: đúng cửa hàng đang chọn, không phải đơn huỷ. Không đạt thì trả
về lỗi \emph{không tìm thấy} chứ không chuyển hướng.

\mhien

Không phải màn --- trả thẳng về tệp PDF, tên tệp là mã đơn.

\mloc

Không có tham số nào.

\mkiem{Tải PDF}

\begin{kiembang}
1 & Đã chọn được cửa hàng & Lỗi \emph{không tìm thấy} \\
2 & Đơn thuộc đúng cửa hàng đang chọn và chưa huỷ & Lỗi \emph{không tìm thấy} \\
3 & Mẫu in đơn hàng của Odoo còn tồn tại trong hệ thống & Lỗi \emph{không tìm thấy} \\
\end{kiembang}

\mghi

Không ghi gì vào dữ liệu nghiệp vụ.

\mhoi

\begin{ask}
\textbf{Tệp PDF đang dùng mẫu in Báo giá/Đơn bán mặc định của Odoo}, chưa phải mẫu có nhận diện
Wujia Tea (logo, thông tin công ty, điều khoản). Anh có muốn đặt mẫu riêng cho cửa hàng nhượng
quyền không --- và nếu có thì PDF nên hiện những mục gì?
\end{ask}

\section{Màn Theo dõi giao hàng}
\rt{/portal/delivery} · \rt{/portal/delivery/results}
\medskip

\mvao

Mục \emph{Giao hàng} trên thanh điều hướng. Đường dẫn \rt{/results} là phần kết quả gọi ngầm khi
bấm chip lọc hoặc đổi bộ lọc.

\mai

\begin{itemize}[leftmargin=*]
  \item Phải đăng nhập. \textbf{Chưa chọn cửa hàng thì vẫn xem được} --- lấy chuyến của
        \textbf{mọi} cửa hàng tài khoản thuộc (khác hẳn màn Lịch sử mua hàng).
  \item Tài khoản không thuộc cửa hàng nào $\rightarrow$ màn trống, không báo lỗi.
  \item Một chuyến xe có thể chở hàng cho \textbf{nhiều cửa hàng}. Màn chỉ hiện \textbf{phần
        phiếu giao thuộc cửa hàng của người xem}; phiếu của cửa hàng khác trong cùng chuyến
        \textbf{không} hiện.
  \item Vai trò không giới hạn gì.
\end{itemize}

\mhien

\begin{hienbang}
Danh sách chuyến &
  \rt{stock.picking.batch} &
  Chuyến có phiếu giao thuộc cửa hàng đang xét --- nhận theo \textbf{hai đường}: phiếu gắn
  thẳng cửa hàng, \textbf{hoặc} phiếu gắn đơn hàng của cửa hàng. &
  \emph{Giờ khởi hành dự kiến} muộn nhất trước; 20 chuyến mỗi trang &
  \uat{3} chuyến toàn hệ, \uat{2} chuyến có cửa hàng nhìn thấy (một chuyến chở chung TP HCM
  Q1 + Đống Đa); Cầu Giấy \uat{0} \\ \addlinespace
Chip lọc trạng thái &
  \rt{stock.picking.batch.delivery_batch_status} &
  Bốn chip: \emph{Tất cả} / \emph{Sắp giao} (nháp, đã gán xe, đang chất hàng) / \emph{Đang giao}
  / \emph{Đã giao}. Số trên chip đếm trên \textbf{chính tập đã lọc} theo ngày và từ khoá, nên
  chọn một chip không làm các chip khác về 0. &
  --- &
  \uat{1} nháp · \uat{1} đã giao · \uat{1} huỷ chuyến \\ \addlinespace
Giờ khởi hành trên thẻ &
  \rt{stock.picking.batch} &
  Đã xuất phát thì hiện \textbf{giờ thực tế} kèm nhãn \emph{Xuất phát (thực tế)}; chưa đi thì
  hiện \textbf{giờ dự kiến} kèm nhãn \emph{Xuất phát (dự kiến)}. &
  --- & --- \\ \addlinespace
Xe và tài xế &
  \rt{stock.picking.batch.vehicle_id} &
  Tên xe, tên tài xế, biển số. Thiếu thì hiện dấu ``---''. &
  --- & --- \\ \addlinespace
Đơn liên quan trên thẻ &
  \rt{sale.order} &
  Mã các đơn thuộc cửa hàng người xem trong chuyến đó, rút gọn dạng \emph{S00035, S00036 +3} &
  --- & --- \\
\end{hienbang}

\textbf{Lưu ý về chip ``Huỷ chuyến'':} trạng thái \emph{Hủy chuyến} \textbf{không có chip
riêng}. Chuyến bị huỷ vẫn nằm trong danh sách khi chọn \emph{Tất cả}, nhưng không lọc riêng
được. Trên UAT đang có \uat{1} chuyến như vậy.

\mloc

\begin{itemize}[leftmargin=*]
  \item \textbf{Từ khoá} --- tìm theo \textbf{mã chuyến} \textbf{hoặc} \textbf{mã đơn hàng}
        trong chuyến.
  \item \textbf{Khoảng ngày} --- lọc theo \emph{giờ khởi hành dự kiến}, quy đổi theo múi giờ
        tài khoản. Ngày ngược $\rightarrow$ không chạy truy vấn, báo tại thanh lọc.
        \textbf{Lưu ý}: chuyến \textbf{chưa đặt lịch khởi hành} sẽ bị loại khỏi kết quả khi
        dùng bộ lọc ngày.
  \item \textbf{Chip trạng thái} --- một chip tại một thời điểm.
  \item \textbf{Số chuyến mỗi trang} --- chỉ nhận \textbf{10, 20, 50}; mặc định 20.
  \item Có tham số kỹ thuật \rt{_preview} dùng để đội phát triển xem thử bốn trạng thái màn
        (đang tải / lỗi / rỗng / có dữ liệu). \textbf{Không} phải chức năng nghiệp vụ, BA không
        cần lên task.
\end{itemize}

\mkiem{một chuyến}

Không kiểm gì ở danh sách --- sang màn chi tiết, và màn đó kiểm lại quyền.

\mghi

Không ghi gì.

\mhoi

\begin{ask}
\begin{enumerate}[leftmargin=*]
  \item \textbf{Cửa hàng đông đơn nhất lại không có chuyến nào.} Trên UAT, Cầu Giấy có \uat{8}
        đơn nhưng \uat{0} chuyến giao; cả hệ thống mới \uat{3}/47 đơn được xếp chuyến. Nghĩa là
        màn Giao hàng gần như trống với phần lớn cửa hàng. Đây là do dữ liệu UAT chưa đủ, hay
        nghiệp vụ thật cũng chỉ một phần đơn đi qua chuyến?
  \item \textbf{Chuyến chưa đặt lịch bị mất khi lọc theo ngày.} Người dùng lọc ``tháng này'' sẽ
        không thấy chuyến chưa có giờ khởi hành, dù chúng đang chờ giao. Anh muốn giữ nguyên
        hay luôn hiện chúng?
  \item \textbf{Không lọc riêng được chuyến bị huỷ.} Có cần thêm chip \emph{Huỷ chuyến} không?
\end{enumerate}
\end{ask}

\section{Màn Chi tiết chuyến giao}
\rt{/portal/delivery/<int:batch_id>}
\medskip

\mvao

Bấm một thẻ chuyến ở màn Theo dõi giao hàng.

\mai

\begin{itemize}[leftmargin=*]
  \item Phải đăng nhập và thuộc ít nhất một cửa hàng.
  \item Chuyến phải có \textbf{phiếu giao gắn thẳng cửa hàng} của người xem. Không đạt thì
        \textbf{đưa về màn danh sách}, không báo lỗi.
  \item \textbf{Điều kiện này HẸP HƠN màn danh sách}: danh sách nhận chuyến theo \emph{hai}
        đường (phiếu gắn cửa hàng \textbf{hoặc} phiếu gắn đơn của cửa hàng), còn chi tiết chỉ
        nhận \emph{một} đường. Về lý thuyết có thể xảy ra trường hợp chuyến hiện ở danh sách
        nhưng bấm vào lại bị đẩy về. Trên UAT hiện \textbf{chưa} xảy ra (cả hai cách đều ra
        \uat{2} chuyến), nhưng đang có \uat{1}/18 phiếu giao chưa gắn cửa hàng --- đúng loại dữ
        liệu sẽ kích hoạt lệch này. Xem mục 7.
\end{itemize}

\mhien

\begin{hienbang}
Khối đầu chuyến &
  \rt{stock.picking.batch} &
  Mã chuyến, trạng thái, giờ khởi hành (thực tế hoặc dự kiến), giờ cập nhật gần nhất, tên cửa
  hàng, kho xuất hàng &
  --- & --- \\ \addlinespace
Thanh tiến trình 3 bước &
  \rt{stock.picking.batch.delivery_batch_status} &
  \emph{Đã xuất phát} $\rightarrow$ \emph{Đang giao} $\rightarrow$ \emph{Đã giao xong}. Bước
  được tô sáng suy ra từ trạng thái chuyến và giờ khởi hành thực tế. &
  --- & --- \\ \addlinespace
Danh sách sản phẩm trong chuyến &
  \rt{stock.move} &
  Gộp theo sản phẩm, cộng số lượng --- \textbf{chỉ} phần thuộc phiếu giao của cửa hàng người
  xem &
  Theo thứ tự gặp &
  --- \\ \addlinespace
Mã đơn liên quan &
  \rt{sale.order} &
  Các đơn của cửa hàng người xem nằm trong chuyến &
  --- & --- \\
\end{hienbang}

\textbf{Về giờ ở thanh tiến trình:} hai mốc \emph{Đang giao} và \emph{Đã giao xong} hiện
\textbf{giờ cập nhật bản ghi gần nhất}, không phải giờ nghiệp vụ thật của từng bước --- hệ thống
hiện chưa lưu riêng ``giờ bắt đầu giao'' và ``giờ giao xong''. Xem mục 7.

\mloc

Không có bộ lọc.

\mkiem{Thêm vào lịch}

Xem màn tiếp theo.

\mghi

Không ghi gì.

\mhoi

\begin{ask}
\begin{enumerate}[leftmargin=*]
  \item \textbf{Điều kiện xem chi tiết hẹp hơn điều kiện hiện ở danh sách.} Nên sửa cho giống
        nhau --- em ghi nhận là việc kỹ thuật, không cần BA quyết, chỉ báo để anh biết có một
        đường dẫn sẽ được vá ở cụm F.
  \item \textbf{Chưa có giờ thật cho hai bước cuối.} Thanh tiến trình đang mượn ``giờ sửa bản
        ghi'' để hiển thị. Nếu anh cần cửa hàng thấy đúng \emph{mấy giờ xe bắt đầu giao} và
        \emph{mấy giờ giao xong} thì phải bổ sung hai trường ở phân hệ giao vận --- đây là
        thay đổi có khối lượng, cần anh xác nhận trước.
\end{enumerate}
\end{ask}

\section{Tải lịch giao hàng vào ứng dụng lịch}
\rt{/portal/delivery/<int:batch_id>.ics}
\medskip

\mvao

Nút \emph{Thêm vào lịch} ở màn chi tiết chuyến. Tệp tải về mở được bằng Google Calendar,
Outlook, Lịch của iPhone.

\mai

Cùng điều kiện với màn chi tiết chuyến. Không đạt thì trả lỗi \emph{không tìm thấy}.

\mhien

Không phải màn. Tệp lịch chứa: tiêu đề \emph{``Giao hàng + mã chuyến''}, mô tả liệt kê các
phiếu giao, địa điểm là tên các cửa hàng nhận hàng.

\mloc

Không có tham số nào.

\mkiem{Thêm vào lịch}

\begin{kiembang}
1 & Tài khoản thuộc ít nhất một cửa hàng & Lỗi \emph{không tìm thấy} \\
2 & Chuyến có phiếu giao gắn cửa hàng của người xem & Lỗi \emph{không tìm thấy} \\
\end{kiembang}

\mghi

Không ghi gì.

\mhoi

\begin{ask}
\textbf{Giờ trong tệp lịch lấy từ một trường khác với giờ hiển thị trên màn.} Màn hiện
\emph{giờ khởi hành} (thực tế hoặc dự kiến), còn tệp lịch lấy \emph{ngày giờ dự kiến của phiếu
giao} và mặc định kéo dài \textbf{2 tiếng}. Hai con số này có thể lệch nhau. Anh chốt: sự kiện
trong lịch nên bắt đầu lúc nào và kéo dài bao lâu?
\end{ask}

\section{Màn Báo cáo đặt hàng}
\rt{/portal/reports/orders}
\medskip

\mvao

Mục \emph{Báo cáo} trên thanh điều hướng. \textbf{Nhân viên không nhìn thấy mục này} và vào
thẳng đường dẫn cũng bị đẩy về Trang chủ.

\mai

\begin{itemize}[leftmargin=*]
  \item Phải đăng nhập và thuộc ít nhất một cửa hàng.
  \item \textbf{Phải là Chủ tiệm hoặc Quản lý.} Nhân viên $\rightarrow$ đẩy về Trang chủ, không
        thông báo. Trên UAT có \uat{4}/10 bản ghi thành viên là Nhân viên.
  \item \textbf{Cách xét vai trò ở đây rộng hơn phạm vi dữ liệu}: hệ thống lấy vai trò
        \textbf{cao nhất trên MỌI cửa hàng} tài khoản thuộc, rồi cho xem số liệu của
        \textbf{cửa hàng đang chọn}. Nghĩa là người chỉ là Nhân viên ở cửa hàng A nhưng là Quản
        lý ở cửa hàng B vẫn xem được báo cáo của cửa hàng A. Trên UAT có đúng \uat{1} tài khoản
        rơi vào tình huống này (Quản lý ở Đống Đa, Nhân viên ở Cầu Giấy và TP HCM Q1). Xem mục 7.
  \item Chưa chọn cửa hàng $\rightarrow$ \textbf{cộng gộp mọi cửa hàng} tài khoản thuộc.
\end{itemize}

\mhien

\begin{hienbang}
4 thẻ số liệu &
  \rt{sale.order} &
  Trong khoảng ngày đang lọc, của cửa hàng đang xét: \emph{Tổng số đơn} và \emph{Tổng doanh
  thu} (\textbf{loại} đơn huỷ), \emph{Đơn đã huỷ}, \emph{Đơn hoàn thành}. &
  --- &
  \uat{31} đơn không huỷ · \uat{16} đơn huỷ · \uat{0} đơn ``hoàn thành'' \\ \addlinespace
Biểu đồ cột theo tháng &
  \rt{sale.order} &
  Số đơn và doanh thu theo từng tháng, \textbf{loại} đơn huỷ. Khoảng luôn được nới ra
  \textbf{ít nhất 12 tháng gần nhất} để biểu đồ không trống, kể cả khi người dùng lọc một
  khoảng ngắn. &
  Theo tháng tăng dần &
  --- \\ \addlinespace
Bảng Top 10 sản phẩm &
  \rt{sale.order.line} &
  Dòng hàng thuộc đơn \textbf{đã xác nhận} của cửa hàng đang xét, trong khoảng ngày. Gộp theo
  sản phẩm, cộng số lượng và thành tiền. &
  Số lượng nhiều nhất trước; 10 dòng &
  --- \\ \addlinespace
Biểu đồ tròn + bảng phân bố trạng thái &
  \rt{sale.order} &
  Đếm đơn theo từng trạng thái kèm tỷ lệ \%. \textbf{Có tính} cả đơn huỷ. &
  --- &
  \uat{17} Nháp · \uat{0} Đã gửi · \uat{14} Đã xác nhận · \uat{0} Hoàn thành · \uat{16} Đã huỷ \\
\end{hienbang}

\mloc

Hai ô \emph{Từ ngày} / \emph{Đến ngày}, lọc theo \emph{ngày đặt hàng}. Bỏ trống thì mặc định là
\textbf{từ 01/01 năm nay đến hôm nay}. Ngày ngược $\rightarrow$ giữ nguyên chữ đã gõ, báo tại
thanh lọc, mọi số liệu và biểu đồ về rỗng.

\mkiem{Tải Excel}

Xem màn tiếp theo.

\mghi

Không ghi gì.

\mhoi

\begin{ask}
\begin{enumerate}[leftmargin=*]
  \item \textbf{Thẻ ``Đơn hoàn thành'' luôn bằng 0.} Odoo 19 chỉ có bốn trạng thái đơn
        (\emph{Nháp}, \emph{Đã gửi}, \emph{Đã xác nhận}, \emph{Đã huỷ}) --- \textbf{không có}
        trạng thái \emph{Hoàn thành}. Thẻ này và lát ``Hoàn thành'' trên biểu đồ tròn vì vậy
        không bao giờ có số. Anh muốn (a) bỏ thẻ, hay (b) định nghĩa lại ``hoàn thành'' = đơn
        đã xác nhận và chuyến giao đã xong (giống nhãn \emph{Hoàn tất} ở màn Lịch sử)?
  \item \textbf{Số tiền trên báo cáo đang dán sai ký hiệu tiền tệ.} Màn Báo cáo định dạng mọi
        số tiền theo \textbf{loại tiền của công ty}, trong khi màn Lịch sử mua hàng dùng
        \textbf{loại tiền của từng đơn}. Trên UAT, công ty đang để \textbf{VND} còn cả
        \uat{47}/47 đơn ghi bằng \textbf{USD} --- nghĩa là báo cáo đang hiện \emph{con số USD
        kèm ký hiệu ₫}. Đây đúng loại lỗi WJ-ORD-025 đã xử lý ở phân hệ Đặt hàng. Em đề xuất
        sửa theo hướng dùng loại tiền của đơn; anh xác nhận giúp.
  \item \textbf{Bộ lọc ngày của Báo cáo không quy đổi múi giờ}, khác với màn Lịch sử mua hàng
        và Giao hàng (hai màn này đã quy đổi). Hệ quả: chọn cùng một khoảng ngày, hai màn có thể
        ra số đơn khác nhau ở hai ngày biên. Đây là lỗi kỹ thuật cần vá --- em báo để anh biết,
        không cần BA quyết.
  \item \textbf{Vai trò được xét trên mọi cửa hàng, nhưng dữ liệu chỉ của cửa hàng đang chọn.}
        Người là Quản lý ở một cửa hàng đọc được báo cáo doanh thu của cửa hàng mà họ chỉ là
        Nhân viên. Anh muốn siết theo đúng cửa hàng đang xem không?
  \item \textbf{Biểu đồ luôn nới ra 12 tháng} kể cả khi người dùng lọc một tuần --- nên con số
        trên biểu đồ không khớp bốn thẻ số liệu phía trên. Anh muốn biểu đồ bám đúng khoảng
        lọc, hay giữ nguyên và ghi chú ``12 tháng gần nhất''?
\end{enumerate}
\end{ask}

\section{Tải báo cáo dạng Excel}
\rt{/portal/reports/orders/export.xlsx}
\medskip

\mvao

Nút \emph{Xuất Excel} ở màn Báo cáo.

\mai

Phải là Chủ tiệm hoặc Quản lý --- nhưng \textbf{xét theo cửa hàng đang chọn}, không phải theo
mọi cửa hàng như màn Báo cáo. Nghĩa là có trường hợp \textbf{xem được màn nhưng bấm xuất Excel
thì bị đẩy về Trang chủ} mà không có lời giải thích nào. Xem mục 7.

\mhien

Không phải màn --- trả về tệp Excel một trang tính tên \emph{Orders}, tên tệp kèm khoảng ngày.
Bảy cột: \emph{Mã đơn}, \emph{Ngày đặt}, \emph{Cửa hàng}, \emph{Khách hàng}, \emph{Trạng thái},
\emph{Số dòng}, \emph{Tổng tiền}.

\mloc

Nhận cùng hai ô ngày với màn Báo cáo. Khác ở chỗ: khoảng ngày ngược \textbf{không báo lỗi} mà
bị lặng lẽ kéo về đúng một ngày.

\mkiem{Xuất Excel}

\begin{kiembang}
1 & Tài khoản thuộc ít nhất một cửa hàng & Đưa về Trang chủ \\
2 & Vai trò tại \textbf{cửa hàng đang chọn} là Chủ tiệm hoặc Quản lý &
    Đưa về Trang chủ, không thông báo \\
\end{kiembang}

\mghi

Không ghi gì.

\mhoi

\begin{ask}
\begin{enumerate}[leftmargin=*]
  \item \textbf{Tệp Excel chứa cả đơn đã huỷ, màn hình thì không.} Người dùng đối chiếu hai bên
        sẽ thấy lệch \uat{16} đơn trên dữ liệu UAT. Anh chốt: tệp Excel nên có hay không có đơn
        huỷ?
  \item \textbf{Cột \emph{Số dòng} trong Excel đếm cả dòng tiêu đề/ghi chú}, còn màn hình chỉ
        đếm dòng sản phẩm thật. Hai nơi sẽ ra số khác nhau với đơn có dòng ghi chú.
  \item \textbf{Cột \emph{Tổng tiền} trong Excel là số trần, không kèm loại tiền} --- với dữ
        liệu USD/VND lẫn lộn như hiện nay thì người đọc không biết đơn vị. Nên thêm một cột
        \emph{Loại tiền}, anh đồng ý chứ?
  \item \textbf{Điều kiện vai trò của nút Xuất khác của màn} (đã nêu ở mục 2). Cần thống nhất
        theo cách anh chốt ở câu 4 mục 5.8.
\end{enumerate}
\end{ask}

\section{Đường dẫn cũ của Lịch sử mua hàng}
\rt{/portal/purchase_history}
\medskip

Đường dẫn thời bản v14 (viết gạch dưới), được giữ lại làm chuyển hướng vĩnh viễn sang
\rt{/portal/purchase-history}. Không hiện màn, không đọc dữ liệu. BA không cần lên task.

\chapter{Công nợ --- Đổi trả · Bù hàng}

Chương này gồm \textbf{8 đường dẫn} gom thành \textbf{7 màn} cộng một đường dẫn cũ, thuộc hai
mô-đun: \rt{wujia_portal_debt} (3) và \rt{wujia_portal_return} (5).

\textbf{Bốn điểm cần biết trước khi đọc:}

\begin{enumerate}[leftmargin=*]
  \item \textbf{Hai phân hệ trong cùng một chương hiểu ``cửa hàng'' ngược nhau.} Công nợ chỉ
        làm việc với \textbf{đúng một cửa hàng đang chọn}; chưa chọn thì màn rỗng. Đổi trả thì
        \textbf{gộp mọi cửa hàng} tài khoản truy cập được --- một tài khoản ba cửa hàng mở màn
        Đổi trả sẽ thấy lẫn phiếu của cả ba. Cùng với ba kiểu đã nêu ở chương trước, portal
        hiện có \textbf{bốn cách} xác định ``cửa hàng đang xem''.
  \item \textbf{Công nợ là phân hệ chặn theo vai trò chặt nhất.} Nhân viên \textbf{không vào
        được cả ba màn} công nợ --- không phải ẩn số liệu, mà là một trang ``không có quyền''
        riêng. Và vai trò được xét tại \textbf{đúng cửa hàng đang chọn}, khác với Báo cáo
        (chương 5) xét vai trò cao nhất trên mọi cửa hàng.
  \item \textbf{Trên UAT chỉ 1 trong 4 cửa hàng có chứng từ kế toán.} Toàn bộ \uat{6} chứng từ
        đã ghi sổ đều thuộc \emph{TP HCM Quận 1}. Ba cửa hàng còn lại mở màn Công nợ lúc nào
        cũng thấy trạng thái rỗng --- đó là do dữ liệu, không phải do luật lọc sai.
  \item \textbf{Phiếu đổi trả lưu nháp là ngõ cụt.} Màn tạo yêu cầu trên máy tính có nút
        \emph{Lưu nháp}, nhưng portal \textbf{không có đường nào} để gửi một phiếu nháp đã lưu:
        màn chi tiết chỉ đọc, không có nút gửi. UAT đang có \uat{3} phiếu nháp nằm vĩnh viễn ở
        đó. Xem mục 7 của màn Tạo yêu cầu.
\end{enumerate}

Quy ước chung toàn portal (phân trang, định dạng ngày giờ và tiền, huy hiệu trạng thái, vai
trò, đính kèm) đã viết một lần ở mục~\ref{sec:quy-uoc-chung}; chương này chỉ nói phần khác
biệt.

% =============================================================================
\section{Màn Công nợ theo tuần}
\rt{/portal/debt}
\medskip

\mvao

Mục \emph{Công nợ} trên thanh điều hướng bên trái (máy tính) hoặc dòng \emph{Công nợ \& thanh
toán} trong bảng \emph{Thêm} (điện thoại). Màn nhận các tham số trên thanh địa chỉ:
\rt{week} (tuần đang xem), \rt{all} (điện thoại: bung hết hoá đơn), \rt{page} ·
\rt{page_size} (phân trang bảng máy tính), \rt{notice} (băng thông báo khi vừa bị đẩy về từ
màn chuyển khoản).

\mai

\begin{itemize}[leftmargin=*]
  \item Phải đăng nhập.
  \item \textbf{Chỉ Chủ tiệm và Quản lý}. Nhân viên mở màn này sẽ nhận một \textbf{trang thông
        báo không có quyền} --- không phải lỗi hệ thống, và không lộ bất kỳ con số nào.
  \item Vai trò được xét tại \textbf{đúng cửa hàng đang chọn}. Một người vừa là Quản lý cửa
        hàng A vừa là Nhân viên cửa hàng B: đang chọn A thì vào được, bấm đổi sang B thì chính
        tài khoản đó bị chặn. Trên UAT có đúng \uat{1} tài khoản như vậy (Phạm Thị Dung ---
        Quản lý \emph{Đống Đa}, Nhân viên \emph{TP HCM Quận 1} và \emph{Cầu Giấy}).
  \item \textbf{Chưa chọn cửa hàng thì không bị chặn.} Không có cửa hàng thì không có vai trò
        để xét, nên hệ thống cho vào và hiện trạng thái rỗng ``chưa chọn cửa hàng''. Nghĩa là
        một Nhân viên nhiều cửa hàng chưa bấm chọn vẫn mở được màn --- chỉ là không thấy số.
  \item Số liệu đọc bằng quyền hệ thống; giới hạn duy nhất là \rt{franchise_id} bằng đúng cửa
        hàng đang chọn.
  \item Thực tế UAT: \uat{4} cửa hàng, \uat{10} lượt tham gia, trong đó \uat{6} lượt đạt Chủ
        tiệm/Quản lý. Nhưng chỉ \emph{TP HCM Quận 1} có chứng từ, nên rốt cuộc chỉ \uat{2} tài
        khoản thực sự nhìn thấy số: chủ cửa hàng đó và Administrator (Quản lý cùng cửa hàng).
\end{itemize}

\mhien

\begin{hienbang}
Ô chọn tuần &
  --- (máy tự sinh) &
  \textbf{6 tuần}: tuần hiện tại và 5 tuần trước, tính từ thứ Hai. Nhãn dạng
  \emph{``10/08 -- 16/08/2026''}. Tuần cũ hơn \textbf{không chọn tới được}. &
  Mới nhất trước &
  Tại ngày chụp: từ 14/09 lùi về 10/08/2026 \\ \addlinespace
Tuần mở sẵn &
  \rt{account.move} &
  Thứ tự ưu tiên: (1) tuần \textbf{quá hạn cũ nhất còn dư nợ}; (2) nếu không có, tuần liền
  trước nếu tuần đó có chứng từ; (3) nếu vẫn không, tuần hiện tại. &
  --- &
  TP HCM Q1 mở sẵn tuần \uat{10/08 -- 16/08} \\ \addlinespace
Thẻ tổng của tuần &
  \rt{account.move} &
  Hoá đơn bán và giấy báo có \textbf{đã ghi sổ} (\rt{state = posted},
  \rt{move_type in (out_invoice, out_refund)}) của cửa hàng, có \textbf{ngày hoá đơn} rơi vào
  thứ Hai--Chủ nhật của tuần. Giấy báo có \textbf{trừ} vào tổng. Chứng từ nháp không tính. &
  Ngày hoá đơn, rồi số thứ tự &
  \uat{6} chứng từ ghi sổ / \uat{6} chứng từ còn nháp \\ \addlinespace
Số \emph{Còn phải trả} &
  \rt{account.move.amount_residual} &
  Tổng số dư ròng, \textbf{cắt sàn ở 0}: âm thì đưa sang ô \emph{Dư có}. Không bao giờ in số
  nợ âm. &
  --- &
  Tuần 10/08: còn phải trả \uat{0} · dư có \uat{\$72.448,85} \\ \addlinespace
Trạng thái tuần &
  suy từ số dư &
  Năm nhãn, xét theo thứ tự: hết dư $\rightarrow$ \emph{Đã thanh toán}; dư âm $\rightarrow$
  \emph{Dư có}; có hoá đơn quá hạn $\rightarrow$ \emph{Có quá hạn}; trả được một phần
  $\rightarrow$ \emph{Thanh toán một phần}; còn lại \emph{Chưa thanh toán}. &
  --- &
  Tuần 10/08 hiện \emph{Dư có} \\ \addlinespace
Hạn thanh toán &
  máy tự tính &
  Thứ Hai của tuần \textbf{cộng 10 ngày} --- tức thứ Năm tuần kế. Không đọc hạn ghi trên hoá
  đơn. &
  --- &
  Tuần 10/08 $\Rightarrow$ hạn \uat{20/08/2026} \\ \addlinespace
Bảng hoá đơn (máy tính) &
  \rt{account.move} &
  Mỗi dòng: số chứng từ, ngày, hạn, trạng thái, \textbf{tổng}, \textbf{đã trả}, \textbf{còn
  phải trả}. Giấy báo có in số âm. &
  \textbf{10} dòng mỗi trang &
  Tuần 10/08 có \uat{5} chứng từ \\ \addlinespace
Khối hoá đơn (điện thoại) &
  \rt{account.move} &
  Chỉ \textbf{2 hoá đơn} đầu kèm dòng ``còn \emph{n} hoá đơn''; bung hết bằng \rt{?all=1}. &
  --- & --- \\ \addlinespace
Trạng thái từng hoá đơn &
  suy từ số dư &
  \emph{Đã thanh toán} (hết dư, hoặc kế toán đã đánh dấu tất toán) · \emph{Giấy báo có} ·
  \emph{Quá hạn} · \emph{Một phần} (đã trả được ít nhất một đồng) · \emph{Chưa thanh toán}. &
  --- &
  Tuần 10/08: \uat{3} Đã thanh toán · \uat{1} Giấy báo có · \uat{1} Quá hạn \\ \addlinespace
Hộp chuyển khoản &
  \rt{res.partner.bank} &
  Chỉ dựng khi \textbf{còn phải trả lớn hơn 0}. Lấy tài khoản của công ty có bật cờ nhận tiền
  qua portal, số thứ tự nhỏ nhất. &
  --- &
  \uat{1}/1 tài khoản đã bật cờ \\
\end{hienbang}

\textbf{Về loại tiền:} ký hiệu tiền lấy theo \textbf{loại tiền của chính các chứng từ trong
tuần}; tuần trộn nhiều loại tiền thì lùi về tiền công ty. Trên UAT toàn bộ chứng từ ghi bằng
USD nên màn hiển thị \$ --- đúng, khác với màn Báo cáo (chương 5) đang in ký hiệu đồng Việt Nam cho số
USD.

\mloc

\begin{itemize}[leftmargin=*]
  \item \textbf{Tuần} --- chọn trong đúng 6 khoá tuần kể trên. Gõ khoá lạ trên thanh địa chỉ
        thì \textbf{không báo lỗi}, hệ thống lặng lẽ quay về tuần mặc định.
  \item \textbf{Bung hết hoá đơn} (điện thoại) --- \rt{?all=1}; giá trị khác coi như không bung.
  \item \textbf{Số dòng mỗi trang} (máy tính) --- chỉ nhận \textbf{10, 20, 50}.
  \item \textbf{Băng thông báo} --- chỉ chấp nhận đúng một giá trị \rt{no_due}; mọi chuỗi khác
        bị bỏ qua, nên không ai nhét được câu chữ tuỳ ý vào trang qua thanh địa chỉ.
\end{itemize}

\mkiem{Thanh toán}

\begin{kiembang}
1 & Đã chọn được cửa hàng & Màn rỗng ``chưa chọn cửa hàng'', nút không hiện \\
2 & Vai trò tại cửa hàng đó là Chủ tiệm hoặc Quản lý & Trang ``không có quyền'' \\
3 & Tuần đang xem \textbf{còn phải trả lớn hơn 0} &
    Bị đẩy về chính màn này kèm băng báo không còn khoản phải trả --- \textbf{không} in số âm,
    \textbf{không} dựng mã QR \\
4 & Công ty đã khai tài khoản nhận tiền bật cho portal &
    Màn chuyển khoản hiện thông báo chưa cấu hình, không hiện ô số tài khoản trống \\
\end{kiembang}

\mghi

\textbf{Không ghi gì.} Xem số tài khoản hay mở mã QR \textbf{không tạo phiếu thanh toán}, không
đánh dấu hoá đơn. Nội dung chuyển khoản do máy chủ tự sinh từ mã cửa hàng, số tuần và số tiền
--- \textbf{không tin} số tiền trình duyệt gửi lên.

\mhoi

\begin{ask}
\begin{enumerate}[leftmargin=*]
  \item \textbf{Số công nợ ở Trang chủ và số ở màn này là hai phép tính khác nhau.} Thẻ Trang
        chủ và huy hiệu đọc \textbf{ô tổng lưu sẵn} trên cửa hàng: tính bằng \textbf{tiền công
        ty}, \textbf{không giới hạn tuần}, cập nhật bằng tác vụ nền hằng ngày. Màn này tính
        trực tiếp bằng \textbf{tiền của chứng từ} và \textbf{chỉ trong một tuần}. UAT chưa khai
        tỷ giá nên hai số đang trùng nhau; khi khai tỷ giá thật, hai chỗ sẽ ra hai con số khác
        nhau cho cùng một cửa hàng.
  \item \textbf{Dropdown chỉ 6 tuần.} UAT có \uat{1} hoá đơn ngày 05/08/2026 nằm ngoài khung
        --- cửa hàng \textbf{không có cách nào} xem lại. Anh muốn giữ 6 tuần, hay thêm ô chọn
        khoảng ngày tự do?
  \item \textbf{``Dư có'' đè lên ``Có quá hạn''.} Tuần 10/08 của TP HCM Q1 có một hoá đơn quá
        hạn \$1,15 nhưng thẻ tuần hiện \emph{Dư có} vì giấy báo có lớn hơn, và nút Thanh toán
        biến mất. Cửa hàng nhìn vào tưởng không còn gì phải xử lý. Anh chốt giúp: quá hạn có
        được ưu tiên hiện không?
  \item \textbf{Trạng thái rỗng nói ``không còn công nợ''.} Câu này dựa trên ô tổng lưu sẵn có
        lớn hơn 0 hay không. TP HCM Q1 đang có ô tổng \textbf{âm}, nên mọi tuần không có chứng
        từ đều báo ``không còn công nợ'' --- trong khi vẫn còn một hoá đơn quá hạn.
  \item \textbf{Hạn thanh toán do portal tự tính} (thứ Hai + 10 ngày), \textbf{không} đọc hạn
        ghi trên hoá đơn. UAT có hoá đơn hạn 12/08 trong khi portal nói hạn tuần là 20/08.
        Lấy theo cái nào?
  \item \textbf{Tuần mở sẵn} = tuần quá hạn cũ nhất còn dư. Anh xác nhận cách mở này.
  \item \textbf{Ảnh mã QR vẫn đang treo} --- hiện chỉ có số tài khoản và nội dung chuyển khoản.
        Anh chốt giúp: mã QR tĩnh dán sẵn hay mã động sinh theo số tiền từng tuần?
  \item \textbf{Nhân viên bị chặn khỏi toàn bộ phân hệ Công nợ.} Xác nhận đúng ý anh.
\end{enumerate}
\end{ask}

% =============================================================================
\section{Màn Lịch sử thanh toán}
\rt{/portal/debt/payment-history}
\medskip

\mvao

Thẻ \emph{Lịch sử thanh toán} ngay trong phân hệ Công nợ. Tham số: \rt{month} (kỳ),
\rt{q} (tìm kiếm), \rt{page} · \rt{page_size}. Ngoài ra còn nhận \rt{date_from} ·
\rt{date_to} nhưng giao diện \textbf{chưa có ô nào} sinh ra hai tham số này.

\mai

Giống hệt màn Công nợ: chỉ Chủ tiệm / Quản lý của \textbf{đúng cửa hàng đang chọn}; Nhân viên
nhận trang không có quyền; chưa chọn cửa hàng thì vào được nhưng rỗng.

\mhien

\begin{hienbang}
Ô chọn kỳ &
  --- (máy tự sinh) &
  \textbf{6 kỳ}, mỗi kỳ là \textbf{trọn một tháng}: tháng hiện tại và 5 tháng trước. Nhãn
  \emph{``01/09 -- 30/09/2026''}. &
  Mới nhất trước &
  Tại ngày chụp: 09/2026 lùi về 04/2026 \\ \addlinespace
Bảng giao dịch &
  \rt{account.payment} &
  Phiếu thanh toán của cửa hàng \textbf{đã được xác nhận}
  (\rt{state in (in_process, paid)}), có \textbf{ngày thanh toán} nằm trong kỳ. Phiếu nháp,
  huỷ, từ chối \textbf{không hiện}. &
  Ngày mới nhất trước; \textbf{10} dòng mỗi trang &
  \uat{5}/6 phiếu đã xác nhận \\ \addlinespace
Cột \emph{Số tiền} &
  \rt{account.payment.amount} &
  Tiền cửa hàng trả vào hiện số dương; \textbf{tiền Ngô Gia chi ngược trả lại} hiện
  \textbf{số âm}, để không bị cộng nhầm vào phần đã trả. Ký hiệu tiền lấy theo
  \textbf{từng giao dịch}. &
  --- &
  UAT chưa có giao dịch chi ngược \\ \addlinespace
Cột \emph{Giờ} &
  \rt{create_date} &
  Ngày thanh toán trong kế toán \textbf{không có giờ}, nên giờ hiển thị lấy từ
  \textbf{lúc phiếu được tạo}, đổi theo múi giờ tài khoản. Hai mốc này có thể lệch nhau. &
  --- & --- \\ \addlinespace
Khối tổng kỳ &
  \rt{account.payment} &
  Tổng \textbf{toàn kỳ} (không phải tổng trang đang xem), \textbf{tách riêng từng loại tiền}
  --- \textbf{không quy đổi tỷ giá}, không cộng gộp số khác loại tiền. &
  --- &
  Kỳ 08/2026 ra \textbf{2 dòng}: \uat{4} giao dịch USD + \uat{1} giao dịch VND \\
\end{hienbang}

\mloc

\begin{itemize}[leftmargin=*]
  \item \textbf{Kỳ} --- chọn trong 6 kỳ; khoá lạ thì lặng lẽ về tháng hiện tại.
  \item \textbf{Tìm kiếm} --- khớp một phần theo \textbf{mã phiếu} hoặc \textbf{nội dung
        thanh toán}. Không tìm theo số tiền, không tìm theo số hoá đơn được gán.
  \item \textbf{Khoảng ngày} (gọi thẳng qua thanh địa chỉ) --- nếu từ ngày lớn hơn đến ngày thì
        hệ thống \textbf{tự hoán đổi hai đầu} rồi chạy tiếp, \textbf{không báo lỗi}. Khác hẳn
        năm màn còn lại của portal (Lịch sử, Giao hàng, Báo cáo, Đổi trả) --- những màn đó dừng
        lại và báo tại thanh lọc.
  \item \textbf{Số dòng mỗi trang} --- 10, 20, 50.
\end{itemize}

\mkiem{một dòng giao dịch}

Không có hành động nào --- bảng chỉ đọc, không mở màn chi tiết.

\mghi

\textbf{Không ghi gì.}

\mhoi

\begin{ask}
\begin{enumerate}[leftmargin=*]
  \item \textbf{Kỳ mở sẵn là tháng hiện tại, và trên UAT tháng đó rỗng.} Toàn bộ \uat{5} giao
        dịch nằm ở tháng 08/2026, còn tháng 09 có \uat{0} $\Rightarrow$ cửa hàng mở màn ra thấy
        trắng, phải tự nghĩ đến việc đổi kỳ. Anh muốn mặc định mở \textbf{kỳ gần nhất có giao
        dịch} không?
  \item \textbf{Ngày ngược thì tự hoán đổi} thay vì báo lỗi --- lệch với bốn màn lọc khác.
        Thống nhất theo cách nào?
  \item \textbf{Tổng tách theo loại tiền, không quy đổi.} Kỳ 08/2026 trên UAT hiện hai dòng
        tổng. Anh xác nhận đây là cách trình bày muốn giữ.
  \item \uat{1}/6 phiếu thanh toán trên UAT \textbf{chưa gắn cửa hàng} nên không cửa hàng nào
        thấy. Cần chặn ở khâu kế toán không?
\end{enumerate}
\end{ask}

% =============================================================================
\section{Màn Chuyển khoản}
\rt{/portal/debt/pay}
\medskip

\mvao

Nút \emph{Thanh toán} ở màn Công nợ. Nhận tham số \rt{week} để biết đang trả cho tuần nào.

\mai

Như hai màn trên: chỉ Chủ tiệm / Quản lý của cửa hàng đang chọn.

\mhien

\begin{hienbang}
Khối số tiền &
  \rt{account.move} &
  Số \emph{còn phải trả} của tuần đang chọn --- tính lại tại chỗ, không tin số trình duyệt
  gửi lên &
  --- &
  Tuần mặc định của TP HCM Q1 đang bằng \uat{0} \\ \addlinespace
Khối tài khoản &
  \rt{res.partner.bank} &
  Tài khoản của công ty có bật cờ nhận tiền qua portal, số thứ tự nhỏ nhất: tên ngân hàng,
  chủ tài khoản, số tài khoản &
  --- &
  \uat{1} tài khoản (Vietcombank -- NGO GIA UAT DEBT) \\ \addlinespace
Nội dung chuyển khoản &
  máy tự sinh &
  \textbf{Mã cửa hàng + K + số tuần + số tiền} (làm tròn xuống, không bao giờ âm). Ví dụ
  tuần 10/08 là tuần ISO 33 $\Rightarrow$ \emph{HCM-01 K33 ...} &
  --- & --- \\ \addlinespace
Ảnh mã QR &
  \textbf{chưa có} &
  Đang chờ BA chốt loại mã (tĩnh hay động). Màn hiện số tài khoản và nội dung, không có ảnh. &
  --- & --- \\
\end{hienbang}

\mloc

Không có ô nhập nào. Chỉ tham số \rt{week}; khoá lạ thì về tuần mặc định.

\mkiem{vào màn này}

\begin{kiembang}
1 & Đã chọn được cửa hàng & Trạng thái rỗng \\
2 & Vai trò Chủ tiệm / Quản lý tại cửa hàng đó & Trang ``không có quyền'' \\
3 & Tuần còn phải trả lớn hơn 0 &
    Quay về màn Công nợ đúng tuần đó kèm băng báo không còn khoản phải trả.
    \textbf{Trên UAT đây là nhánh luôn xảy ra}: tuần mặc định đang dư có \\
4 & Có tài khoản ngân hàng bật portal & Màn báo chưa cấu hình thay vì hiện ô trống \\
\end{kiembang}

\mghi

\textbf{Không ghi gì.} Đây là màn hướng dẫn chuyển khoản, không phải cổng thanh toán --- hệ
thống không tạo phiếu thu, không đánh dấu hoá đơn, không gọi sang ngân hàng.

\mhoi

\begin{ask}
\begin{enumerate}[leftmargin=*]
  \item \textbf{Ảnh mã QR} --- điểm treo lâu nhất của phân hệ. Mã tĩnh (dán sẵn một ảnh cho mọi
        cửa hàng) hay mã động (sinh theo số tiền và nội dung từng tuần)?
  \item \textbf{Nội dung chuyển khoản làm tròn xuống} về số nguyên. Với tiền USD có phần lẻ
        (UAT đang có hoá đơn \$1,15) thì nội dung sẽ ghi \emph{1} --- lệch số thật. Cần giữ
        phần lẻ không?
  \item \textbf{Sau khi cửa hàng chuyển khoản thì portal không biết gì} --- phải đợi kế toán
        Ngô Gia ghi nhận rồi tác vụ nền cập nhật lại. Cửa hàng có cần chỗ báo ``tôi đã chuyển''
        kèm ảnh chụp màn hình không?
\end{enumerate}
\end{ask}

% =============================================================================
\section{Màn Danh sách yêu cầu đổi trả · bù hàng}
\rt{/portal/return}
\medskip

\mvao

Dòng \emph{Đổi trả / Bù hàng} trong bảng \emph{Thêm} (điện thoại); trên Trang chủ có một thẻ số
liệu và một khối danh sách cùng trỏ vào đây. Tiêu đề màn ghi là \emph{Bù hàng}.

\mai

\begin{itemize}[leftmargin=*]
  \item Phải đăng nhập. \textbf{Không chặn vai trò} --- Nhân viên xem và tạo yêu cầu bình
        thường, ngang Chủ tiệm.
  \item Phạm vi là \textbf{mọi cửa hàng tài khoản truy cập được}, \textbf{không} phải cửa hàng
        đang chọn. Người phụ trách ba cửa hàng thấy phiếu của cả ba trộn chung một danh sách,
        và không có bộ lọc theo cửa hàng.
  \item Không truy cập được cửa hàng nào thì màn hiện trạng thái rỗng.
  \item Dữ liệu đọc bằng quyền hệ thống. Hệ thống \textbf{có} khai một luật giới hạn cho người
        dùng portal (chỉ thấy phiếu của các cửa hàng mình tham gia) nhưng luật đó không được
        dùng tới ở màn này; điều kiện lọc trong bảng dưới cho ra đúng cùng kết quả.
\end{itemize}

\mhien

\begin{hienbang}
Danh sách phiếu &
  \rt{wujia.return.request} &
  Phiếu thuộc \textbf{mọi cửa hàng của tôi} (\rt{franchise_id in ...}). Không loại trạng
  thái nào --- kể cả nháp, từ chối, đã huỷ. &
  Ngày tạo mới nhất trước; \textbf{20} phiếu mỗi trang &
  \uat{14} phiếu: Cầu Giấy \uat{5} · Đống Đa \uat{2} · TP HCM Q1 \uat{7} \\ \addlinespace
Huy hiệu trạng thái &
  \rt{state} + tiến độ bù &
  Tám nhãn, trong đó hai nhãn là \textbf{ghép}: \emph{Đang xử lý} gộp hai trạng thái
  (\emph{đang xét} và \emph{đang thực hiện}); \emph{Đang bù một phần} \textbf{không phải một
  trạng thái} mà là ``đang thực hiện + đã bù được một phần''. &
  --- &
  \uat{3} Nháp · \uat{2} Đã gửi · \uat{5} Đã duyệt · \uat{2} Hoàn tất · \uat{2} Từ chối \\ \addlinespace
Cột sản phẩm &
  \rt{sale.order.line} &
  Sản phẩm của \textbf{dòng} trên đơn gốc --- mỗi phiếu đúng \textbf{một} sản phẩm &
  --- &
  \uat{10}/14 phiếu chưa gắn đơn gốc nên cột này trống \\ \addlinespace
Cột đơn gốc · chuyến &
  \rt{sale.order} · \rt{stock.picking.batch} &
  Mã đơn đặt hàng và mã chuyến giao lấy theo đơn đó &
  --- &
  \uat{4}/14 phiếu có đơn gốc \\ \addlinespace
Ảnh minh chứng &
  \rt{ir.attachment} &
  Số ảnh đã đính kèm phiếu &
  --- &
  \uat{1}/14 phiếu có ảnh (3 ảnh) \\
\end{hienbang}

\mloc

\begin{itemize}[leftmargin=*]
  \item \textbf{Trạng thái} --- 8 nhãn kể trên, mặc định \emph{``-- Tất cả trạng thái --''}.
        Giá trị lạ bị bỏ qua.
  \item \textbf{Từ khoá} --- rộng nhất trong portal: khớp một phần theo \textbf{mã yêu cầu},
        \textbf{mã đơn gốc}, \textbf{mã chuyến giao}, \textbf{tên sản phẩm} hoặc \textbf{mã
        sản phẩm}.
  \item \textbf{Khoảng ngày} --- lọc theo \emph{ngày tạo yêu cầu}, quy đổi đúng múi giờ tài
        khoản. Ngày gõ sai định dạng $\rightarrow$ băng báo ở đầu trang; ngày ngược
        $\rightarrow$ báo ngay tại thanh lọc và \textbf{không chạy truy vấn}.
  \item \textbf{Số phiếu mỗi trang} --- 10, 20, 50; mặc định \textbf{20}.
\end{itemize}

\mkiem{Tạo yêu cầu}

Không kiểm gì ở màn danh sách --- sang màn tạo yêu cầu, và màn đó kiểm lại toàn bộ từ đầu.

\mghi

\textbf{Không ghi gì.} Màn chỉ đọc.

\mhoi

\begin{ask}
\begin{enumerate}[leftmargin=*]
  \item \textbf{Danh sách gộp mọi cửa hàng mà không có bộ lọc cửa hàng}, cũng không có cột cửa
        hàng trên bảng. Người phụ trách nhiều cửa hàng khó biết phiếu nào của đâu.
  \item \textbf{Nhãn ``Nháp'' có trong bộ lọc} dù phiếu nháp không gửi được từ portal (xem màn
        sau).
  \item \uat{10}/14 phiếu trên UAT là dữ liệu mẫu nhập thẳng trong ERP, thiếu đơn gốc, sản phẩm
        và loại lỗi --- nếu cần bản chụp màn hình sạch để trình bày thì phải dọn trước.
\end{enumerate}
\end{ask}

% =============================================================================
\section{Màn Tạo yêu cầu đổi trả · bù hàng}
\rt{/portal/return/new}
\medskip

\mvao

Nút \emph{Tạo yêu cầu} ở màn danh sách. Cùng một đường dẫn phục vụ cả việc mở form (mở trang)
lẫn việc gửi form (bấm nút).

\mai

Phải đăng nhập và truy cập được ít nhất một cửa hàng; không thì bị đưa về danh sách kèm băng
báo chưa chọn cửa hàng. \textbf{Không chặn vai trò}.

\mhien

\begin{hienbang}
Ô \emph{Cửa hàng} &
  \rt{wujia.franchise.management} &
  \textbf{Mọi cửa hàng tôi truy cập được}. Chỉ có một cửa hàng thì ô này ẩn đi và gài sẵn. &
  --- &
  Administrator thấy \uat{2} cửa hàng \\ \addlinespace
Ô \emph{Loại lỗi} &
  \rt{wujia.return.issue.type} &
  Các loại lỗi \textbf{đang bật} (\rt{active = true}) &
  Theo số thứ tự cấu hình &
  \uat{5} loại: Hỏng bao bì · Sai sản phẩm · Hết hạn/cận date · Thiếu số lượng · Lỗi khác \\ \addlinespace
Ô \emph{Đơn hàng gốc} &
  \rt{sale.order} &
  Đơn \textbf{đã xác nhận} (\rt{state in (sale, done)}) của \textbf{mọi cửa hàng của tôi}, có
  \textbf{ngày đặt trong 10 ngày} trở lại. \textbf{Không} lọc lại theo cửa hàng vừa chọn ở ô
  bên cạnh. &
  Ngày đặt mới nhất trước &
  \uat{0} đơn tại ngày chụp --- đơn xác nhận gần nhất là 05/09, đã 15 ngày \\ \addlinespace
Ô \emph{Sản phẩm} &
  \rt{sale.order.line} &
  Các dòng hàng của đơn vừa chọn. Máy chủ dựng sẵn danh sách dòng hàng của \textbf{tất cả} đơn
  đủ điều kiện rồi nhúng vào trang, nên đổi đơn là đổi danh sách ngay, không phải gọi lại. &
  Theo thứ tự dòng trong đơn &
  --- \\ \addlinespace
Khối minh chứng &
  \rt{ir.attachment} &
  \textbf{3--5 ảnh} JPG/PNG, mỗi ảnh $\leq$ 5 MB; \textbf{tối đa 1 video} MP4/MOV
  $\leq$ 10 MB; tổng toàn bộ $\leq$ 30 MB &
  --- & --- \\
\end{hienbang}

\mloc

\begin{itemize}[leftmargin=*]
  \item \textbf{Số lượng yêu cầu} --- số thực lớn hơn 0 (khác màn Đặt hàng: ở đây \textbf{cho
        phép số lẻ}, không có bước nhảy, không có trần).
  \item \textbf{Thời gian mở hàng} --- \textbf{bắt buộc}, ô chọn ngày kèm giờ. Hệ thống nhận ba
        kiểu ghi ngày giờ; không đọc được thì báo lỗi và giữ lại những gì đã nhập.
  \item \textbf{Ngày sản xuất} --- không bắt buộc.
  \item \textbf{Ghi chú} --- cắt còn 5.000 ký tự, phần dư bị bỏ lặng lẽ.
  \item \textbf{Hai nút gửi khác nhau giữa hai nền tảng}: máy tính có \emph{Lưu nháp} và
        \emph{Gửi yêu cầu}; điện thoại \textbf{chỉ có} \emph{Gửi}.
\end{itemize}

\mkiem{Gửi yêu cầu}

\begin{kiembang}
1 & Truy cập được ít nhất một cửa hàng & Đưa về danh sách kèm băng báo chưa chọn cửa hàng \\
2 & Cửa hàng gửi lên là số và \textbf{nằm trong danh sách của tôi} &
    ``Cửa hàng không truy cập được.'' \\
3 & Có chọn đơn gốc và sản phẩm, cả hai là số &
    ``Vui lòng chọn đơn hàng gốc và sản phẩm.'' / ``Đơn hàng / sản phẩm không hợp lệ.'' \\
4 & Đơn đó \textbf{vẫn đủ điều kiện} --- kiểm lại ở máy chủ: đúng cửa hàng vừa chọn, đã xác
    nhận, trong 10 ngày &
    ``Đơn hàng không hợp lệ hoặc đã quá thời hạn 10 ngày.'' \\
5 & Dòng sản phẩm thuộc \textbf{đúng đơn đó} &
    ``Sản phẩm phải thuộc đơn hàng gốc của cửa hàng.'' \\
6 & Sản phẩm \textbf{đã được cấu hình chính sách bù hàng} (bật cờ bù, có đơn vị khai hao hụt,
    và tỷ lệ quy đổi hợp lệ) &
    ``Sản phẩm chưa được cấu hình chính sách bù hàng. Vui lòng liên hệ Ngô Gia.'' ---
    trên UAT \uat{4}/5 sản phẩm portal rơi vào nhánh này \\
7 & Loại lỗi còn bật & ``Vui lòng chọn loại lỗi.'' \\
8 & Số lượng lớn hơn 0 & ``Số lượng yêu cầu phải lớn hơn 0.'' \\
9 & Thời gian mở hàng đọc được & ``Vui lòng nhập thời gian mở hàng hợp lệ.'' \\
10 & Số ảnh trong khoảng \textbf{3--5} &
    ``Cần tải từ 3 đến 5 ảnh minh chứng.'' --- bấm \emph{Lưu nháp} thì \textbf{bỏ qua mức tối
    thiểu}, vẫn chặn mức tối đa \\
11 & Nhiều nhất 1 video & ``Chỉ được tải tối đa 1 video minh chứng.'' \\
12 & Mỗi tệp \textbf{đúng định dạng thật} và đúng dung lượng &
    ``Tệp không đúng định dạng hoặc vượt quá dung lượng cho phép.'' --- cùng một câu cho cả hai
    lý do, cố ý không cho biết đã sai cái nào \\
13 & Tổng dung lượng $\leq$ 30 MB & ``Tổng dung lượng minh chứng không được vượt quá 30 MB.'' \\
14 & Ghi được phiếu &
    ``Không thể gửi yêu cầu. Vui lòng kiểm tra lại thông tin và thử lại.'' ---
    \textbf{không} lộ chi tiết kỹ thuật ra màn \\
15 & Gắn được tệp vào phiếu &
    Phiếu vừa tạo \textbf{bị xoá hẳn}, quay lại form kèm lỗi --- không để lại phiếu rỗng \\
16 & (chỉ khi bấm Gửi) Đủ \textbf{3 ảnh} lúc chuyển sang \emph{Đã gửi} &
    Phiếu \textbf{bị xoá hẳn}, quay lại form \\
\end{kiembang}

\textbf{Về bước 12:} hệ thống đọc \textbf{nội dung thật} của tệp để biết đó là ảnh hay video,
không tin đuôi tệp và không tin loại tệp trình duyệt khai. Đổi tên một tệp bất kỳ thành
\rt{.jpg} không qua được bước này.

\mghi

\begin{itemize}[leftmargin=*]
  \item Tạo \textbf{một} bản ghi yêu cầu (\rt{wujia.return.request}), mã tự sinh theo dãy số
        (UAT đang ở \emph{RTN/26/00004}).
  \item Trạng thái ban đầu luôn là \textbf{Nháp}; bấm \emph{Gửi} thì hệ thống chuyển tiếp sang
        \textbf{Đã gửi} và ghi một dòng \emph{``Yêu cầu đã được gửi''} vào nhật ký trao đổi của
        phiếu.
  \item Người tạo = người đang đăng nhập; ngày tạo = lúc bấm nút. Cửa hàng ghi trên phiếu là
        cửa hàng chọn trong form, \textbf{không} phải cửa hàng đang chọn ở cookie.
  \item Đơn vị tính dùng để khai hao hụt lấy theo \textbf{cấu hình bù của sản phẩm}; sản phẩm
        chưa khai thì lùi về đơn vị ghi trên đơn gốc.
  \item Tạo các bản ghi tệp đính kèm, tách riêng nhóm ảnh và nhóm video.
  \item \textbf{Không} tạo đơn hàng, \textbf{không} động vào kho, \textbf{không} gửi email.
  \item Bất kỳ bước nào sau khi tạo mà hỏng thì phiếu bị \textbf{xoá hẳn} --- hệ thống không để
        lại phiếu mồ côi hay tệp mồ côi.
\end{itemize}

\mhoi

\begin{ask}
\begin{enumerate}[leftmargin=*]
  \item \textbf{Cửa sổ 10 ngày đang làm màn này không dùng được.} Mốc tính từ \emph{ngày đặt}
        của đơn đã xác nhận, đếm \textbf{theo giờ} chứ không theo ngày. Trên UAT có \uat{14} đơn
        đã xác nhận nhưng \uat{0} đơn lọt cửa sổ (đơn gần nhất 05/09, cách 15 ngày)
        $\Rightarrow$ \textbf{không ai tạo được yêu cầu nào}, ô chọn đơn trống kèm dòng đỏ.
        Anh chốt giúp: giữ 10 ngày hay nới ra, và tính từ ngày đặt hay ngày \textbf{giao hàng}?
  \item \textbf{Ô chọn đơn không lọc theo cửa hàng ở ô bên cạnh.} Người nhiều cửa hàng chọn
        lệch cặp cửa hàng--đơn sẽ nhận câu ``Đơn hàng không hợp lệ hoặc đã quá thời hạn 10
        ngày'' --- câu này \textbf{nói sai nguyên nhân}.
  \item \textbf{Phiếu nháp là ngõ cụt.} Máy tính có nút \emph{Lưu nháp}, nhưng portal không có
        chỗ nào gửi phiếu nháp đã lưu: màn chi tiết chỉ đọc. UAT đang có \uat{3} phiếu kẹt.
        Anh muốn (a) bỏ nút Lưu nháp, hay (b) thêm nút \emph{Gửi} ở màn chi tiết?
  \item \textbf{Điện thoại không có nút Lưu nháp} --- hai nền tảng khác luật cho cùng một việc.
  \item \textbf{Sản phẩm chưa cấu hình bù thì chặn cả đổi và trả}, dù hai phương án đó không
        cần chính sách bù. Trên UAT \uat{4}/5 sản phẩm đang hiện trên màn Đặt hàng rơi vào đây
        --- cửa hàng đặt được nhưng không khiếu nại được.
  \item \textbf{Mỗi phiếu đúng một sản phẩm.} Một đơn hỏng ba mặt hàng phải làm ba phiếu, mỗi
        phiếu tự chụp đủ 3 ảnh. Anh xác nhận giữ nguyên.
  \item \textbf{Video tải lên không bao giờ hiện lại cho cửa hàng.} Hệ thống nhận, kiểm và lưu
        video, nhân sự Ngô Gia xem được trong ERP, nhưng màn chi tiết trên portal chỉ hiện ảnh.
        Cửa hàng không có cách nào xem lại video mình vừa gửi.
  \item \textbf{Cửa hàng không sửa và không huỷ được phiếu sau khi gửi} --- portal không có nút
        nào cho việc đó.
\end{enumerate}
\end{ask}

% =============================================================================
\section{Màn Chi tiết yêu cầu}
\rt{/portal/return/<int:request_id>}
\medskip

\mvao

Bấm một phiếu ở màn danh sách; hoặc được đưa tới ngay sau khi tạo yêu cầu thành công (kèm băng
báo vừa tạo xong).

\mai

Phiếu phải thuộc \textbf{một trong các cửa hàng của tôi}. Không thoả thì bị đưa về danh sách
kèm băng \emph{không tìm thấy} --- \textbf{cùng một câu} cho cả trường hợp phiếu không tồn tại
lẫn trường hợp phiếu của cửa hàng khác, cố ý không cho dò số phiếu.

\mhien

\begin{hienbang}
Đầu phiếu &
  \rt{wujia.return.request} &
  Mã yêu cầu, ngày tạo, cửa hàng, huy hiệu trạng thái (tám nhãn như màn danh sách) &
  --- & --- \\ \addlinespace
Khối nội dung &
  \rt{wujia.return.request} &
  Loại lỗi, thời gian mở hàng, ngày sản xuất, số lượng yêu cầu, ghi chú của cửa hàng &
  --- & --- \\ \addlinespace
Khối minh chứng &
  \rt{ir.attachment} &
  \textbf{Chỉ ảnh}; mỗi ảnh mở bằng đường dẫn tải riêng của phiếu. Video \textbf{không} hiện. &
  --- &
  \uat{1}/14 phiếu có ảnh \\ \addlinespace
Khối phản hồi &
  \rt{reject_reason} · \rt{approval_note} &
  \textbf{Lý do từ chối} khi bị từ chối, hoặc \textbf{ghi chú duyệt} khi đã duyệt. Ghi chú nội
  bộ của Ngô Gia \textbf{không} lộ ra portal. Chưa có gì thì ghi ``chưa có phản hồi''. &
  --- & --- \\ \addlinespace
Khối tiến độ bù &
  \rt{wujia.compensation.allocation} &
  Chỉ hiện khi Ngô Gia \textbf{đã chốt phương án}. Với phương án \emph{Bù hàng}: số lượng
  duyệt, đã lên đơn, đã bù, còn thiếu, thanh phần trăm (\textbf{chặn trần 100\%}), và danh sách
  đơn bù kèm nhãn giao hàng. &
  --- &
  \uat{0} bản ghi phân bổ $\Rightarrow$ khối luôn trống, kể cả \uat{3} phiếu đã chốt phương án
  bù \\ \addlinespace
Nhãn giao của đơn bù &
  \rt{stock.picking} &
  Năm nhãn: \emph{Chưa tạo phiếu giao} · \emph{Chưa giao} · \emph{Đã giao n/m phiếu} ·
  \emph{Đã giao đủ} · \emph{Đơn bù đã bị hủy} &
  --- & --- \\
\end{hienbang}

Khi \textbf{mọi đơn bù đều đã bị huỷ}, màn hiện thêm một cảnh báo riêng --- để cửa hàng biết
quyền lợi đã đóng lại và phải tạo yêu cầu mới, thay vì ngồi chờ hàng.

\mloc

Không có ô lọc hay ô nhập nào.

\mkiem{bất kỳ nút nào}

Màn \textbf{chỉ đọc} --- không có nút nào ghi dữ liệu. Duy nhất các ảnh minh chứng là liên kết
mở sang đường dẫn tải (màn sau).

\mghi

\textbf{Không ghi gì.}

\mhoi

\begin{ask}
\begin{enumerate}[leftmargin=*]
  \item \textbf{Cửa hàng không có nút nào ở màn này} --- không gửi được phiếu nháp, không huỷ,
        không bổ sung ảnh, không trả lời lại Ngô Gia. Anh muốn thêm gì trước khi chạy thật?
  \item \textbf{Thanh tiến độ bù chưa chạy được trên UAT} vì chưa có bản ghi phân bổ nào. Cần
        dựng một ca mẫu (duyệt $\rightarrow$ lên đơn bù $\rightarrow$ giao một phần) để nghiệm
        thu đúng nhãn \emph{Đang bù một phần}.
  \item \textbf{Thanh phần trăm bị chặn trần 100\%} --- bù dư hơn số duyệt vẫn hiện đúng 100\%.
\end{enumerate}
\end{ask}

% =============================================================================
\section{Xem · tải minh chứng}
\rt{/portal/return/<int:request_id>/attachment/<int:att_id>}
\medskip

\mvao

Không phải màn --- đây là đường dẫn trả về chính tệp ảnh, gọi ngầm khi màn chi tiết hiện ảnh
thu nhỏ, và mở ra tab mới khi người dùng bấm vào ảnh.

\mai

Ba lớp kiểm, theo thứ tự:

\begin{kiembang}
1 & Truy cập được ít nhất một cửa hàng & Từ chối truy cập \\
2 & Phiếu tồn tại và thuộc \textbf{cửa hàng của tôi} & Báo không tìm thấy \\
3 & Tệp \textbf{thuộc đúng phiếu đó} (nằm trong nhóm ảnh hoặc nhóm video của phiếu) &
    Từ chối truy cập --- \textbf{không} cho mượn đường dẫn này để đọc tệp bất kỳ trong hệ thống \\
4 & Tệp còn tồn tại & Báo không tìm thấy \\
\end{kiembang}

\mhien

Trả về đúng nội dung tệp, đặt chế độ \textbf{xem thẳng trong trình duyệt} chứ không ép tải
xuống. Video được phép đi qua đường dẫn này nhưng \textbf{không màn nào trỏ tới} --- xem mục 7
màn Tạo yêu cầu.

\mloc

Không có.

\mkiem{một ảnh}

Xem bảng ba lớp ở trên.

\mghi

\textbf{Không ghi gì.}

\mhoi

\begin{ask}
Ảnh minh chứng đi qua đường dẫn có kiểm quyền nên \textbf{không} đọc được nếu chỉ có số tệp.
Không có điểm nào cần BA chốt ở đường dẫn này.
\end{ask}

% =============================================================================
\section{Đường dẫn cũ của phân hệ Đổi trả}
\rt{/portal/return-request-list}
\medskip

Đường dẫn từ bản v14, được giữ lại để chuyển hướng \textbf{vĩnh viễn} sang \rt{/portal/return}.
Khoảng 1.500 tài khoản nhượng quyền có thể còn dấu trang hoặc email cũ trỏ vào đó.

Đường dẫn này \textbf{không đòi đăng nhập} --- ai mở cũng được chuyển hướng, rồi mới bị màn
đích đòi đăng nhập. Đây là cố ý: chuyển hướng trước, hỏi đăng nhập sau, để người dùng không
mất đường dẫn đích sau khi đăng nhập xong.

\begin{ask}
Không có điểm cần chốt --- nhưng nếu BA muốn khai tử đường dẫn cũ sau một mốc thời gian nào đó
thì cho biết mốc, dev sẽ gỡ.
\end{ask}

\chapter{Hỗ trợ --- Yêu cầu thông tin --- Kiến thức --- Thông báo}

Chương này gồm \textbf{22 đường dẫn} gom thành \textbf{19 mục}, thuộc bốn mô-đun:
\rt{wujia_portal_support} (5), \rt{wujia_portal_info_request} (5),
\rt{wujia_portal_knowledge} (4), \rt{wujia_portal_notification} (8).

\textbf{Bốn điểm cần biết trước khi đọc:}

\begin{enumerate}[leftmargin=*]
  \item \textbf{Phiếu Hỗ trợ lọc theo NGƯỜI TẠO, không theo cửa hàng.} Đây là phân hệ duy nhất
        trong portal làm như vậy. Đồng nghiệp cùng cửa hàng \textbf{không thấy} phiếu của nhau;
        Chủ tiệm \textbf{không thấy} phiếu nhân viên mình gửi. Trong khi đó Yêu cầu thông tin
        (ngay chương này) lại lọc theo cửa hàng nên cả cửa hàng cùng thấy. Hai phân hệ giống
        nhau về hình thức nhưng ngược nhau về phạm vi.
  \item \textbf{Kiến thức không phân biệt cửa hàng và không phân biệt vai trò.} Mọi tài khoản
        portal thấy đúng cùng một tập bài viết. Đây là phân hệ duy nhất như vậy.
  \item \textbf{Thông báo ghi ``đã đọc'' theo bộ ba người + thông báo + cửa hàng.} Cùng một
        người, đọc ở cửa hàng A rồi bấm sang cửa hàng B thì thông báo đó \textbf{lại thành chưa
        đọc}. UAT hiện có \uat{6} dấu đã đọc cũ \textbf{chưa gắn cửa hàng} (dữ liệu trước khi đổi
        cách ghi) --- những dấu này không khớp với bất kỳ cửa hàng nào.
  \item \textbf{Yêu cầu cập nhật thông tin chưa từng được dùng.} UAT có \uat{0} bản ghi. Mọi mô
        tả trong phần đó là đọc từ mã nguồn, chưa có số thật để đối chiếu.
\end{enumerate}

Quy ước chung toàn portal xem mục~\ref{sec:quy-uoc-chung}.

% =============================================================================
\section{Màn Danh sách phiếu hỗ trợ}
\rt{/portal/support}
\medskip

\mvao

Mục \emph{Hỗ trợ} trên thanh điều hướng. Tham số: \rt{page}, \rt{state}, \rt{q}.

\mai

\begin{itemize}[leftmargin=*]
  \item Phải đăng nhập. \textbf{Không chặn vai trò.}
  \item Chỉ thấy phiếu \textbf{do chính tài khoản mình tạo} (\rt{created_by_id = tôi}).
        \textbf{Không} lọc theo cửa hàng --- nghĩa là người rời cửa hàng vẫn thấy phiếu cũ của
        mình, còn người mới vào cửa hàng không thấy gì.
  \item Phiếu bị Ngô Gia tắt cờ hiện trên portal thì \textbf{biến mất hoàn toàn}, kể cả với
        người đã tạo ra nó.
  \item Thực tế UAT: \uat{15} phiếu, \uat{0} phiếu do tài khoản cửa hàng thật tạo --- toàn bộ
        do tài khoản nội bộ tạo. Vì vậy \textbf{mọi tài khoản cửa hàng trên UAT đều mở ra màn
        rỗng}, và bản chụp màn hình phải làm bằng tài khoản Administrator.
\end{itemize}

\mhien

\begin{hienbang}
Danh sách phiếu &
  \rt{wujia.support.ticket} &
  Phiếu do tôi tạo, còn bật cờ hiện trên portal, và \textbf{không phải phiếu đã huỷ}
  (\rt{state != 'cancelled'}) &
  Ngày tạo mới nhất trước; \textbf{20} phiếu mỗi trang, \textbf{cố định} &
  \uat{15} phiếu toàn hệ, \uat{2} đã huỷ (bị giấu) \\ \addlinespace
Huy hiệu trạng thái &
  \rt{state} &
  Sáu nhãn: \emph{Mới} · \emph{Đang xử lý} · \emph{Chờ phản hồi} · \emph{Đã giải quyết} ·
  \emph{Đã đóng} · \emph{Đã huỷ} (nhãn cuối chỉ gặp ở màn chi tiết) &
  --- &
  \uat{4} Mới · \uat{3} Đang xử lý · \uat{1} Chờ phản hồi · \uat{3} Đã giải quyết ·
  \uat{2} Đã đóng \\ \addlinespace
Huy hiệu mức độ &
  \rt{priority} &
  Hai mức: \emph{Bình thường} và \emph{Khẩn} &
  --- &
  \uat{4}/15 phiếu đánh dấu Khẩn \\ \addlinespace
Cột danh mục &
  \rt{wujia.support.category} &
  Danh mục do Ngô Gia cấu hình &
  --- &
  \uat{7} danh mục đang bật \\
\end{hienbang}

\mloc

\begin{itemize}[leftmargin=*]
  \item \textbf{Trạng thái} --- một trong sáu nhãn; giá trị lạ bị bỏ qua.
  \item \textbf{Từ khoá} --- khớp một phần theo \textbf{mã phiếu} hoặc \textbf{tiêu đề}. Không
        tìm trong nội dung mô tả, không tìm trong các lời nhắn.
  \item \textbf{Không có bộ lọc ngày} và \textbf{không đổi được số dòng mỗi trang} --- khác mọi
        màn danh sách khác trong portal.
\end{itemize}

\mkiem{Tạo phiếu}

Không kiểm gì ở màn danh sách.

\mghi

\textbf{Không ghi gì.}

\mhoi

\begin{ask}
\begin{enumerate}[leftmargin=*]
  \item \textbf{Phiếu chỉ người tạo mới thấy.} Nhân viên nghỉ việc thì phiếu đi theo tài khoản
        đó; Chủ tiệm không nắm được cửa hàng mình đã gửi những gì. Anh muốn đổi sang
        \textbf{lọc theo cửa hàng} như các phân hệ còn lại không? Đây là điểm đã nêu ở bản tài
        liệu cũ và tới nay chưa chốt.
  \item \textbf{Không có bộ lọc ngày và cỡ trang cố định 20} --- lệch với năm màn danh sách
        khác. Có cần đồng bộ không?
  \item \textbf{Phiếu Đã huỷ bị giấu khỏi danh sách nhưng màn chi tiết vẫn mở được} nếu còn giữ
        đường dẫn. UAT có \uat{2} phiếu như vậy.
\end{enumerate}
\end{ask}

% =============================================================================
\section{Màn Tạo phiếu hỗ trợ}
\rt{/portal/support/new}
\medskip

\mvao

Nút \emph{Tạo phiếu} ở màn danh sách. Cùng một đường dẫn dùng cho cả mở form lẫn gửi form.

\mai

Phải đăng nhập và \textbf{truy cập được ít nhất một cửa hàng}; không thì bị đưa thẳng về danh
sách (không có thông báo gì). Không chặn vai trò.

\mhien

\begin{hienbang}
Ô \emph{Cửa hàng} &
  \rt{wujia.franchise.management} &
  Mọi cửa hàng tôi truy cập được &
  --- & --- \\ \addlinespace
Ô \emph{Danh mục} &
  \rt{wujia.support.category} &
  Danh mục đang bật (\rt{active = true}) &
  Theo số thứ tự, rồi tên &
  \uat{7} danh mục \\ \addlinespace
Ô \emph{Mức độ} &
  --- &
  Hai lựa chọn: \emph{Bình thường} (mặc định) và \emph{Khẩn} &
  --- & --- \\ \addlinespace
Khối tệp đính kèm &
  \rt{ir.attachment} &
  Tối đa \textbf{6 tệp}, mỗi tệp $\leq$ \textbf{5 MB}; chỉ nhận PNG, JPEG, WEBP, PDF &
  --- & --- \\
\end{hienbang}

\mloc

\begin{itemize}[leftmargin=*]
  \item \textbf{Tiêu đề} --- bắt buộc. Hệ thống nhận cả tên ô cũ (\rt{subject}) lẫn tên ô mới
        (\rt{title}) để không làm hỏng liên kết cũ.
  \item \textbf{Mô tả} --- không bắt buộc, không giới hạn độ dài ở tầng portal.
  \item \textbf{Mức độ} --- giá trị lạ tự hạ về \emph{Bình thường}.
\end{itemize}

\mkiem{Gửi phiếu}

\begin{kiembang}
1 & Cửa hàng và danh mục gửi lên là số & Quay lại form kèm mã lỗi \rt{invalid_input} \\
2 & Có tiêu đề, có cửa hàng, có danh mục & Quay lại form kèm \rt{missing_fields} \\
3 & Cửa hàng \textbf{nằm trong danh sách của tôi} & Quay lại form kèm \rt{invalid_franchise} \\
4 & Danh mục còn đang bật & Quay lại form kèm \rt{invalid_input} \\
5 & Tệp đính kèm đúng số lượng, dung lượng, định dạng &
    Phiếu vừa tạo \textbf{bị xoá hẳn}, quay lại form kèm \rt{invalid_attachment} \\
\end{kiembang}

Bốn bước đầu chạy \textbf{trước} khi ghi; bước 5 chạy \textbf{sau} khi đã ghi phiếu, nên hệ
thống phải xoá lại phiếu. Các mã lỗi trên được trả về dưới dạng tham số trên thanh địa chỉ và
màn dịch ra câu tiếng Việt.

\mghi

\begin{itemize}[leftmargin=*]
  \item Tạo một phiếu \rt{wujia.support.ticket}, mã tự sinh (UAT đang ở \emph{WJ-TK/26/00016}),
        trạng thái \textbf{Mới}, cờ hiện trên portal \textbf{bật sẵn}.
  \item Người tạo = người đang đăng nhập; cửa hàng = cửa hàng chọn trong form.
  \item Tạo các bản ghi tệp đính kèm gắn vào phiếu.
  \item \textbf{Không} gửi email, \textbf{không} tạo việc cần làm cho ai. Phiếu chỉ hiện trong
        ERP để Ngô Gia tự vào xem.
\end{itemize}

\mhoi

\begin{ask}
\begin{enumerate}[leftmargin=*]
  \item \textbf{Định dạng tệp chỉ kiểm theo lời khai của trình duyệt}, không đọc nội dung thật
        --- khác với màn Đổi trả (chương 6) đã đọc nội dung. Đây là mục đã nằm trong danh sách
        chuẩn hoá F1 của dev; ghi ở đây để BA biết, không cần ra task riêng.
  \item \textbf{Gửi phiếu xong không ai được báo.} Không email, không thông báo trong ERP. Ngô
        Gia phải chủ động mở danh sách. Anh có muốn gắn thông báo cho người phụ trách không?
  \item \textbf{Chưa có giới hạn số phiếu gửi trong ngày} --- một tài khoản bấm liên tục sẽ tạo
        liên tục.
\end{enumerate}
\end{ask}

% =============================================================================
\section{Màn Chi tiết phiếu hỗ trợ}
\rt{/portal/support/<int:ticket_id>}
\medskip

\mvao

Bấm một dòng ở màn danh sách, hoặc được đưa tới ngay sau khi tạo phiếu.

\mai

Phiếu phải \textbf{do chính tôi tạo} và \textbf{còn bật cờ hiện trên portal}. Sai thì bị đưa về
danh sách, không có thông báo. \textbf{Lưu ý}: màn này \textbf{không} loại phiếu đã huỷ --- phiếu
huỷ bị giấu ở danh sách nhưng vẫn mở được bằng đường dẫn.

\mhien

\begin{hienbang}
Đầu phiếu &
  \rt{wujia.support.ticket} &
  Mã phiếu, tiêu đề, danh mục, mức độ, trạng thái, cửa hàng, ngày tạo &
  --- & --- \\ \addlinespace
Nội dung &
  \rt{description} &
  Mô tả người dùng gõ lúc tạo &
  --- & --- \\ \addlinespace
Khối trao đổi &
  \rt{mail.message} &
  \textbf{Mọi lời nhắn} trên phiếu (loại bỏ các dòng nhật ký hệ thống). Nghĩa là lời Ngô Gia gõ
  vào ô trao đổi trong ERP \textbf{hiện thẳng} cho cửa hàng. Ô \emph{Ghi chú nội bộ} thì không. &
  Cũ trước, mới sau &
  --- \\ \addlinespace
Tệp đính kèm &
  \rt{ir.attachment} &
  Tệp gắn vào phiếu, tải qua đường dẫn có kiểm quyền &
  --- &
  \uat{1} tệp trên toàn bộ 15 phiếu \\
\end{hienbang}

\mloc

Chỉ một ô: khung soạn lời nhắn.

\mkiem{Gửi lời nhắn}

Xem mục kế tiếp.

\mghi

Mở màn \textbf{không ghi gì}.

\mhoi

\begin{ask}
\begin{enumerate}[leftmargin=*]
  \item \textbf{Ranh giới nội bộ / công khai nằm ở chỗ nhân sự Ngô Gia chọn ô nào.} Gõ vào ô
        trao đổi thì cửa hàng đọc được; muốn ghi riêng phải dùng ô \emph{Ghi chú nội bộ} hoặc
        chức năng ghi chú. Cần dặn lại đội vận hành điểm này.
  \item \textbf{Phiếu Đã đóng / Đã huỷ vẫn nhận được lời nhắn mới} --- không có chốt chặn theo
        trạng thái.
\end{enumerate}
\end{ask}

% =============================================================================
\section{Gửi lời nhắn vào phiếu hỗ trợ}
\rt{/portal/support/<int:ticket_id>/reply}
\medskip

\mvao

Không phải màn --- là hành động của nút \emph{Gửi} trong khung trao đổi ở màn chi tiết.

\mai

Cùng điều kiện với màn chi tiết: phiếu do tôi tạo và còn hiện trên portal.

\mhien

Không hiện gì --- xong việc thì quay lại chính màn chi tiết.

\mloc

Đúng một ô nội dung. \textbf{Không} đính kèm được tệp trong lời nhắn.

\mkiem{Gửi}

\begin{kiembang}
1 & Phiếu tồn tại, do tôi tạo, còn hiện trên portal & Đưa về danh sách \\
2 & Nội dung \textbf{không rỗng} sau khi bỏ khoảng trắng &
    \textbf{Không báo lỗi} --- lặng lẽ quay về màn chi tiết, không ghi gì \\
\end{kiembang}

\mghi

\begin{itemize}[leftmargin=*]
  \item Ghi một lời nhắn vào phiếu, \textbf{đứng tên chính người dùng} (không phải tên hệ
        thống), loại \emph{bình luận} nên Ngô Gia nhận được trong ERP.
  \item Việc ghi lời nhắn \textbf{kích hoạt phép tính thời gian phản hồi} của phiếu --- dùng
        cho báo cáo nội bộ.
  \item \textbf{Không} đổi trạng thái phiếu. Phiếu đang \emph{Chờ phản hồi} mà cửa hàng trả lời
        thì vẫn nằm nguyên ở \emph{Chờ phản hồi} cho tới khi Ngô Gia tự đổi.
\end{itemize}

\mhoi

\begin{ask}
\begin{enumerate}[leftmargin=*]
  \item \textbf{Trả lời không tự đổi trạng thái.} Anh có muốn phiếu \emph{Chờ phản hồi} tự về
        \emph{Đang xử lý} khi cửa hàng trả lời không?
  \item \textbf{Không đính kèm được ảnh trong lời nhắn} --- cửa hàng muốn gửi thêm ảnh chụp
        phải mở phiếu mới.
\end{enumerate}
\end{ask}

% =============================================================================
\section{Tải tệp đính kèm của phiếu hỗ trợ}
\rt{/portal/support/<int:ticket_id>/attachment/<int:att_id>}
\medskip

\mvao

Bấm tên tệp ở màn chi tiết.

\mai

\begin{kiembang}
1 & Phiếu tồn tại, do tôi tạo, còn hiện trên portal & Báo không tìm thấy \\
2 & Tệp \textbf{thuộc đúng phiếu đó} &
    Từ chối truy cập --- không mượn được đường dẫn này để đọc tệp khác trong hệ thống \\
\end{kiembang}

\mhien

Trả về tệp ở chế độ \textbf{tải xuống} (khác Đổi trả --- bên đó mở thẳng trong trình duyệt).

\mloc

Không có.

\mkiem{tên tệp}

Xem bảng trên.

\mghi

\textbf{Không ghi gì.}

\mhoi

\begin{ask}
Không có điểm cần chốt.
\end{ask}

% =============================================================================
\section{Màn Danh sách yêu cầu cập nhật thông tin}
\rt{/portal/info-request}
\medskip

\mvao

Mục \emph{Yêu cầu cập nhật thông tin} trên thanh điều hướng. Tham số: \rt{page},
\rt{state}, \rt{request_type}, \rt{page_size}.

\mai

\begin{itemize}[leftmargin=*]
  \item Phải đăng nhập. \textbf{Màn danh sách không chặn vai trò} --- Nhân viên xem được;
        chỉ màn \emph{tạo} mới chặn.
  \item Phạm vi là \textbf{mọi cửa hàng tôi truy cập được}, không lọc theo người tạo --- ngược
        hẳn với phân hệ Hỗ trợ ở đầu chương.
  \item Không truy cập được cửa hàng nào thì màn rỗng.
\end{itemize}

\mhien

\begin{hienbang}
Danh sách yêu cầu &
  \rt{wujia.info.update.request} &
  Yêu cầu thuộc mọi cửa hàng của tôi &
  Ngày tạo mới nhất trước; \textbf{20} yêu cầu mỗi trang (10/20/50) &
  \uat{0} bản ghi --- phân hệ chưa được dùng lần nào \\ \addlinespace
Huy hiệu trạng thái &
  \rt{state} &
  Năm nhãn: \emph{Nháp} · \emph{Đã gửi} · \emph{Đang xem} · \emph{Đã duyệt} · \emph{Từ chối}.
  \textbf{Không có} nhãn ``Đã huỷ''. &
  --- & --- \\ \addlinespace
Cột loại thông tin &
  \rt{request_type} &
  Bảy loại cố định trong mã nguồn: Địa chỉ · Số điện thoại · Email · Tên chủ cửa hàng ·
  Thông tin ngân hàng · Người đại diện · Khác &
  --- & --- \\
\end{hienbang}

\mloc

Lọc theo \textbf{trạng thái} và \textbf{loại thông tin}; giá trị lạ bị bỏ qua. Số dòng mỗi
trang 10/20/50. \textbf{Không} có ô tìm kiếm, \textbf{không} có bộ lọc ngày.

\mkiem{Tạo yêu cầu}

Xem màn kế tiếp.

\mghi

\textbf{Không ghi gì.}

\mhoi

\begin{ask}
\begin{enumerate}[leftmargin=*]
  \item \textbf{Nhân viên xem được mọi yêu cầu của cửa hàng} (kể cả yêu cầu đổi thông tin ngân
        hàng) nhưng \textbf{không tạo được}. Đúng ý anh chứ?
  \item \textbf{Không có nhãn ``Đã huỷ''} --- yêu cầu người dùng tự huỷ sẽ hiện thành
        \emph{Từ chối}; xem mục Huỷ yêu cầu bên dưới.
  \item \textbf{Bảy loại thông tin đang cố định trong mã nguồn}, muốn thêm loại phải sửa mã.
        Anh xác nhận danh sách này đủ.
\end{enumerate}
\end{ask}

% =============================================================================
\section{Màn Tạo yêu cầu cập nhật thông tin}
\rt{/portal/info-request/new}
\medskip

\mvao

Nút \emph{Tạo yêu cầu} ở màn danh sách. Cùng đường dẫn cho cả mở form lẫn gửi form.

\mai

\begin{itemize}[leftmargin=*]
  \item Phải truy cập được ít nhất một cửa hàng.
  \item \textbf{Chỉ Chủ tiệm / Quản lý}. Nhân viên nhận \textbf{trang lỗi 403 thô của máy chủ},
        \textbf{không} phải trang thông báo thân thiện như phân hệ Công nợ. Hai phân hệ cùng
        chặn vai trò nhưng chặn ra hai kiểu khác nhau.
  \item Vai trò xét theo \textbf{mức cao nhất trên MỌI cửa hàng} của tài khoản --- giống màn Báo
        cáo (chương 5), khác Công nợ (chương 6, xét tại cửa hàng đang chọn).
\end{itemize}

\mhien

\begin{hienbang}
Ô \emph{Cửa hàng} &
  \rt{wujia.franchise.management} &
  Mọi cửa hàng tôi truy cập được &
  --- & --- \\ \addlinespace
Ô \emph{Loại thông tin} &
  --- &
  Bảy loại cố định &
  Theo thứ tự khai trong mã &
  --- \\ \addlinespace
Ô \emph{Giá trị hiện tại} &
  \rt{wujia.franchise.management} &
  Đọc thẳng từ cửa hàng theo loại thông tin đang chọn; gọi ngầm mỗi lần đổi ô
  (xem mục \emph{Lấy giá trị hiện tại}) &
  --- & --- \\ \addlinespace
Khối tệp đính kèm &
  \rt{ir.attachment} &
  Tối đa \textbf{6 tệp}, mỗi tệp $\leq$ \textbf{5 MB}; PNG, JPEG, WEBP, PDF &
  --- & --- \\
\end{hienbang}

\mloc

\begin{itemize}[leftmargin=*]
  \item \textbf{Giá trị mới} --- bắt buộc, cắt còn 5.000 ký tự.
  \item \textbf{Tên trường cần đổi} --- chỉ hiện khi chọn loại \emph{Khác}; bắt buộc khi đó.
        Đây là \textbf{tên kỹ thuật} của trường, không phải tên tiếng Việt.
  \item \textbf{Ghi chú} --- cắt còn 2.000 ký tự.
  \item \textbf{Mức độ} --- Bình thường / Khẩn; giá trị lạ hạ về Bình thường.
  \item \textbf{Hai nút}: \emph{Lưu nháp} và \emph{Gửi}.
\end{itemize}

\mkiem{Gửi yêu cầu}

\begin{kiembang}
1 & Truy cập được ít nhất một cửa hàng & Đưa về danh sách \\
2 & Vai trò Chủ tiệm / Quản lý (tính trên mọi cửa hàng) & Trang lỗi 403 của máy chủ \\
3 & Cửa hàng gửi lên là số và nằm trong danh sách của tôi &
    ``Cửa hàng không hợp lệ.'' / ``Cửa hàng không truy cập được.'' \\
4 & Loại thông tin nằm trong bảy loại & ``Loại thông tin không hợp lệ.'' \\
5 & Có nhập giá trị mới & ``Vui lòng nhập giá trị mới.'' \\
6 & Chọn \emph{Khác} thì phải nhập tên trường & ``Khi chọn 'Khác' phải nhập tên field.'' \\
7 & Ghi được bản ghi &
    ``Không thể gửi yêu cầu. Vui lòng kiểm tra lại thông tin và thử lại.'' \\
8 & Tệp đính kèm hợp lệ & Bản ghi vừa tạo \textbf{bị xoá hẳn}, quay lại form \\
\end{kiembang}

\mghi

\begin{itemize}[leftmargin=*]
  \item Tạo một bản ghi \rt{wujia.info.update.request}, mã tự sinh, trạng thái \textbf{Nháp}.
  \item Bấm \emph{Gửi} thì chuyển sang \textbf{Đã gửi}, ghi mốc thời gian gửi và một dòng vào
        khung trao đổi của bản ghi.
  \item \textbf{Không} sửa gì trên hồ sơ cửa hàng --- đây chỉ là đề nghị.
\end{itemize}

\mhoi

\begin{ask}
\begin{enumerate}[leftmargin=*]
  \item \textbf{Quyền xét trên mọi cửa hàng nhưng gửi được cho từng cửa hàng.} Một người là
        Quản lý cửa hàng A và Nhân viên cửa hàng B vẫn gửi được yêu cầu đổi thông tin
        \textbf{của cửa hàng B}. Anh muốn siết theo đúng cửa hàng như Công nợ không?
  \item \textbf{Nhân viên bị chặn bằng trang lỗi kỹ thuật 403}, không phải trang thông báo.
        Cần làm giống Công nợ.
  \item \textbf{Loại ``Khác'' bắt người dùng gõ tên kỹ thuật của trường} --- không ai ngoài dev
        biết gõ gì. Anh muốn đổi thành ô mô tả tự do không?
  \item \textbf{Duyệt yêu cầu KHÔNG tự cập nhật hồ sơ cửa hàng.} Bấm duyệt chỉ đổi trạng thái;
        nhân sự Ngô Gia vẫn phải tự vào sửa tay. Anh xác nhận giữ vậy, hay muốn duyệt là áp
        luôn?
  \item \textbf{Phiếu nháp ở đây gửi được} (khác Đổi trả) --- nhưng chỉ khi tạo lại từ đầu; màn
        chi tiết cũng không có nút gửi. Cần rà lại cho thống nhất.
\end{enumerate}
\end{ask}

% =============================================================================
\section{Màn Chi tiết yêu cầu cập nhật thông tin}
\rt{/portal/info-request/<int:req_id>}
\medskip

\mvao

Bấm một dòng ở màn danh sách, hoặc sau khi tạo xong.

\mai

Yêu cầu phải thuộc \textbf{một trong các cửa hàng của tôi} --- \textbf{không} đòi phải là người
tạo. Sai thì về danh sách. Không chặn vai trò.

\mhien

\begin{hienbang}
Đầu yêu cầu &
  \rt{wujia.info.update.request} &
  Mã, cửa hàng, loại thông tin, trạng thái, mức độ, ngày tạo, ngày gửi &
  --- & --- \\ \addlinespace
Giá trị cũ &
  \rt{wujia.franchise.management} &
  \textbf{Đọc lại từ cửa hàng ngay lúc mở màn} --- không phải ảnh chụp lúc gửi. Ngô Gia sửa hồ
  sơ xong thì ô ``giá trị cũ'' của yêu cầu cũ \textbf{đổi theo}, và trông như cũ đã bằng mới. &
  --- & --- \\ \addlinespace
Giá trị mới &
  \rt{new_value} &
  Chữ người dùng đã gõ &
  --- & --- \\ \addlinespace
Tệp đính kèm · lý do từ chối &
  \rt{ir.attachment} · \rt{refuse_reason} &
  Tệp minh chứng và lý do khi bị từ chối &
  --- & --- \\
\end{hienbang}

\mloc

Không có ô lọc.

\mkiem{Huỷ yêu cầu}

Xem mục kế tiếp.

\mghi

Mở màn \textbf{không ghi gì}.

\mhoi

\begin{ask}
\textbf{``Giá trị cũ'' không phải ảnh chụp.} Mô tả trong mã nguồn ghi là ảnh chụp lúc mở form,
nhưng thực tế nó đọc lại mỗi lần mở màn. Sau khi Ngô Gia áp thay đổi, hồ sơ lịch sử không còn
cho biết trước đó là gì. Anh có cần lưu lại giá trị cũ thật không?
\end{ask}

% =============================================================================
\section{Huỷ yêu cầu cập nhật thông tin}
\rt{/portal/info-request/<int:req_id>/cancel}
\medskip

\mvao

Nút \emph{Huỷ} ở màn chi tiết.

\mai

Yêu cầu phải thuộc cửa hàng của tôi \textbf{và do chính tôi tạo}. Người khác cùng cửa hàng xem
được nhưng không huỷ được.

\mhien

Không hiện gì --- quay về màn chi tiết kèm băng thông báo.

\mloc

Không có ô nhập.

\mkiem{Huỷ}

\begin{kiembang}
1 & Yêu cầu thuộc cửa hàng của tôi và do tôi tạo & Đưa về danh sách \\
2 & Trạng thái đang là \emph{Nháp} hoặc \emph{Đã gửi} &
    Băng báo ``Chỉ có thể huỷ yêu cầu ở trạng thái Nháp/Đã gửi.'' \\
\end{kiembang}

\mghi

Chuyển trạng thái sang \textbf{Từ chối} và ghi lý do \emph{``User huỷ''}.

\mhoi

\begin{ask}
\textbf{Huỷ và bị từ chối trông giống nhau.} Hệ thống không có trạng thái \emph{Đã huỷ} nên
yêu cầu tự huỷ hiện đúng huy hiệu đỏ \emph{Từ chối} như khi Ngô Gia từ chối; chỉ khác dòng lý
do. Anh muốn tách thành trạng thái riêng không?
\end{ask}

% =============================================================================
\section{Lấy giá trị hiện tại của cửa hàng}
\rt{/portal/info-request/franchise/<int:fid>/values}
\medskip

\mvao

Không phải màn --- gọi ngầm từ form tạo yêu cầu mỗi lần người dùng đổi ô \emph{Loại thông tin},
để điền ô \emph{Giá trị hiện tại}.

\mai

Chỉ kiểm \textbf{một điều}: cửa hàng phải nằm trong danh sách tôi truy cập được.
\textbf{Không kiểm vai trò} --- khác với màn tạo yêu cầu vốn chỉ cho Chủ tiệm / Quản lý.

\mhien

Trả về đúng một giá trị: nội dung của trường tương ứng trên hồ sơ cửa hàng. Bốn loại thông tin
có sẵn ánh xạ (Địa chỉ, Số điện thoại, Email, Tên chủ cửa hàng); loại \emph{Khác} thì lấy theo
\textbf{tên trường người gọi truyền lên}.

\mloc

Hai tham số: loại thông tin và tên trường.

\mkiem{gọi}

\begin{kiembang}
1 & Cửa hàng nằm trong danh sách của tôi & Trả mã \rt{forbidden} \\
2 & Cửa hàng còn tồn tại & Trả mã \rt{not_found} \\
3 & Tên trường có thật trên hồ sơ cửa hàng & Trả chuỗi rỗng \\
\end{kiembang}

\mghi

\textbf{Không ghi gì.}

\mhoi

\begin{ask}
\textbf{Đường dẫn này đọc được BẤT KỲ trường nào của hồ sơ cửa hàng}, chỉ cần biết tên kỹ thuật
--- vì loại \emph{Khác} nhận thẳng tên trường do người gọi truyền lên và không có danh sách
trường cho phép. Nó lại \textbf{không kiểm vai trò}, nên một Nhân viên đọc được cả những số mà
màn Công nợ đã chặn họ (ví dụ số dư công nợ lưu trên hồ sơ cửa hàng). Dev sẽ đưa vào đợt chuẩn
hoá F1 --- ghi ở đây để BA nắm, \textbf{không cần lên task riêng}.
\end{ask}

% =============================================================================
\section{Màn Danh sách bài viết kiến thức}
\rt{/portal/knowledge}
\medskip

\mvao

Mục \emph{Kiến thức} trên thanh điều hướng. Tham số: \rt{page}, \rt{category_id},
\rt{tag_id}, \rt{keyword}, \rt{page_size}, \rt{notice}.

\mai

\begin{itemize}[leftmargin=*]
  \item Chỉ cần \textbf{đăng nhập}. Không cần chọn cửa hàng, \textbf{không} chặn vai trò,
        \textbf{không} lọc theo cửa hàng.
  \item Mọi tài khoản portal thấy \textbf{đúng cùng một tập bài viết} --- phân hệ duy nhất
        trong portal như vậy.
\end{itemize}

\mhien

\begin{hienbang}
Lưới bài viết &
  \rt{wujia.knowledge.article} &
  Bài \textbf{đang hiện trên portal} (đã phát hành, chưa lưu trữ, chưa hết hạn hiệu lực)
  \textbf{và} đã tới ngày phát hành (\rt{publish_date} rỗng hoặc đã qua) &
  Số thứ tự, rồi ngày phát hành mới nhất; \textbf{12} bài mỗi trang (12/24/48) &
  \uat{21}/27 bài đang hiện \\ \addlinespace
Trạng thái bài không hiện &
  \rt{state} &
  \uat{3} bài còn nháp · \uat{1} bài đã lưu trữ · \uat{2} bài \textbf{đã phát hành nhưng quá
  hạn hiệu lực} &
  --- &
  \uat{23} bài ở trạng thái đã phát hành \\ \addlinespace
Cột danh mục &
  \rt{wujia.knowledge.category} &
  Danh mục đang bật &
  Số thứ tự, rồi tên &
  \uat{9} danh mục --- trong đó có \textbf{tên trùng nhau} \\ \addlinespace
Khối thẻ &
  \rt{wujia.knowledge.tag} &
  Thẻ đang bật &
  Theo tên &
  \uat{6} thẻ \\ \addlinespace
Lượt xem &
  \rt{view_count} &
  Đếm mỗi lần mở màn chi tiết &
  --- &
  Bài cao nhất \uat{65} lượt \\
\end{hienbang}

\mloc

\begin{itemize}[leftmargin=*]
  \item \textbf{Từ khoá} --- khớp một phần theo \textbf{tiêu đề}, \textbf{tóm tắt} hoặc
        \textbf{toàn văn nội dung}. Đây là ô tìm sâu nhất trong portal.
  \item \textbf{Danh mục} và \textbf{Thẻ} --- chọn một cái mỗi loại. Nếu danh mục hay thẻ vừa bị
        Ngô Gia lưu trữ (hoặc số gửi lên không đọc được) thì màn \textbf{quay về danh sách gốc}
        kèm băng báo ``danh mục/thẻ không còn'' --- cố ý không im lặng bỏ lọc.
  \item \textbf{Số bài mỗi trang} --- \textbf{12, 24, 48} (bội của 3 cho lưới 3 cột), khác bậc
        10/20/50 của các màn khác.
\end{itemize}

\mkiem{một bài viết}

Không kiểm gì ở màn danh sách.

\mghi

\textbf{Không ghi gì.}

\mhoi

\begin{ask}
\begin{enumerate}[leftmargin=*]
  \item \textbf{Kiến thức không phân biệt cửa hàng và vai trò.} Mọi nhân viên của mọi cửa hàng
        đọc được mọi bài. Nếu sau này có tài liệu chỉ dành cho Chủ tiệm thì phải làm thêm.
  \item \textbf{Danh mục đang trùng tên trên UAT} --- \emph{Vận hành} và \emph{Sản phẩm} mỗi cái
        xuất hiện hai lần (hai bản ghi khác nhau). Thanh lọc bên trái vì thế hiện hai dòng
        giống hệt. Cần dọn dữ liệu.
  \item \textbf{Bài hết hạn không biến mất ngay.} Cờ ``đang hiện'' được lưu sẵn và chỉ tính lại
        bằng tác vụ nền chạy hằng ngày, nên bài đặt hạn 10 giờ sáng vẫn đọc được tới lần chạy
        kế tiếp. Anh chấp nhận độ trễ này chứ?
  \item \textbf{Lượt xem đếm mỗi lần mở}, không chống bấm lặp --- một người bấm lại 10 lần thì
        cộng 10. Con số \uat{65} lượt trên UAT nên hiểu theo nghĩa đó.
\end{enumerate}
\end{ask}

% =============================================================================
\section{Màn Chi tiết bài viết}
\rt{/portal/knowledge/<string:slug>}
\medskip

\mvao

Bấm một thẻ bài ở màn danh sách, hoặc từ khối \emph{Kiến thức} trên Trang chủ, hoặc từ ô gợi ý
tìm kiếm. Đường dẫn dùng \textbf{tên rút gọn} của bài chứ không dùng số --- ví dụ
\rt{/portal/knowledge/checklist-mo-cua-hang-buoi-sang}.

\mai

Bài phải \textbf{đang hiện trên portal} và đã tới ngày phát hành. Không thoả --- kể cả khi tên
rút gọn gõ sai --- thì về danh sách kèm băng ``bài viết không còn''.

\mhien

\begin{hienbang}
Nội dung bài &
  \rt{wujia.knowledge.article} &
  Tiêu đề, tóm tắt, ngày phát hành, danh mục, thẻ, toàn văn nội dung &
  --- & --- \\ \addlinespace
Tệp đính kèm &
  \rt{ir.attachment} &
  Tệp gắn vào bài, tải qua đường dẫn có kiểm quyền &
  --- & --- \\ \addlinespace
Bài liên quan &
  \rt{wujia.knowledge.article} &
  Bài \textbf{cùng danh mục}, đang hiện, trừ chính bài đang đọc &
  Ngày phát hành mới nhất; \textbf{tối đa 4} bài &
  --- \\
\end{hienbang}

\mloc

Không có ô lọc.

\mkiem{một bài liên quan}

Không kiểm gì --- mở sang màn chi tiết của bài đó, và màn đó kiểm lại từ đầu.

\mghi

\textbf{Có ghi} --- đây là màn chỉ-đọc duy nhất trong portal có ghi dữ liệu: mỗi lần mở,
\textbf{lượt xem của bài cộng thêm 1}. Phép cộng chạy thẳng dưới cơ sở dữ liệu nên hai người mở
cùng lúc không ghi đè nhau. Không ghi ai đã đọc bài --- \textbf{không có} khái niệm ``đã đọc''
cho Kiến thức.

\mhoi

\begin{ask}
\begin{enumerate}[leftmargin=*]
  \item \textbf{Không biết ai đã đọc bài nào.} Nếu Ngô Gia cần chắc rằng cửa hàng đã đọc tài
        liệu bắt buộc thì phải làm thêm phần xác nhận đã đọc.
  \item \textbf{Đường dẫn dùng tên rút gọn.} Đổi tiêu đề bài \textbf{không} đổi tên rút gọn, nên
        liên kết cũ vẫn sống --- nhưng tên rút gọn sẽ không còn khớp tiêu đề.
\end{enumerate}
\end{ask}

% =============================================================================
\section{Tải tệp đính kèm của bài viết}
\rt{/portal/knowledge/<string:slug>/attachment/<int:att_id>}
\medskip

\mvao

Bấm tên tệp ở màn chi tiết bài.

\mai

\begin{kiembang}
1 & Bài đang hiện trên portal & Về danh sách kèm băng ``bài viết không còn'' \\
2 & Tệp thuộc đúng bài đó & Từ chối truy cập \\
\end{kiembang}

Ở đây \textbf{cố ý không} kiểm theo cửa hàng --- kiến thức là tài liệu chung cho mọi tài khoản
portal.

\mhien

Trả về tệp ở chế độ \textbf{tải xuống}.

\mloc

Không có.

\mkiem{tên tệp}

Xem bảng trên.

\mghi

\textbf{Không ghi gì} --- tải tệp không cộng lượt xem.

\mhoi

\begin{ask}
Không có điểm cần chốt.
\end{ask}

% =============================================================================
\section{Gợi ý tìm kiếm kiến thức}
\rt{/portal/knowledge/search}
\medskip

\mvao

Không phải màn --- gọi ngầm khi người dùng đang gõ trong ô tìm kiếm, sau khi ngừng gõ khoảng
0,3 giây.

\mai

Chỉ cần đăng nhập. Cùng phạm vi với màn danh sách: mọi bài đang hiện.

\mhien

Trả về danh sách gợi ý, mỗi dòng gồm tiêu đề, danh mục và đường dẫn tới bài. Tìm theo tiêu đề,
tóm tắt và toàn văn nội dung.

\mloc

\begin{itemize}[leftmargin=*]
  \item \textbf{Từ khoá} --- \textbf{dưới 2 ký tự thì không gọi}, trả về rỗng.
  \item \textbf{Số gợi ý} --- mặc định 10, tối đa \textbf{20}; số lạ tự về 10.
\end{itemize}

\mkiem{gõ phím}

Không kiểm gì ngoài độ dài từ khoá.

\mghi

\textbf{Không ghi gì} --- gợi ý không cộng lượt xem.

\mhoi

\begin{ask}
\textbf{Chưa có giới hạn số lần gọi.} Gõ nhanh và liên tục sẽ tạo nhiều truy vấn tìm toàn văn;
với 1.500 tài khoản cần theo dõi khi chạy thật.
\end{ask}

% =============================================================================
\section{Màn Danh sách thông báo}
\rt{/portal/notification} · \rt{/portal/notification/results}
\medskip

\mvao

Mục \emph{Thông báo} trên thanh điều hướng, hoặc dòng \emph{Xem tất cả} trong hộp chuông.
Đường dẫn \rt{/results} không phải màn riêng --- là phần kết quả gọi ngầm khi đổi bộ lọc.
Tham số: \rt{page}, \rt{limit}, \rt{type_id}, \rt{keyword}, \rt{priority},
\rt{read_status} (hoặc \rt{unread} kiểu cũ), \rt{date_from}, \rt{date_to}.

\mai

\begin{itemize}[leftmargin=*]
  \item Phải đăng nhập. \textbf{Không chặn vai trò.}
  \item Thấy thông báo \textbf{gửi toàn hệ} \emph{hoặc} \textbf{gửi đích danh một trong các cửa
        hàng của tôi}. Chưa chọn cửa hàng thì vẫn thấy phần gửi toàn hệ.
  \item Danh sách là \textbf{lịch sử}: gồm cả thông báo \textbf{đã hết hiệu lực} (hộp chuông thì
        không).
  \item Thực tế UAT: \uat{3} thông báo, \uat{2} đang phát hành --- cả hai đều gửi toàn hệ;
        \uat{0} thông báo gửi đích danh cửa hàng đang hiện ra portal. Nghĩa là luật lọc theo
        cửa hàng \textbf{chưa được dữ liệu UAT kiểm chứng}.
\end{itemize}

\mhien

\begin{hienbang}
Danh sách thông báo &
  \rt{wujia.notification} &
  Đã phát hành ra portal (\rt{is_published_portal = true}), \textbf{đã tới giờ phát hành},
  và đúng đối tượng nhận &
  \textbf{Ghim lên đầu}, rồi ngày phát hành mới nhất; \textbf{10} thông báo mỗi trang
  (10/20/50) &
  \uat{2} đang phát hành · \uat{1} chưa phát hành · \uat{0} được ghim \\ \addlinespace
Huy hiệu loại &
  \rt{wujia.notification.type} &
  Năm loại kèm màu riêng: Khẩn cấp · Thông báo chung · Khuyến mãi · Hệ thống · Khác &
  Theo số thứ tự &
  \uat{5} loại đang bật \\ \addlinespace
Huy hiệu mức độ &
  \rt{priority} &
  Ba nhãn: \emph{Thông thường} · \emph{Quan trọng} · \emph{Cần làm} &
  --- & --- \\ \addlinespace
Dấu đã đọc &
  \rt{wujia.notification.read} &
  Đánh dấu theo bộ ba \textbf{người + thông báo + cửa hàng đang chọn}. Chưa chọn cửa hàng thì
  bỏ điều kiện cửa hàng khi tra. &
  --- &
  \uat{6} dấu cũ chưa gắn cửa hàng \\ \addlinespace
Huy hiệu hết hiệu lực &
  \rt{expired_date} &
  Thông báo quá hạn vẫn nằm trong danh sách nhưng được đánh dấu riêng &
  --- &
  \uat{0} thông báo đã hết hiệu lực \\
\end{hienbang}

\mloc

\begin{itemize}[leftmargin=*]
  \item \textbf{Từ khoá} --- tiêu đề, \textbf{số công văn}, hoặc tóm tắt.
  \item \textbf{Loại} và \textbf{Mức độ} --- mỗi thứ chọn một; giá trị lạ bị bỏ qua.
  \item \textbf{Trạng thái đọc} --- \emph{Tất cả} / \emph{Đã đọc} / \emph{Chưa đọc}. Chọn
        \emph{Chưa đọc} thì hệ thống \textbf{tự loại luôn thông báo đã hết hiệu lực} --- quá hạn
        không bị tính là còn nợ đọc.
  \item \textbf{Khoảng ngày} --- lọc theo \emph{ngày phát hành}, quy đổi đúng múi giờ tài khoản.
        Ngày ngược thì \textbf{chặn truy vấn}, giữ chữ đã gõ và báo tại thanh lọc --- cố ý không
        bỏ lọc rồi trả về toàn bộ.
  \item \textbf{Số dòng mỗi trang} --- 10, 20, 50.
\end{itemize}

\mkiem{một thông báo}

Không kiểm gì ở danh sách.

\mghi

\textbf{Không ghi gì} --- xem danh sách không đánh dấu đã đọc.

\mhoi

\begin{ask}
\begin{enumerate}[leftmargin=*]
  \item \textbf{``Đã đọc'' tính theo từng cửa hàng.} Người phụ trách 3 cửa hàng phải đọc cùng
        một thông báo toàn hệ \textbf{ba lần} mới sạch chuông. Anh muốn giữ vậy, hay đọc một
        lần là xong cho mọi cửa hàng?
  \item \uat{6} dấu đã đọc cũ \textbf{chưa gắn cửa hàng} nên không khớp cửa hàng nào --- chúng
        chỉ có tác dụng khi người dùng \textbf{chưa chọn} cửa hàng. Cần dọn dữ liệu.
  \item \textbf{Thẻ đếm chưa đọc trên Trang chủ tính khác màn này} --- Trang chủ bỏ qua điều
        kiện hết hiệu lực và điều kiện đã tới giờ phát hành, lại trừ \textbf{mọi} dấu đã đọc.
        Đây là chênh lệch đã nêu ở chương 4; nhắc lại để BA gom vào một quyết định.
  \item \textbf{UAT chưa có thông báo gửi đích danh cửa hàng nào đang phát hành} --- cần dựng ca
        mẫu trước khi nghiệm thu phần lọc theo cửa hàng.
\end{enumerate}
\end{ask}

% =============================================================================
\section{Màn Chi tiết thông báo}
\rt{/portal/notification/<int:notification_id>}
\medskip

\mvao

Bấm một dòng ở danh sách hoặc một dòng trong hộp chuông.

\mai

Thông báo phải \textbf{đã phát hành, đã tới giờ, đúng đối tượng nhận}. Thông báo
\textbf{đã hết hiệu lực vẫn mở được} --- cố ý, để tra lại lịch sử. Không thoả thì về danh sách.

\mhien

\begin{hienbang}
Nội dung &
  \rt{wujia.notification} &
  Tiêu đề, số công văn, loại, mức độ, ngày phát hành, toàn văn nội dung &
  --- & --- \\ \addlinespace
Tệp đính kèm &
  \rt{ir.attachment} &
  Tệp gắn vào thông báo, tải qua đường dẫn có kiểm quyền &
  --- &
  \uat{0}/3 thông báo có tệp \\ \addlinespace
Ghi chú nội bộ &
  --- &
  \textbf{Không hiện} --- ô ghi chú của Ngô Gia không ra portal &
  --- & --- \\
\end{hienbang}

\mloc

Không có ô lọc.

\mkiem{mở màn}

\begin{kiembang}
1 & Thông báo đã phát hành, đã tới giờ, đúng đối tượng & Về danh sách \\
2 & Đã chọn được cửa hàng &
    \textbf{Vẫn cho đọc nội dung}, chỉ \textbf{không ghi} dấu đã đọc --- tránh tạo dấu không
    gắn cửa hàng nào \\
\end{kiembang}

\mghi

\begin{itemize}[leftmargin=*]
  \item Lần đầu mở tại một cửa hàng: tạo một dấu đã đọc cho bộ ba \textbf{người + thông báo +
        cửa hàng}, ghi cả \emph{lần đọc đầu} và \emph{lần mở gần nhất}.
  \item Mở lại: \textbf{chỉ cập nhật lần mở gần nhất}, giữ nguyên lần đọc đầu.
  \item Chưa chọn cửa hàng: \textbf{không ghi gì}, và chuông vẫn coi là chưa đọc.
\end{itemize}

\mhoi

\begin{ask}
\textbf{Mở thông báo khi chưa chọn cửa hàng thì không được tính là đã đọc} --- người dùng đọc
xong vẫn thấy chuông báo đỏ. Cần nói rõ trên màn hay tự chọn cửa hàng giúp họ?
\end{ask}

% =============================================================================
\section{Hộp chuông trên đầu trang}
\rt{/portal/notification/recent} · \rt{/portal/notification/unread-count}
\medskip

\mvao

Không phải màn --- hai đường dẫn phục vụ biểu tượng chuông có mặt trên \textbf{mọi trang}
portal: một cái nạp nội dung khi bấm mở, một cái chỉ trả về con số trên chuông.

\mai

Phải đăng nhập. Phạm vi giống màn danh sách, nhưng \textbf{chỉ thông báo còn hiệu lực} --- quá
hạn thì rụng khỏi chuông dù vẫn nằm trong lịch sử.

\mhien

\begin{hienbang}
Danh sách trong hộp &
  \rt{wujia.notification} &
  Thông báo \textbf{còn hiệu lực} đúng đối tượng nhận &
  Ghim lên đầu, rồi mới nhất; \textbf{đúng 5} dòng &
  \uat{2} thông báo còn hiệu lực trên UAT \\ \addlinespace
Số trên chuông &
  \rt{wujia.notification.read} &
  \textbf{Số thông báo còn hiệu lực trừ đi số dấu đã đọc} của tôi tại cửa hàng đang chọn; sàn ở
  0 &
  --- &
  Toàn hệ có \uat{18} dấu đã đọc, trong đó \uat{6} dấu chưa gắn cửa hàng \\
\end{hienbang}

Hai đường dẫn này chạy trên mọi trang nên được viết để \textbf{chỉ đếm, không kéo dữ liệu về} ---
điều kiện ``còn hiệu lực'' được đẩy xuống thành câu hỏi phụ trong cơ sở dữ liệu thay vì nạp
danh sách mã ra rồi so ở tầng ứng dụng.

\mloc

Không có ô nhập.

\mkiem{mở chuông}

Không kiểm gì ngoài đăng nhập.

\mghi

\textbf{Không ghi gì} --- mở chuông không đánh dấu đã đọc.

\mhoi

\begin{ask}
\textbf{Hộp chuông cố định 5 dòng} và không có nút ``xem thêm'' trong hộp --- phải sang màn danh
sách. Anh xác nhận.
\end{ask}

% =============================================================================
\section{Đánh dấu đã đọc}
\rt{/portal/notification/mark-all-read} · \rt{/portal/notification/mark-read}
\medskip

\mvao

Nút \emph{Đánh dấu tất cả đã đọc} trong hộp chuông (đường dẫn thứ nhất); đường dẫn thứ hai giữ
lại cho các màn cũ gửi lên một danh sách mã cụ thể.

\mai

Phải đăng nhập \textbf{và phải đang chọn một cửa hàng}. Chưa chọn thì cả hai trả về lỗi
\emph{``Vui lòng chọn cửa hàng trước khi thao tác''} và \textbf{không ghi gì}.

\mhien

Không hiện gì --- trả về số dấu vừa tạo và số chưa đọc còn lại.

\mloc

Đường dẫn thứ hai nhận một danh sách mã thông báo; mã không đọc được thì trả lỗi
\rt{invalid_ids}.

\mkiem{Đánh dấu tất cả đã đọc}

\begin{kiembang}
1 & Đang chọn một cửa hàng & ``Vui lòng chọn cửa hàng trước khi thao tác.'' \\
2 & (đường dẫn thứ hai) Mã gửi lên đều là số & \rt{invalid_ids} \\
3 & (đường dẫn thứ hai) Thông báo \textbf{đúng đối tượng nhận} &
    Mã lạ bị \textbf{lặng lẽ bỏ qua} --- không báo lỗi, không đánh dấu \\
\end{kiembang}

\mghi

\begin{itemize}[leftmargin=*]
  \item \textbf{Đánh dấu tất cả}: tạo dấu đã đọc cho \textbf{mọi thông báo còn hiệu lực} chưa
        có dấu, tại \textbf{cửa hàng đang chọn}. \textbf{Không} ghi mốc ``lần mở gần nhất'' ---
        người dùng chưa thực sự mở nội dung.
  \item \textbf{Đánh dấu theo danh sách}: cùng cách, nhưng có ghi cả mốc lần mở.
  \item Cả hai đều \textbf{chạy lại được nhiều lần không sinh trùng} --- bộ ba người + thông báo
        + cửa hàng là duy nhất.
  \item \textbf{Chỉ đánh dấu cho cửa hàng đang chọn}: người ba cửa hàng bấm \emph{tất cả} vẫn
        còn chuông đỏ ở hai cửa hàng kia.
\end{itemize}

\mhoi

\begin{ask}
\textbf{``Đánh dấu tất cả'' không phải là tất cả} --- chỉ là tất cả \emph{tại cửa hàng đang
chọn} và chỉ những thông báo \emph{còn hiệu lực}. Nhãn nút nên nói rõ điều đó.
\end{ask}

% =============================================================================
\section{Tải tệp đính kèm của thông báo}
\rt{/portal/notification/<int:notification_id>/attachment/<int:attachment_id>}
\medskip

\mvao

Bấm tên tệp ở màn chi tiết thông báo.

\mai

\begin{kiembang}
1 & Thông báo đã phát hành, đã tới giờ, \textbf{đúng đối tượng nhận của tôi} & Báo không tìm thấy \\
2 & Tệp thuộc đúng thông báo đó &
    Từ chối truy cập --- đây chính là lý do portal không dùng đường dẫn tải tệp mặc định của
    Odoo, vì đường dẫn đó chỉ cần biết số tệp là tải được \\
\end{kiembang}

\mhien

Trả về tệp ở chế độ \textbf{tải xuống}.

\mloc

Không có.

\mkiem{tên tệp}

Xem bảng trên.

\mghi

\textbf{Không ghi gì.}

\mhoi

\begin{ask}
Không có điểm cần chốt.
\end{ask}

\chapter{Đăng ký thi}

Chương này gồm \textbf{7 đường dẫn} của mô-đun \rt{wujia_portal_exam}, gom thành \textbf{7 mục}
(cộng một mục cho đường dẫn cũ). Một đường dẫn --- \rt{/portal/exam/register} --- \textbf{có hai
hàm}: một hàm mở màn, một hàm nhận phiếu gửi lên; chương này tách thành hai mục vì hai việc khác
hẳn nhau.

\textbf{Bốn điểm cần biết trước khi đọc:}

\begin{enumerate}[leftmargin=*]
  \item \textbf{Trên UAT hiện KHÔNG ai đăng ký thi được.} Có \uat{3} khóa thi đã phát hành và
        \uat{2} kỳ thi còn ở trạng thái mở, nhưng ngày thi của chúng là \textbf{16/08} và
        \textbf{25/08} --- đã qua --- và hạn đăng ký là 09:00 chính ngày thi. Vì vậy cả ba khóa
        đều hiện chữ \emph{``Đã đóng''} và lịch \textbf{không có ngày nào bấm được}. Giống hệt
        tình trạng của Đổi trả ở chương 6, cần dựng dữ liệu mẫu trước khi nghiệm thu.
  \item \textbf{Không chặn vai trò.} Nhân viên đăng ký thi được như Chủ tiệm. Đây là phân hệ
        ghi dữ liệu duy nhất ngoài Đặt hàng không có cổng vai trò.
  \item \textbf{Phạm vi là CỬA HÀNG ĐANG CHỌN, không phải mọi cửa hàng.} Người phụ trách nhiều
        cửa hàng phải bấm đổi cửa hàng mới xem được phiếu của cửa hàng kia --- giống Công nợ,
        khác Đổi trả và Yêu cầu thông tin.
  \item \textbf{Gửi phiếu xong là chốt.} Portal không có nút sửa, không có nút huỷ, không gửi
        lại được phiếu bị từ chối. Mọi thay đổi phải nhờ Ngô Gia làm trong ERP.
\end{enumerate}

Số thật trên UAT: \uat{1} phiếu đăng ký (của cửa hàng \emph{TP HCM Quận 1}, do Administrator
tạo, đã được xác nhận, \uat{1} người dự thi, kết quả \emph{Đạt}), \uat{3} kỳ thi, \uat{3} ca giờ,
\uat{0} phiếu chờ duyệt, \uat{0} phiếu bị từ chối, \uat{0} phiếu bị huỷ.

% =============================================================================
\section{Màn Danh sách phiếu đăng ký thi}
\rt{/portal/exam}
\medskip

\mvao

Mục \emph{Đăng ký thi} trên thanh điều hướng. Tham số: \rt{page}, \rt{limit}, \rt{state},
\rt{result}, \rt{q}, \rt{date_from}, \rt{date_to}.

\mai

\begin{itemize}[leftmargin=*]
  \item Phải đăng nhập. \textbf{Không chặn vai trò} --- mọi thành viên cửa hàng đều xem được.
  \item Chỉ thấy phiếu của \textbf{cửa hàng đang chọn}. \textbf{Chưa chọn cửa hàng} thì màn
        hiện đúng câu \emph{``Chưa có đăng ký thi''} --- \textbf{không} báo là do chưa chọn
        cửa hàng, nên người dùng dễ tưởng cửa hàng mình chưa từng đăng ký.
  \item Thực tế UAT: \uat{1} phiếu duy nhất, thuộc cửa hàng \emph{TP HCM Quận 1}. Ba cửa hàng
        còn lại mở ra màn rỗng.
\end{itemize}

\mhien

\begin{hienbang}
Danh sách phiếu &
  \rt{wujia.exam.registration} &
  Phiếu có cửa hàng \textbf{đúng bằng cửa hàng đang chọn} &
  Ngày nộp phiếu mới nhất trước; \textbf{10} phiếu mỗi trang (10/20/50) &
  \uat{1} phiếu toàn hệ \\ \addlinespace
Huy hiệu trạng thái &
  \rt{state} &
  Bốn nhãn: \emph{Chờ xác nhận} · \emph{Đã đăng ký} · \emph{Từ chối} · \emph{Đã hủy}.
  Bản điện thoại gọi nhãn đầu là \emph{``Chờ duyệt''} --- \textbf{cùng một trạng thái, hai chữ
  khác nhau trên hai giao diện}. &
  --- &
  \uat{1} Đã đăng ký · \uat{0} ở ba trạng thái còn lại \\ \addlinespace
Cột kết quả &
  \rt{wujia.exam.session} &
  Chỉ hiện khi \textbf{kỳ thi đã công bố kết quả}; khi đó đếm \emph{n Đạt · n Không đạt} trong
  phiếu. Chưa công bố thì ghi \emph{``Chưa công bố''}. &
  --- &
  \uat{1} kỳ đã công bố · \uat{1} người Đạt · \uat{0} Không đạt \\ \addlinespace
Dòng trên bản điện thoại &
  --- &
  Khi đã công bố, huy hiệu đổi thành \emph{``Có kết quả''} và \textbf{đè lên} trạng thái phiếu
  --- phiếu bị từ chối thuộc kỳ đã công bố cũng hiện \emph{``Có kết quả''} &
  --- & --- \\
\end{hienbang}

\mloc

\begin{itemize}[leftmargin=*]
  \item \textbf{Từ khoá} --- khớp một phần theo \textbf{mã phiếu} hoặc \textbf{tên khóa thi}.
        Không tìm theo tên thí sinh.
  \item \textbf{Trạng thái} --- một trong bốn nhãn; giá trị lạ bị bỏ qua.
  \item \textbf{Kết quả} --- \emph{Đã công bố} / \emph{Chưa công bố}, lọc theo cờ của
        \textbf{kỳ thi} chứ không theo từng người.
  \item \textbf{Khoảng ngày} --- lọc theo \textbf{ngày nộp phiếu}, quy đổi đúng múi giờ tài
        khoản. Ngày ngược thì \textbf{chặn truy vấn}, giữ nguyên hai ô đã gõ và báo tại thanh
        lọc --- giống Thông báo và Giao hàng, khác Công nợ (bên đó lặng lẽ đảo hai ngày).
  \item \textbf{Số dòng mỗi trang} --- 10, 20, 50; gõ số ngoài danh sách thì tự về 10.
\end{itemize}

\mkiem{Đăng ký mới}

Không kiểm gì ở màn danh sách --- mọi phép kiểm nằm ở màn đăng ký.

\mghi

\textbf{Không ghi gì.}

\mhoi

\begin{ask}
\begin{enumerate}[leftmargin=*]
  \item \textbf{Chưa chọn cửa hàng nhưng màn báo ``Chưa có đăng ký thi''} --- sai nguyên nhân.
        Nên đổi thành lời nhắc chọn cửa hàng như các màn khác.
  \item \textbf{Cùng một trạng thái, hai chữ khác nhau giữa máy tính và điện thoại}
        (\emph{Chờ xác nhận} · \emph{Chờ duyệt}). Anh chọn một chữ để thống nhất.
  \item \textbf{Huy hiệu ``Có kết quả'' đè lên trạng thái phiếu} trên bản điện thoại --- phiếu
        bị từ chối hay bị huỷ nhưng thuộc kỳ đã công bố vẫn hiện \emph{Có kết quả}. Cần sửa.
  \item \textbf{Chỉ xem được cửa hàng đang chọn}, trong khi luật dữ liệu của hệ thống cho phép
        xem mọi cửa hàng của mình. Anh muốn gộp như Đổi trả không?
\end{enumerate}
\end{ask}

% =============================================================================
\section{Màn Đăng ký thi --- chọn khóa và lịch}
\rt{/portal/exam/register} \note{(mở màn)}
\medskip

\mvao

Nút \emph{Đăng ký mới} ở màn danh sách. Tham số: \rt{course_id}, \rt{year}, \rt{month}.

\mai

\begin{itemize}[leftmargin=*]
  \item Phải đăng nhập \textbf{và phải đang chọn một cửa hàng} --- chưa chọn thì bị đưa thẳng
        về danh sách, \textbf{không có thông báo}.
  \item Không chặn vai trò.
\end{itemize}

\mhien

\begin{hienbang}
Danh sách khóa thi &
  \rt{wujia.exam.course} &
  Khóa \textbf{đã phát hành} và còn hiệu lực. Khóa còn nháp thì portal không thấy. &
  Theo tên &
  \uat{3}/3 khóa đã phát hành \\ \addlinespace
Nhãn trạng thái khóa &
  \rt{wujia.exam.session} &
  \emph{``Còn lịch''} khi còn ít nhất một kỳ đăng ký được; \emph{``Hết chỗ''} khi còn kỳ trong
  hạn nhưng hết ghế; \emph{``Đã đóng''} khi không còn kỳ nào &
  --- &
  Cả \uat{3} khóa đang hiện \textbf{Đã đóng} \\ \addlinespace
Dòng mô tả khóa &
  \rt{registration_horizon_days} &
  \emph{``n kỳ thi $\bullet$ Trong 60 ngày tới''} --- chỉ đếm kỳ \textbf{đang mở} có ngày thi
  \textbf{từ hôm nay} tới hết cửa sổ &
  --- &
  Cả 3 khóa đặt cửa sổ \uat{60} ngày; hiện đếm ra \uat{0} kỳ \\ \addlinespace
Lịch tháng &
  \rt{wujia.exam.session} &
  Kỳ thi của khóa đang chọn, trong tháng đang xem, \textbf{chưa bị huỷ}. Mỗi ngày mang một
  trong ba trạng thái: \emph{còn chỗ} · \emph{hết chỗ} · \emph{không có lịch}. &
  Tuần bắt đầu từ Thứ 2 &
  \uat{0} ngày bấm được \\ \addlinespace
Bảng tóm tắt (máy tính) &
  --- &
  Tên khóa, tên cửa hàng đang chọn, \emph{``0 / n''} số người đã thêm trên giới hạn mỗi phiếu.
  Ngày thi · địa điểm · hạn đăng ký để trống, do chọn khung giờ mới có. &
  --- & --- \\
\end{hienbang}

Một ngày được coi là \textbf{còn chỗ} khi kỳ thi hôm đó \textbf{đang mở}, \textbf{chưa quá hạn
đăng ký}, \textbf{còn ghế trống}, và ngày đó nằm trong khoảng \textbf{từ hôm nay tới hết cửa sổ
60 ngày}. Thiếu một trong bốn điều kiện thì ngày đó không bấm được.

\mloc

Chọn \textbf{khóa thi} và bấm mũi tên đổi \textbf{tháng}. Chọn khóa không có trong danh sách
đã phát hành thì màn \textbf{tự quay về khóa đầu tiên}. Gõ tháng lạ trên thanh địa chỉ có thể
làm màn báo lỗi kỹ thuật --- số tháng chưa được kiểm.

\mkiem{một ngày trên lịch}

Bấm ngày chỉ mở bảng khung giờ; mọi phép kiểm thật nằm ở bước gửi phiếu.

\mghi

\textbf{Không ghi gì.}

\mhoi

\begin{ask}
\begin{enumerate}[leftmargin=*]
  \item \textbf{UAT không có kỳ thi nào đăng ký được} --- cần tạo ít nhất một kỳ ngày tương lai
        để nghiệm thu được toàn bộ luồng.
  \item \textbf{Chưa chọn cửa hàng thì bị đá về danh sách không lời giải thích.}
  \item \textbf{Số tháng gõ trên thanh địa chỉ chưa được kiểm} --- dev sẽ đưa vào đợt chuẩn hoá,
        không cần task riêng.
  \item \textbf{``Cửa sổ 60 ngày'' đang đặt ở khóa thi}, mỗi khóa một giá trị. Anh xác nhận để
        Ngô Gia tự chỉnh theo từng khóa, hay muốn chốt chung một con số?
\end{enumerate}
\end{ask}

% =============================================================================
\section{Lấy lịch tháng của một khóa thi}
\rt{/portal/exam/calendar}
\medskip

\mvao

Không phải màn --- gọi ngầm khi người dùng đổi khóa thi hoặc bấm mũi tên sang tháng khác.

\mai

Phải đăng nhập \textbf{và đang chọn cửa hàng}; chưa chọn thì trả về lời nhắc \emph{``Vui lòng
chọn cửa hàng trước khi thao tác''}. Không chặn vai trò.

\mhien

Trả về ma trận tuần của tháng, mỗi ô một ngày kèm trạng thái, đúng như mô tả ở màn trên.

\mloc

Ba tham số: khóa thi, năm, tháng. Thiếu năm/tháng thì lấy tháng hiện tại.

\mkiem{gọi}

\begin{kiembang}
1 & Đang chọn cửa hàng & ``Vui lòng chọn cửa hàng trước khi thao tác.'' \\
2 & Khóa thi tồn tại \textbf{và đang ở trạng thái đã phát hành} & ``Khóa thi không hợp lệ.'' \\
\end{kiembang}

\textbf{Không} kiểm cửa hàng có liên quan gì tới khóa thi --- lịch thi là dữ liệu chung, cửa
hàng nào cũng xem được như nhau.

\mghi

\textbf{Không ghi gì.}

\mhoi

\begin{ask}
Không có điểm cần chốt.
\end{ask}

% =============================================================================
\section{Lấy khung giờ thi của một ngày}
\rt{/portal/exam/slots}
\medskip

\mvao

Không phải màn --- gọi ngầm khi người dùng bấm một ngày còn chỗ trên lịch.

\mai

Giống mục trên: phải đăng nhập và đang chọn cửa hàng.

\mhien

\begin{hienbang}
Danh sách khung giờ &
  \rt{wujia.exam.session} &
  Kỳ thi của khóa đang chọn, \textbf{đúng ngày đó}, chưa bị huỷ &
  Theo giờ bắt đầu &
  \uat{3} ca giờ đang bật: 08--10, 14--16, 16--17 \\ \addlinespace
Nhãn tình trạng &
  --- &
  Bốn khả năng theo thứ tự xét: \emph{Đã đóng} (kỳ không còn mở) $\rightarrow$ \emph{Hết hạn}
  (quá hạn đăng ký) $\rightarrow$ \emph{Hết chỗ} (không còn ghế) $\rightarrow$ \emph{Còn n chỗ} &
  --- & --- \\ \addlinespace
Thông tin kèm theo &
  \rt{wujia.exam.session} &
  Địa điểm, hạn đăng ký, số ghế còn, \textbf{giới hạn người mỗi phiếu} &
  --- &
  Kỳ 25/08 cho \uat{3} người/phiếu trong khi khóa của nó đặt \uat{2} \\
\end{hienbang}

\textbf{Giới hạn người mỗi phiếu} lấy của \textbf{kỳ thi} trước; kỳ để trống mới lấy của khóa.
Đây là nguồn duy nhất dùng cho cả câu hướng dẫn trên màn, phép đếm hạn mức, lẫn phép kiểm lúc
gửi --- nên ba chỗ không bao giờ lệch nhau.

\mloc

Hai tham số: khóa thi và ngày. Ngày nhận cả dạng \emph{2026-08-25} lẫn \emph{25/08/2026}.

\mkiem{gọi}

\begin{kiembang}
1 & Đang chọn cửa hàng & ``Vui lòng chọn cửa hàng trước khi thao tác.'' \\
2 & Khóa thi tồn tại và đã phát hành, ngày đọc được & ``Lịch thi không hợp lệ.'' \\
\end{kiembang}

\textbf{Không} kiểm ngày có nằm trong cửa sổ 60 ngày hay không --- màn lịch đã chặn từ trước
bằng cách không cho bấm những ngày đó.

\mghi

\textbf{Không ghi gì.}

\mhoi

\begin{ask}
\textbf{Kỳ thi đặt giới hạn người riêng, khác giới hạn của khóa.} Trên UAT kỳ 25/08 cho 3 người
mỗi phiếu còn khóa của nó ghi 2. Anh xác nhận cách này (kỳ đè khóa) là đúng ý.
\end{ask}

% =============================================================================
\section{Gửi phiếu đăng ký thi}
\rt{/portal/exam/register} \note{(gửi phiếu)}
\medskip

\mvao

Không phải màn --- nút \emph{Xác nhận đăng ký} ở bước cuối của màn đăng ký.

\mai

Phải đăng nhập và đang chọn cửa hàng. Không chặn vai trò. Cửa hàng ghi vào phiếu là
\textbf{cửa hàng đang chọn}, không phải số mà màn gửi lên --- nên không sửa được bằng cách can
thiệp dữ liệu gửi đi.

\mhien

Gửi xong thì chuyển thẳng sang màn chi tiết của phiếu vừa tạo.

\mloc

\begin{itemize}[leftmargin=*]
  \item \textbf{Khung giờ} --- chọn từ bảng khung giờ.
  \item \textbf{Danh sách người dự thi} --- mỗi người gồm \textbf{họ tên} (bắt buộc),
        \textbf{số điện thoại} (bắt buộc), \textbf{năm sinh}, \textbf{chức vụ}, \textbf{ảnh}.
        Người dự thi \textbf{gõ tay hoàn toàn}, không chọn từ danh sách nhân sự nào.
  \item \textbf{Ghi chú} --- một dòng, không bắt buộc.
  \item \textbf{Ảnh} --- chỉ nhận PNG hoặc JPEG, tối đa \textbf{5 MB} mỗi ảnh; ảnh được
        \textbf{thu nhỏ về tối đa 1920 điểm ảnh} trước khi lưu.
\end{itemize}

\mkiem{Xác nhận đăng ký}

\begin{kiembang}
1 & Đang chọn cửa hàng & ``Vui lòng chọn cửa hàng trước khi thao tác.'' \\
2 & Khung giờ còn tồn tại và khóa của nó đã phát hành &
    ``Khung giờ đã thay đổi hoặc không còn. Vui lòng chọn lại lịch thi.'' \\
3 & Có ít nhất 1 người dự thi & ``Cần ít nhất 1 người dự thi.'' \\
4 & Số người không vượt giới hạn mỗi phiếu & ``Tối đa n người mỗi phiếu.'' \\
5 & Mỗi người đủ họ tên và số điện thoại & ``Mỗi người dự thi cần có họ tên và số điện thoại.'' \\
6 & Số điện thoại đúng dạng (bắt đầu 0 hoặc +84, 9--11 chữ số) &
    ``Số điện thoại '…' không hợp lệ (vd 0901234567).'' \\
7 & Năm sinh là số & ``Năm sinh '…' không hợp lệ.'' \\
8 & Ảnh đúng định dạng, đúng dung lượng, không hỏng &
    ``Ảnh nhân viên không đúng định dạng hoặc vượt dung lượng cho phép.'' /
    ``Ảnh nhân viên không hợp lệ hoặc bị hỏng.'' \\
9 & Kỳ thi \textbf{đang mở đăng ký} & ``Kỳ thi '…' không mở đăng ký.'' \\
10 & \textbf{Chưa quá hạn đăng ký} & ``Kỳ thi '…' đã quá hạn đăng ký.'' \\
11 & \textbf{Còn đủ ghế} sau khi cộng số người của phiếu này &
    ``Kỳ thi '…' đã hết chỗ (sức chứa n, đã giữ m).'' \\
12 & Không có phiếu khác đang được ghi cùng lúc cho kỳ này &
    ``Hệ thống đang xử lý đăng ký khác cho kỳ thi này, vui lòng thử lại.'' \\
\end{kiembang}

Bước 9--12 là phần đáng chú ý nhất: \textbf{hệ thống khoá dòng kỳ thi lại rồi mới đếm ghế},
nên hai cửa hàng bấm gửi cùng lúc không thể cùng chiếm chỗ cuối. Nếu kỳ thi hết chỗ ngay lúc
đó, phiếu \textbf{không được ghi} và người dùng nhận câu báo rõ ràng thay vì lỗi kỹ thuật.

Năm sinh còn bị kiểm thêm ở tầng dữ liệu: phải nằm trong khoảng \textbf{1900 đến năm hiện tại}.

\mghi

\begin{itemize}[leftmargin=*]
  \item Tạo một phiếu \rt{wujia.exam.registration}, mã tự sinh dạng \emph{WJ-EXR/26/00004},
        trạng thái \textbf{Chờ xác nhận} --- \textbf{không có trạng thái nháp}.
  \item Người đăng ký = người đang đăng nhập; cửa hàng = cửa hàng đang chọn; hệ thống còn
        \textbf{tự tìm tư cách thành viên} của người đó tại cửa hàng để gắn vào phiếu.
  \item Tạo một dòng \rt{wujia.exam.registration.line} cho mỗi người dự thi, kết quả để
        \emph{Chưa có}.
  \item \textbf{Ghi ngay số ghế đã giữ} của kỳ thi --- phiếu chờ duyệt đã chiếm chỗ, không đợi
        Ngô Gia xác nhận.
  \item \textbf{Không} gửi email, \textbf{không} tạo việc cần làm. Ngô Gia phải tự vào ERP xem.
\end{itemize}

\mhoi

\begin{ask}
\begin{enumerate}[leftmargin=*]
  \item \textbf{Phiếu chờ duyệt đã chiếm ghế.} Cửa hàng gửi phiếu rồi bỏ đó sẽ khoá chỗ của
        cửa hàng khác cho tới khi Ngô Gia từ chối. Anh muốn giữ vậy, hay chỉ tính ghế khi đã
        xác nhận?
  \item \textbf{Người dự thi gõ tay hoàn toàn}, không nối với danh sách nhân sự của cửa hàng.
        Gõ sai tên thì không có gì đối chiếu. Anh có cần nối vào danh sách thành viên cửa hàng
        không?
  \item \textbf{Gửi phiếu xong không ai được báo} --- giống Hỗ trợ ở chương 7.
  \item \textbf{Không có bước xem lại trước khi gửi ở bản điện thoại} sau khi bấm xác nhận ---
        bấm là gửi luôn.
\end{enumerate}
\end{ask}

% =============================================================================
\section{Màn Chi tiết phiếu đăng ký · kết quả thi}
\rt{/portal/exam/registration/<int:reg_id>}
\medskip

\mvao

Bấm một dòng ở màn danh sách, hoặc được đưa tới ngay sau khi gửi phiếu.

\mai

Phiếu phải thuộc \textbf{cửa hàng đang chọn}. Phiếu của cửa hàng khác --- kể cả cửa hàng mà tôi
cũng truy cập được --- \textbf{không mở được} cho tới khi đổi sang cửa hàng đó. Sai thì về danh
sách, không có thông báo. Không chặn vai trò.

\mhien

\begin{hienbang}
Thông tin phiếu &
  \rt{wujia.exam.registration} &
  Mã phiếu, khóa thi, ngày giờ thi, địa điểm, cửa hàng, người đăng ký, ngày nộp, số người &
  --- &
  Phiếu UAT: \uat{1} người, nộp 03/08 \\ \addlinespace
Băng thông báo theo trạng thái &
  \rt{state} &
  Bốn trạng thái ra bốn băng khác nhau: chờ xác nhận (nhắc portal chưa hiện kết quả) · đã xác
  nhận · từ chối (kèm \textbf{lý do từ chối}) · đã huỷ (kèm \textbf{lý do huỷ}) &
  --- & --- \\ \addlinespace
Danh sách người dự thi &
  \rt{wujia.exam.registration.line} &
  Họ tên, điện thoại, năm sinh, chức vụ, ảnh &
  Theo thứ tự nhập &
  \uat{0}/1 người có ảnh \\ \addlinespace
Cột kết quả &
  \rt{result} &
  \textbf{Chỉ hiện khi kỳ thi đã công bố kết quả}; chưa công bố thì mọi người đều hiện
  \emph{Chưa có} dù Ngô Gia đã nhập điểm trong ERP &
  --- &
  \uat{1} người Đạt \\ \addlinespace
Bảng kết quả (máy tính) &
  --- &
  Chỉ hiện khi phiếu \textbf{đã được xác nhận}. Phiếu bị từ chối hoặc bị huỷ không có bảng này. &
  --- & --- \\ \addlinespace
Lời nhắc thi lại &
  --- &
  Chỉ hiện khi đã công bố \textbf{và} thật sự có người không đạt &
  --- & --- \\
\end{hienbang}

\mloc

Không có ô lọc, không có ô nhập.

\mkiem{--- }

\textbf{Màn này không có nút nào ghi dữ liệu.} Không sửa, không huỷ, không gửi lại.

\mghi

\textbf{Không ghi gì.}

\mhoi

\begin{ask}
\begin{enumerate}[leftmargin=*]
  \item \textbf{Portal không huỷ được phiếu.} Cửa hàng đăng ký nhầm phải gọi Ngô Gia; trong
        ERP người huỷ còn bắt buộc nhập lý do. Anh có muốn cho cửa hàng tự huỷ khi phiếu còn
        \emph{Chờ xác nhận} không?
  \item \textbf{Phiếu bị từ chối là ngõ cụt} --- không sửa, không gửi lại; phải tạo phiếu mới
        từ đầu, gõ lại toàn bộ người dự thi.
  \item \textbf{Chỉ mở được phiếu của cửa hàng đang chọn}, dù luật dữ liệu cho phép xem mọi cửa
        hàng của mình --- cùng câu hỏi với màn danh sách.
\end{enumerate}
\end{ask}

% =============================================================================
\section{Ảnh thí sinh trong phiếu}
\rt{/portal/exam/line/<int:line_id>/photo}
\medskip

\mvao

Không phải màn --- là ảnh hiện trong bảng người dự thi ở màn chi tiết.

\mai

\begin{kiembang}
1 & Đang chọn cửa hàng & Từ chối truy cập \\
2 & Dòng người dự thi thuộc \textbf{cửa hàng đang chọn} & Báo không tìm thấy \\
3 & Người đó thật sự có ảnh & Báo không tìm thấy \\
\end{kiembang}

\mhien

Trả về ảnh đã thu nhỏ, kèm cơ chế đệm của trình duyệt nên mở lại không tải lại.

\mloc

Không có.

\mkiem{ảnh}

Xem bảng trên.

\mghi

\textbf{Không ghi gì.}

\mhoi

\begin{ask}
Không có điểm cần chốt.
\end{ask}

% =============================================================================
\section{Đường dẫn cũ của phân hệ Đăng ký thi}
\rt{/portal/exam-registration}
\medskip

Đường dẫn của bản portal cũ (v14) được giữ lại và \textbf{chuyển hướng vĩnh viễn} sang
\rt{/portal/exam}. Lý do giữ: khoảng 1.500 tài khoản cửa hàng có thể còn lưu liên kết cũ trong
trình duyệt hoặc trong email Ngô Gia đã gửi.

Đường dẫn này \textbf{không đòi đăng nhập} --- nó chỉ chuyển hướng; màn đích mới hỏi đăng nhập.
\textbf{Không ghi gì.}

\begin{ask}
Không có điểm cần chốt --- đây là phần giữ tương thích, BA không cần lên task.
\end{ask}

\chapter{Khảo sát cửa hàng --- Metabase}

Chương này gồm \textbf{17 đường dẫn} còn lại, chia ba phần:

\begin{itemize}[leftmargin=*]
  \item \textbf{Phần A --- phía cửa hàng} (\rt{wujia_portal_inspection}, 8 đường dẫn): cửa hàng
        xem kết quả khảo sát và gửi khắc phục.
  \item \textbf{Phần B --- phía người đi khảo sát} (\rt{wujia_franchise_inspection}, 7 đường
        dẫn): màn nhân sự Ngô Gia dùng khi xuống cửa hàng chấm điểm. \textbf{Không phải màn
        portal của cửa hàng}, nhưng chạy trên cùng máy chủ và cùng đường dẫn web nên đưa vào
        đây cho đủ bản đồ.
  \item \textbf{Phần C --- Metabase} (\rt{wujia_metabase_connector}, 2 đường dẫn): nhúng bảng
        số liệu từ hệ thống báo cáo ngoài.
\end{itemize}

\textbf{Lưu ý về quyền sở hữu mã:} toàn bộ chương này là mã của anh Thái. Phiên tài liệu chỉ
\textbf{đọc để mô tả}, không sửa một dòng nào. Những điểm dev thấy cần bàn được ghi trong khung
vàng như mọi chương khác --- BA đọc để nắm, phần xử lý dev sẽ trao đổi trực tiếp với anh Thái.

\textbf{Ba điểm cần biết trước khi đọc:}

\begin{enumerate}[leftmargin=*]
  \item \textbf{Phần A phạm vi là MỌI cửa hàng tôi truy cập được} --- không phải cửa hàng đang
        chọn. Đây là cách làm giống Đổi trả, khác Đăng ký thi và Công nợ.
  \item \textbf{Điểm hiện trên portal được tính lại bằng công thức riêng của portal}, không lấy
        thẳng điểm đã chốt trong ERP --- xem mục Màn chi tiết. Hai nơi có thể ra hai con số.
  \item \textbf{Phần B và phần C chỉ kiểm ``đã đăng nhập'' ở phần lớn đường dẫn} --- chi tiết
        ghi trong từng mục.
\end{enumerate}

Số thật trên UAT: \uat{1} phiếu khảo sát (cửa hàng \emph{TP HCM Quận 1}, ngày 01/09, trạng thái
\emph{Chờ phản hồi}, điểm 86, xếp loại B), \uat{4} dòng tiêu chí trong đó \uat{2} không đạt,
\uat{0} tiêu chí đã được cửa hàng gửi khắc phục, \uat{0} dòng điểm danh, \uat{1} bảng Metabase.

% =============================================================================
\section{Màn Danh sách phiếu khảo sát}
\rt{/portal/inspection} · \rt{/portal/inspection/ajax}
\medskip

\mvao

Mục \emph{Khảo sát} trên thanh điều hướng. Đường dẫn \rt{/ajax} không phải màn riêng --- là
phần kết quả gọi ngầm khi đổi thẻ hoặc lọc. Tham số: \rt{tab}, \rt{date_from}, \rt{date_to},
\rt{search}, \rt{page}.

\mai

\begin{itemize}[leftmargin=*]
  \item Phải đăng nhập. \textbf{Không chặn vai trò.}
  \item Thấy phiếu của \textbf{mọi cửa hàng tôi truy cập được}.
  \item Không truy cập được cửa hàng nào thì \textbf{không thấy phiếu nào} --- điều kiện lọc cố
        ý để rỗng thay vì bỏ qua, nên không có kẽ hở xem nhầm cửa hàng khác.
  \item \textbf{Chỉ thấy phiếu đã chấm xong}: trạng thái \emph{Chờ phản hồi} hoặc
        \emph{Hoàn thành}. Phiếu người khảo sát đang làm dở (còn nháp) portal không thấy.
\end{itemize}

\mhien

\begin{hienbang}
Danh sách phiếu &
  \rt{wujia.franchise.inspection} &
  Phiếu thuộc cửa hàng của tôi, trạng thái \emph{Chờ phản hồi} hoặc \emph{Hoàn thành} &
  Ngày khảo sát mới nhất trước; \textbf{5} phiếu mỗi trang, \textbf{cố định} &
  \uat{1} phiếu · \uat{0} phiếu nháp \\ \addlinespace
Huy hiệu trạng thái &
  \rt{state} &
  Hai nhãn: \emph{Chờ phản hồi} (vàng) và \emph{Hoàn thành} (xanh) &
  --- &
  \uat{1} Chờ phản hồi \\ \addlinespace
Điểm và xếp loại &
  \rt{total_score} · \rt{wujia.franchise.inspection.grade} &
  Điểm tổng của phiếu; xếp loại lấy từ phiếu, nếu phiếu chưa gắn xếp loại thì
  \textbf{tra lại theo thang điểm} &
  --- &
  \uat{86} điểm, loại B \\ \addlinespace
Số tiêu chí đạt / không đạt &
  \rt{wujia.franchise.inspection.line} &
  Đếm các dòng tiêu chí; \emph{không đạt} = dòng bị đánh trượt &
  --- &
  \uat{2}/4 tiêu chí không đạt \\ \addlinespace
Nút \emph{Gửi khắc phục} &
  --- &
  Chỉ hiện khi phiếu đang \emph{Chờ phản hồi} và còn tiêu chí trượt \textbf{chưa gửi khắc
  phục}; nút trỏ thẳng tới tiêu chí đầu tiên còn thiếu &
  --- & --- \\ \addlinespace
Tên phiếu &
  \rt{name} &
  Phiếu đặt tên chung chung (\emph{``Khảo sát cửa hàng nhượng quyền''}) thì màn \textbf{tự đổi}
  thành \emph{``Lần dd/mm/yyyy''} cho dễ phân biệt &
  --- & --- \\
\end{hienbang}

\mloc

\begin{itemize}[leftmargin=*]
  \item \textbf{Ba thẻ}: \emph{Tất cả} · \emph{Chờ phản hồi} · \emph{Hoàn thành}.
  \item \textbf{Từ khoá} --- khớp một phần theo \textbf{tên phiếu}, \textbf{tên cửa hàng} hoặc
        \textbf{mã cửa hàng}.
  \item \textbf{Khoảng ngày} --- lọc theo \textbf{ngày khảo sát}; nhận cả dạng
        \emph{dd/mm/yyyy} lẫn \emph{yyyy-mm-dd}.
  \item \textbf{Không đổi được số dòng mỗi trang} --- cố định 5.
\end{itemize}

\mkiem{một phiếu}

Không kiểm gì ở danh sách.

\mghi

\textbf{Không ghi gì.}

\mhoi

\begin{ask}
\begin{enumerate}[leftmargin=*]
  \item \textbf{Gõ ngày ngược thì màn hiện rỗng, không báo gì} --- khác năm màn khác trong
        portal (Giao hàng, Báo cáo, Thông báo, Đăng ký thi, Lịch sử mua) vốn giữ hai ô đã gõ và
        báo tại thanh lọc. Anh có muốn đồng bộ không?
  \item \textbf{Cỡ trang cố định 5} --- nhỏ hơn hẳn các màn khác (10 hoặc 12) và không đổi được.
  \item \textbf{Phiếu còn nháp portal không thấy} --- đúng ý anh chứ? (Nghĩa là cửa hàng chỉ
        biết kết quả sau khi người khảo sát bấm hoàn tất.)
\end{enumerate}
\end{ask}

% =============================================================================
\section{Màn Chi tiết phiếu khảo sát}
\rt{/portal/inspection/detail/<int:inspection_id>}
\medskip

\mvao

Bấm một phiếu ở màn danh sách.

\mai

\begin{kiembang}
1 & Phiếu tồn tại & Về danh sách \\
2 & Trạng thái là \emph{Chờ phản hồi} hoặc \emph{Hoàn thành} & Về danh sách \\
3 & Phiếu thuộc \textbf{một trong các cửa hàng của tôi} & Về danh sách \\
\end{kiembang}

Cả ba trường hợp đều đưa về danh sách \textbf{không có thông báo}.

\mhien

\begin{hienbang}
Đầu phiếu &
  \rt{wujia.franchise.inspection} &
  Mã cửa hàng, ngày khảo sát, điểm tổng, xếp loại &
  --- &
  \uat{86}/100, loại B \\ \addlinespace
Bảng tóm tắt danh mục &
  \rt{wujia.franchise.inspection.line} &
  Mỗi danh mục một dòng. Danh mục thường hiện \emph{điểm đạt / điểm tối đa}; danh mục
  \textbf{vi phạm nghiêm trọng} hiện \textbf{số lỗi} thay cho điểm. &
  Theo thứ tự trong phiếu & --- \\ \addlinespace
Danh sách tiêu chí &
  \rt{wujia.franchise.inspection.line} &
  Nội dung tiêu chí, đạt hay không, điểm trừ, ghi chú của người khảo sát, ảnh hiện trường,
  ghi chú khắc phục và ảnh khắc phục của cửa hàng &
  Theo thứ tự trong phiếu &
  \uat{4} tiêu chí \\ \addlinespace
Nút \emph{Gửi khắc phục} &
  --- &
  Chỉ hiện khi phiếu đang \emph{Chờ phản hồi} và còn tiêu chí trượt chưa phản hồi &
  --- & --- \\
\end{hienbang}

\textbf{Điểm trên màn này được tính lại tại chỗ, không lấy điểm đã chốt trong ERP.} Cách tính
của portal: mỗi tiêu chí đạt được cộng số điểm bằng \textbf{điểm trừ của chính nó}; tiêu chí
\textbf{không khai điểm trừ thì được tính là 4 điểm}; và số điểm trừ lẻ bị \textbf{cắt phần
thập phân} (2,5 thành 2). Cột \emph{điểm tối đa} của mỗi danh mục cũng cộng theo cách đó, còn
\textbf{tổng điểm tối đa của phiếu luôn ghi cứng là 100}.

\textbf{Cách nhận biết ``vi phạm nghiêm trọng''} là \textbf{dò chữ}: hệ thống tìm cụm chữ
\emph{``nghiêm trọng''} trong tên danh mục hoặc trong nội dung tiêu chí. Đổi cách đặt tên danh
mục là cách nhận biết này hỏng.

\mloc

Không có ô lọc.

\mkiem{Gửi khắc phục}

Xem hai mục kế tiếp.

\mghi

Mở màn \textbf{không ghi gì}.

\mhoi

\begin{ask}
\begin{enumerate}[leftmargin=*]
  \item \textbf{Điểm trên portal và điểm trong ERP có thể lệch nhau} --- portal tự cộng lại
        theo quy ước riêng (tiêu chí không khai điểm trừ tính 4 điểm, điểm lẻ bị cắt phần thập
        phân), trong khi ERP dùng công thức chấm điểm của phiếu. Anh cần chốt: portal hiện đúng
        con số ERP, hay giữ cách tính riêng này?
  \item \textbf{Tổng điểm tối đa ghi cứng 100} --- bộ tiêu chí nào có thang khác 100 sẽ hiện
        sai tỉ lệ.
  \item \textbf{``Vi phạm nghiêm trọng'' nhận biết bằng cách dò chữ ``nghiêm trọng''} thay vì
        dùng cờ đánh dấu. Anh có muốn Ngô Gia đánh dấu hẳn một cờ trong danh mục không?
  \item \textbf{Có vài giá trị mẫu còn sót trong màn} --- phiếu thiếu mã cửa hàng, thiếu ngày
        hoặc thiếu điểm thì màn hiện giá trị mẫu (mã \emph{HN\_ST042}, ngày \emph{15/05/2024},
        điểm \emph{85}, loại \emph{B+}) thay vì để trống. Dev sẽ trao đổi với anh Thái, BA
        không cần lên task.
\end{enumerate}
\end{ask}

% =============================================================================
\section{Màn Gửi khắc phục cho một tiêu chí}
\rt{/portal/inspection/remediation/<int:line_id>} ·
\rt{/portal/remediation/<int:line_id>} · \rt{/portal/remediation}
\medskip

\mvao

Nút \emph{Gửi khắc phục} ở màn danh sách hoặc màn chi tiết. Ba đường dẫn cùng một màn:
đường dẫn đầy đủ, đường dẫn rút gọn, và \textbf{đường dẫn thiếu mã tiêu chí} --- gõ
\rt{/portal/remediation} không kèm số thì bị đưa về danh sách.

\mai

\begin{kiembang}
1 & Có mã tiêu chí & Về danh sách \\
2 & Tiêu chí tồn tại và thuộc một phiếu khảo sát có thật & Về danh sách \\
3 & Phiếu thuộc \textbf{một trong các cửa hàng của tôi} & Về danh sách \\
\end{kiembang}

\mhien

\begin{hienbang}
Thông tin tiêu chí &
  \rt{wujia.franchise.inspection.line} &
  Tên danh mục, mã tiêu chí, nội dung tiêu chí, \textbf{ghi chú của người khảo sát} &
  --- & --- \\ \addlinespace
Ảnh hiện trường &
  \rt{evidence_image} &
  Ảnh người khảo sát chụp lúc chấm --- tải qua đường dẫn ảnh chuẩn của hệ thống, có kiểm quyền
  theo cửa hàng &
  --- & --- \\ \addlinespace
Mốc thời gian &
  \rt{create_date} &
  Giờ và ngày \textbf{tạo phiếu khảo sát} --- không phải ngày khảo sát ghi trong phiếu &
  --- & --- \\
\end{hienbang}

\mloc

Hai ô: \textbf{ghi chú khắc phục} và \textbf{ảnh minh chứng} (chụp hoặc chọn từ máy).

\mkiem{Gửi}

Xem mục kế tiếp.

\mghi

Mở màn \textbf{không ghi gì}.

\mhoi

\begin{ask}
\textbf{Mốc thời gian trên màn là giờ tạo phiếu, không phải ngày khảo sát} --- hai con số này
lệch nhau khi người khảo sát tạo phiếu trước rồi vài hôm sau mới đi. Nên đổi sang ngày khảo sát.
\end{ask}

% =============================================================================
\section{Gửi nội dung khắc phục}
\rt{/portal/inspection/remediation/submit} · \rt{/portal/remediation/submit}
\medskip

\mvao

Không phải màn --- nút \emph{Gửi} ở màn khắc phục. Hai đường dẫn dẫn tới cùng một việc.

\mai

Chỉ kiểm \textbf{đã đăng nhập}. \textbf{Không} kiểm tiêu chí có thuộc cửa hàng của người gửi
hay không --- khác hẳn màn hiển thị ngay trên, vốn kiểm đủ ba lớp.

\mhien

Không hiện màn --- trả về câu báo thành công kèm đường dẫn quay lại màn chi tiết.

\mloc

Hai ô như mục trên. Ảnh nhận được ở \textbf{bốn dạng} khác nhau (tệp tải lên, chuỗi ảnh nhúng,
chuỗi mã hoá, dữ liệu nhị phân) --- viết vậy để cả bản điện thoại lẫn bản máy tính dùng chung
một đường dẫn.

\mkiem{Gửi}

\begin{kiembang}
1 & Có mã tiêu chí & ``Thiếu ID tiêu chí.'' \\
2 & Tiêu chí tồn tại & ``Tiêu chí không tồn tại.'' \\
3 & Tiêu chí \textbf{bắt buộc ảnh} thì phải có ảnh (ảnh mới hoặc ảnh đã gửi trước đó) &
    ``Tiêu chí này yêu cầu bắt buộc phải có ảnh minh chứng.'' \\
4 & Phải có \textbf{ít nhất} ghi chú hoặc ảnh &
    ``Vui lòng nhập ghi chú khắc phục hoặc tải lên ảnh minh chứng trước khi gửi.'' \\
\end{kiembang}

\mghi

\begin{itemize}[leftmargin=*]
  \item Ghi \textbf{ghi chú khắc phục} và \textbf{ảnh khắc phục} lên dòng tiêu chí.
  \item Đổi trạng thái khắc phục của \textbf{riêng dòng đó} sang \emph{Đã khắc phục}.
  \item \textbf{Không} đổi trạng thái của cả phiếu --- phiếu vẫn ở \emph{Chờ phản hồi} cho tới
        khi Ngô Gia tự chuyển. Cửa hàng gửi khắc phục đủ cả bốn tiêu chí thì phiếu vẫn hiện
        \emph{Chờ phản hồi}.
  \item \textbf{Không} báo cho ai --- không email, không việc cần làm.
\end{itemize}

\mhoi

\begin{ask}
\begin{enumerate}[leftmargin=*]
  \item \textbf{Đường dẫn gửi khắc phục không kiểm cửa hàng.} Người đăng nhập bất kỳ, nếu biết
        mã của một dòng tiêu chí, có thể ghi ghi chú và ảnh lên tiêu chí của cửa hàng khác.
        Đây là điểm dev sẽ trao đổi với anh Thái --- \textbf{BA không cần lên task}, nhưng cần
        biết vì nó ảnh hưởng tới câu ``ai xem được gì''.
  \item \textbf{Ảnh khắc phục không kiểm định dạng, không kiểm dung lượng} --- khác hẳn Đổi trả
        (chương 6) và Đăng ký thi (chương 8). Cùng nhóm với điểm trên.
  \item \textbf{Gửi khắc phục xong phiếu vẫn ``Chờ phản hồi''} và không ai được báo. Anh muốn
        phiếu tự chuyển trạng thái khi mọi tiêu chí trượt đã có phản hồi không?
  \item \textbf{Gửi rồi vẫn gửi lại được} --- lần gửi sau đè lên lần trước, không lưu lịch sử.
\end{enumerate}
\end{ask}

% =============================================================================
\section{Màn Làm khảo sát tại cửa hàng}
\rt{/franchise/inspection/do/<int:inspection_id>}
\medskip

\mvao

Nhân sự Ngô Gia mở từ phiếu khảo sát trong ERP khi xuống cửa hàng. \textbf{Không nằm trên thanh
điều hướng của portal cửa hàng.}

\mai

\begin{itemize}[leftmargin=*]
  \item Phải đăng nhập.
  \item Mở được nếu là \textbf{quản trị hệ thống}, là \textbf{người được giao khảo sát}, hoặc
        là \textbf{bất kỳ nhân sự nội bộ nào}.
  \item \textbf{Phiếu chưa giao cho ai thì không kiểm gì} --- bất kỳ tài khoản đăng nhập nào
        cũng mở được, kể cả tài khoản cửa hàng. Ghi lại ở khung vàng bên dưới.
\end{itemize}

\mhien

\begin{hienbang}
Danh sách tiêu chí &
  \rt{wujia.franchise.inspection.line} &
  Toàn bộ dòng của phiếu, chia theo danh mục, kèm điểm trừ và các cờ ``bắt buộc ghi chú'' /
  ``bắt buộc ảnh'' &
  Theo thứ tự trong phiếu &
  \uat{4} tiêu chí \\ \addlinespace
Kết quả lần khảo sát trước &
  \rt{previous_line_id} &
  Với mỗi tiêu chí: lần trước đạt hay trượt, ghi chú và ảnh lần trước, tên người khảo sát lần
  trước &
  --- & --- \\ \addlinespace
Bài kiểm tra nhân viên &
  \rt{exam_line_ids} &
  Câu hỏi kiểm tra kiến thức cho nhân viên được chọn &
  --- & --- \\ \addlinespace
Bảng điểm danh &
  \rt{...attendance.line} &
  Nhân sự có mặt lúc khảo sát &
  --- &
  \uat{0} dòng \\ \addlinespace
Báo cáo doanh thu 3 tháng &
  \rt{report_line_ids} &
  Doanh thu, doanh thu bình quân, tỉ lệ hàng hoá, doanh số qua ứng dụng &
  --- & --- \\ \addlinespace
Ngôn ngữ nội dung &
  --- &
  Nội dung tiêu chí lưu \textbf{song ngữ trong cùng một ô}; màn \textbf{tự tách} phần tiếng
  Trung hoặc tiếng Việt theo ngôn ngữ tài khoản &
  --- & --- \\
\end{hienbang}

\mloc

Chấm đạt / không đạt từng tiêu chí, nhập ghi chú, chụp ảnh hiện trường, trả lời bài kiểm tra,
điểm danh nhân sự, ký tên.

\mkiem{Lưu}

Xem mục kế tiếp.

\mghi

Mở màn \textbf{không ghi gì}.

\mhoi

\begin{ask}
\textbf{Phiếu chưa giao cho người khảo sát nào thì ai đăng nhập cũng mở được màn này}, kể cả
tài khoản cửa hàng --- và màn này hiện cả kết quả lần trước lẫn bài kiểm tra. Dev sẽ trao đổi
với anh Thái; BA chỉ cần biết để trả lời khi cửa hàng hỏi.
\end{ask}

% =============================================================================
\section{Lưu bài khảo sát}
\rt{/franchise/inspection/do/<int:inspection_id>/save}
\medskip

\mvao

Không phải màn --- nút \emph{Lưu tạm} và nút \emph{Hoàn tất} ở màn làm khảo sát.

\mai

Chỉ kiểm \textbf{đã đăng nhập} và phiếu chưa khoá. \textbf{Không} kiểm người gửi có phải người
được giao khảo sát hay không.

\mhien

Trả về điểm vừa tính lại, xếp loại, trạng thái phiếu và kết quả bài kiểm tra.

\mloc

Nhận cả gói: các dòng tiêu chí, bài kiểm tra, tên nhân viên được kiểm, thâm niên, tình trạng
mặt bằng, người xác nhận, bảng điểm danh, ảnh chữ ký, báo cáo doanh thu.

\mkiem{Hoàn tất}

\begin{kiembang}
1 & Phiếu tồn tại & ``Inspection sheet does not exist'' \\
2 & Phiếu \textbf{chưa hoàn tất, chưa bị huỷ} &
    ``Inspection sheet is already completed or cancelled and cannot be edited!'' \\
3 & Bấm \emph{Hoàn tất} thì phải có \textbf{tên nhân viên được kiểm} &
    ``Please enter staff name before saving!'' \\
\end{kiembang}

Bấm \emph{Lưu tạm} thì bỏ qua bước 3 và \textbf{không khoá bài kiểm tra}; bấm \emph{Hoàn tất}
thì bài kiểm tra bị chốt, không sửa lại được.

\mghi

\begin{itemize}[leftmargin=*]
  \item Ghi từng dòng tiêu chí --- \textbf{chỉ ghi dòng thật sự thay đổi}, so sánh trước khi
        ghi để không đụng dữ liệu thừa.
  \item Chấm bài kiểm tra và ghi điểm từng câu.
  \item Ghi bảng điểm danh; nhân sự mới trong bảng điểm danh \textbf{được tạo thành thành viên
        cửa hàng} luôn.
  \item Ghi lại bảng doanh thu 3 tháng --- dòng không còn gửi lên thì \textbf{bị xoá}.
  \item Tính lại điểm phần tiêu chí, điểm bài kiểm tra, điểm tổng và xếp loại.
  \item \textbf{Có tiêu chí trượt thì phiếu tự chuyển từ nháp sang \emph{Chờ phản hồi}} ---
        đây chính là lúc phiếu hiện ra trên portal cửa hàng.
\end{itemize}

\mhoi

\begin{ask}
\begin{enumerate}[leftmargin=*]
  \item \textbf{Đường dẫn lưu không kiểm người gửi.} Bất kỳ tài khoản đăng nhập nào biết mã
        phiếu đều ghi được điểm, ghi chú, chữ ký và bảng điểm danh lên phiếu đang mở. Cùng nhóm
        với hai điểm trên --- dev trao đổi với anh Thái.
  \item \textbf{Điểm danh tự tạo thành viên cửa hàng.} Gõ nhầm tên là sinh một thành viên mới
        trong hồ sơ cửa hàng. Anh có muốn bước xác nhận trước khi tạo không?
  \item \textbf{Phiếu tự chuyển sang Chờ phản hồi ngay khi lưu có tiêu chí trượt} --- kể cả khi
        người khảo sát mới chỉ \emph{lưu tạm}. Cửa hàng sẽ thấy phiếu trước khi người khảo sát
        làm xong.
\end{enumerate}
\end{ask}

% =============================================================================
\section{Kéo doanh thu từ PosApp}
\rt{/franchise/inspection/do/<int:inspection_id>/sync_posapp}
\medskip

\mvao

Không phải màn --- nút đồng bộ trong phần báo cáo doanh thu ở màn làm khảo sát.

\mai

Chỉ kiểm đã đăng nhập và phiếu chưa khoá.

\mhien

Trả về ba dòng doanh thu vừa kéo về để màn vẽ lại bảng.

\mloc

Không có ô nhập --- bấm là chạy.

\mkiem{Đồng bộ}

\begin{kiembang}
1 & Phiếu tồn tại & ``Inspection sheet does not exist'' \\
2 & Phiếu chưa hoàn tất, chưa bị huỷ &
    ``Inspection sheet is already completed or cancelled!'' \\
3 & Gọi được sang PosApp & Trả về nguyên văn câu lỗi kỹ thuật của hệ thống \\
\end{kiembang}

\mghi

\textbf{Xoá sạch bảng doanh thu cũ rồi ghi lại bảng mới} --- số nhân sự gõ tay trước đó mất
hết, không hỏi lại.

\mhoi

\begin{ask}
\begin{enumerate}[leftmargin=*]
  \item \textbf{Bấm đồng bộ là xoá sạch số đã gõ tay}, không có bước xác nhận. Anh có muốn hỏi
        lại trước khi ghi đè không?
  \item \textbf{Lỗi kết nối hiện nguyên văn câu lỗi kỹ thuật} --- cùng nhóm chuẩn hoá thông báo
        lỗi mà dev đang làm.
\end{enumerate}
\end{ask}

% =============================================================================
\section{Bảng điểm danh nhân sự}
\rt{.../attendance/add} · \rt{.../attendance/save_member} ·
\rt{.../attendance/deactivate_member} · \rt{.../attendance/delete_line}
\medskip

\mvao

Không phải màn --- bốn nút trong bảng điểm danh ở màn làm khảo sát: thêm dòng, lưu dòng thành
thành viên cửa hàng, ngưng hiệu lực thành viên, xoá dòng.

\mai

Chỉ kiểm \textbf{đã đăng nhập}. Ba đường dẫn sau còn kiểm dòng điểm danh \textbf{đúng thuộc
phiếu ghi trên đường dẫn} --- nên không mượn được để sửa dòng của phiếu khác.

\mhien

Mỗi lần gọi đều trả về \textbf{bảng nhân sự của cửa hàng} và \textbf{số người có mặt} để màn
cập nhật ngay.

\mloc

Họ tên (bắt buộc khi thêm), vai trò, số điện thoại, ghi chú, có mặt hay không.

\mkiem{Thêm nhân sự}

\begin{kiembang}
1 & Phiếu tồn tại và chưa hoàn tất / chưa bị huỷ &
    ``Inspection survey is locked or does not exist!'' \\
2 & Có họ tên & ``Please enter staff full name!'' \\
3 & (ba nút còn lại) Dòng điểm danh tồn tại và thuộc đúng phiếu &
    ``Attendance line does not exist!'' \\
\end{kiembang}

\mghi

\begin{itemize}[leftmargin=*]
  \item \textbf{Thêm}: tạo dòng điểm danh, rồi \textbf{tự tạo thành viên cửa hàng} cho người đó
        --- tạo không được thì bỏ qua, dòng điểm danh vẫn còn.
  \item \textbf{Lưu thành thành viên}: sửa tên, vai trò, điện thoại trên dòng rồi tạo hoặc cập
        nhật thành viên cửa hàng tương ứng.
  \item \textbf{Ngưng hiệu lực}: tắt hiệu lực thành viên cửa hàng --- \textbf{đây là thao tác
        sửa hồ sơ cửa hàng từ màn khảo sát}.
  \item \textbf{Xoá dòng}: xoá hẳn dòng điểm danh; thành viên cửa hàng đã tạo \textbf{vẫn còn}.
\end{itemize}

\mhoi

\begin{ask}
\begin{enumerate}[leftmargin=*]
  \item \textbf{Màn khảo sát sửa được hồ sơ nhân sự của cửa hàng} --- thêm thành viên, ngưng
        hiệu lực thành viên. Anh xác nhận cho phép; nếu không thì cần tách quyền.
  \item \textbf{Xoá dòng điểm danh không gỡ thành viên đã tạo} --- gõ nhầm rồi xoá dòng thì
        người đó vẫn nằm trong danh sách nhân sự cửa hàng.
\end{enumerate}
\end{ask}

% =============================================================================
\section{Bảng số liệu Metabase}
\rt{/metabase/embed/<int:dashboard_id>} · \rt{/metabase/embed_url/<int:dashboard_id>}
\medskip

\mvao

Không nằm trên thanh điều hướng portal. Đường dẫn thứ nhất mở trang có khung nhúng; đường dẫn
thứ hai chỉ trả về địa chỉ nhúng cho màn khác tự gắn vào.

\mai

\begin{kiembang}
1 & Đã đăng nhập & Trang đăng nhập \\
2 & Bảng tồn tại và đang bật & Báo không tìm thấy \\
3 & Bảng cho phép \textbf{công ty} hiện tại &
    Trang báo ``Access Denied'' kèm tên công ty \\
4 & \textbf{Nếu bảng có khai nhóm quyền}: tôi thuộc một trong các nhóm đó (hoặc là quản trị) &
    Trang báo ``Access Denied'' \\
5 & Có khai máy chủ Metabase và máy chủ đang bật & Trang báo ``Configuration Error'' \\
6 & Máy chủ có thư viện ký chữ ký số & Trang báo ``System Error'' \\
\end{kiembang}

\textbf{Bước 4 chỉ chạy khi bảng có khai nhóm quyền.} Bảng \textbf{không khai nhóm nào} thì bỏ
qua hoàn toàn --- mọi tài khoản đăng nhập đều xem được. Trên UAT, bảng duy nhất
(\emph{``Báo cáo nhượng quyền''}) \textbf{không khai nhóm nào}.

\mhien

\begin{hienbang}
Khung nhúng &
  \rt{metabase.dashboard} &
  Bảng số liệu nằm trên máy chủ Metabase ngoài, gọi qua một chữ ký số có hạn &
  --- &
  \uat{1} bảng, đang bật \\ \addlinespace
Hạn của chữ ký &
  \rt{token_expiry_minutes} &
  Mặc định \textbf{10 phút}; hết hạn thì khung nhúng báo lỗi, phải tải lại trang &
  --- & --- \\ \addlinespace
Chiều cao khung &
  \rt{height} &
  Theo cấu hình của bảng, mặc định 800 &
  --- & --- \\
\end{hienbang}

\mloc

Không có ô lọc --- mọi thao tác lọc diễn ra bên trong khung Metabase, hệ thống này không can
thiệp. \textbf{Không truyền tham số lọc nào sang Metabase} --- nghĩa là bảng số liệu
\textbf{không tự giới hạn theo cửa hàng của người xem}.

\mkiem{mở bảng}

Xem bảng ở mục \emph{Ai xem được gì}.

\mghi

\textbf{Không ghi gì.}

\mhoi

\begin{ask}
\begin{enumerate}[leftmargin=*]
  \item \textbf{Bảng không khai nhóm quyền thì ai đăng nhập cũng xem được} --- kể cả tài khoản
        cửa hàng. Bảng duy nhất trên UAT đang ở trạng thái đó. Nếu định mở Metabase cho cửa
        hàng thì phải khai nhóm trước.
  \item \textbf{Không truyền cửa hàng sang Metabase} --- người xem thấy số liệu toàn hệ, không
        phải số liệu cửa hàng mình. Anh xác nhận đây là bảng dành cho nội bộ Ngô Gia.
\end{enumerate}
\end{ask}

\chapter{Tổng hợp --- thứ BA cần để lên task}

Chương này gom lại toàn bộ kết luận của 9 chương trước thành bốn bảng tra:

\begin{enumerate}[leftmargin=*]
  \item \textbf{Đối chiếu CT-001…CT-071} --- cập nhật sau khi soi sâu, nói rõ ``một phần'' là
        thiếu đúng luật nào.
  \item \textbf{Code có nhưng tab CT chưa mô tả} --- để BA bổ sung.
  \item \textbf{Toàn bộ điểm cần BA xác nhận} --- gom từ mục 7 của mọi màn, có số trang để lật
        lại.
  \item \textbf{Việc BA KHÔNG cần lên task} --- phần dev tự chuẩn hoá, và phần chỉ đổi chỗ mã
        nguồn chứ không đổi đường dẫn.
\end{enumerate}

\section{Ba đính chính so với bản tài liệu 19/09}

Bản kiểm kê 20 trang ngày 19/09 có ba chỗ nay đã lạc hậu hoặc cần nói lại cho đúng:

\begin{enumerate}[leftmargin=*]
  \item \textbf{``Công nợ chưa nối hoá đơn kế toán'' --- không còn đúng.} Màn Công nợ
        \textbf{đang đọc hoá đơn thật} (\rt{account.move} đã vào sổ, loại hoá đơn bán và giảm
        giá) và \textbf{phiếu thu thật} (\rt{account.payment}). Điều cần chốt bây giờ không
        phải ``nối dữ liệu'' mà là \textbf{cách tính kỳ và cách quy đổi tiền} --- xem chương 6.
  \item \textbf{``Chưa dùng chung một hàm tải tệp'' --- đúng một nửa.} Hỗ trợ và Yêu cầu thông
        tin \textbf{đã} dùng chung một hàm (đếm tệp, giới hạn dung lượng, danh sách định dạng).
        Đổi trả dùng hàm riêng \textbf{chặt hơn} (đọc nội dung thật của tệp). Khắc phục khảo
        sát \textbf{không kiểm gì}. Vậy là ba mức khác nhau, không phải một.
  \item \textbf{Số đường dẫn: 104, không phải 103.} Máy đếm lại theo cấu trúc mã nguồn cho ra
        \textbf{104 đường dẫn} thuộc 16 mô-đun; trong đó \textbf{2 đường dẫn có hai hàm} (một
        hàm mở màn, một hàm nhận dữ liệu gửi lên): \rt{/portal/exam/register} và
        \rt{/portal/support/new}. Tài liệu này mô tả đủ \textbf{104/104}.
\end{enumerate}

\section{Đối chiếu CT-001…CT-071 --- bản cập nhật}

Cột \emph{Đọc ở} chỉ chương mô tả chi tiết. Cột \emph{Điều chỉnh sau khi soi sâu} chỉ ghi ở
những mục có thay đổi so với bản 19/09; để trống nghĩa là kết luận cũ vẫn đúng.

{\footnotesize
\begin{longtable}{@{}L{1.4cm}L{3.4cm}L{1.6cm}L{1.4cm}L{6.8cm}@{}}
\rowcolor{hdrbg}\textbf{Mục CT} & \textbf{Tên} & \textbf{Trạng thái} & \textbf{Đọc ở} &
\textbf{Điều chỉnh sau khi soi sâu}\\ \midrule \endfirsthead
\rowcolor{hdrbg}\textbf{Mục CT} & \textbf{Tên} & \textbf{Trạng thái} & \textbf{Đọc ở} &
\textbf{Điều chỉnh sau khi soi sâu}\\ \midrule \endhead
CT-001 & Đăng nhập portal & \ok & ch.4 & Màn đăng nhập \textbf{chưa chặn thử mật khẩu hàng loạt}, trong khi màn Quên mật khẩu có chặn. \\
CT-002 & Xác định user hiện tại & \ok & ch.4 & \\
CT-003 & Chọn cửa hàng làm việc & \ok & ch.4 & Lưu bằng \textbf{cookie 30 ngày}; \textbf{đăng xuất không xoá} --- máy dùng chung ở cửa hàng giữ nguyên lựa chọn cho người sau. \\
CT-004 & Lấy cửa hàng hiện tại & \ok & ch.4 & Có \textbf{bốn cách hiểu} ``cửa hàng đang xem'' trong portal --- xem mục \ref{sec:pham-vi}. \\
CT-005 & Hiển thị thông tin tài khoản & \ok & ch.4 & \\
CT-006 & Cập nhật thông tin tài khoản & \ok & ch.4 & Sửa email ở đây \textbf{không} đổi email đăng nhập. \\
CT-007 & Đổi mật khẩu & \ok & ch.4 & \textbf{Không} đăng xuất thiết bị khác; luật mật khẩu lệch nhẹ với màn Đặt lại. \\
CT-008 & Thông tin cửa hàng & \ok & ch.4 & Đang có \textbf{hai bộ màn song song} cho cùng nội dung (bộ Vuexy và bộ giao diện Odoo chuẩn). \\
CT-009 & Danh sách thành viên & \ok & ch.4 & Xác nhận: Nhân viên \textbf{xem được} đủ danh sách đồng nghiệp kèm điện thoại. \\
CT-010 & Dữ liệu tổng quan Home & \ok & ch.4 & \textbf{Chưa chọn cửa hàng thì các số là tổng mọi cửa hàng} --- không phải rỗng. \\
CT-011 & Thông báo mới nhất ở Home & \ok & ch.4/7 & Thẻ đếm chưa đọc ở Home \textbf{tính khác} màn Thông báo. \\
CT-012 & Đơn hàng gần đây & \ok & ch.4 & Mọi khối danh sách Home \textbf{cố định 2 dòng}. \\
CT-013 & Yêu cầu đổi trả gần đây & \ok & ch.4 & Thẻ đếm \textbf{bỏ sót trạng thái Đang xét}. \\
CT-014 & Công nợ tóm tắt & \ok & ch.4/6 & \textbf{Đính chính}: có đọc kế toán thật. Nhưng Home dùng \textbf{ô tổng lưu sẵn, tiền công ty, mọi thời điểm}, còn màn Công nợ tính \textbf{tiền gốc chứng từ, một tuần} --- hai số khác nhau. \\
CT-015 & Danh sách sản phẩm & \ok & ch.3 & Điều kiện đúng như BA nêu: cờ hiện trên portal + còn hiệu lực. \\
CT-016 & Danh mục sản phẩm & \ok & ch.3 & Dùng bảng danh mục riêng; sản phẩm \textbf{chưa gán danh mục vẫn hiện} ở ``Tất cả''. \\
CT-017 & Chi tiết sản phẩm & \ok & ch.3 & Khối ``sản phẩm liên quan'' dựa hoàn toàn vào danh mục portal. \\
CT-018 & Kiểm tra điều kiện đặt & \ok & ch.3 & UAT \textbf{chưa sản phẩm nào đặt trần tối đa}. \\
CT-019 & Khung giờ đặt hàng chung & \ok & ch.3 & Giờ tính theo \textbf{múi giờ tài khoản}, không theo khu vực cửa hàng. \\
CT-020 & Khung giờ theo khu vực & \ok & ch.3 & Khu vực Hà Nội (3/4 cửa hàng UAT) \textbf{chưa có khung riêng}. \\
CT-021 & Lưu giỏ hàng & \ok & ch.3 & Giỏ dùng chung cả cửa hàng; \textbf{không có lịch sử}, người sau đè người trước. \\
CT-022 & Tạo đơn hàng xuống Odoo & \ok & ch.3 & Gửi đơn lần hai trong ngày \textbf{huỷ đơn nháp cũ một cách im lặng}. \\
CT-023 & Trả kết quả đặt hàng & \ok & ch.3 & Máy tính và điện thoại kết thúc ở \textbf{hai màn khác nhau}. \\
CT-024 & Danh sách đơn hàng & \ok & ch.5 & Đơn \textbf{Đã huỷ bị giấu hoàn toàn} (UAT: 16/47 đơn). \\
CT-025 & Chi tiết đơn hàng & \ok & ch.5 & Không hiện lịch sử thay đổi của đơn. \\
CT-026 & Danh sách chuyến giao & \ok & ch.5 & Chuyến \textbf{chưa đặt lịch} biến mất khi lọc theo ngày. \\
CT-027 & Chi tiết giao hàng & \ok & ch.5 & Điều kiện xem chi tiết \textbf{hẹp hơn} điều kiện hiện ở danh sách. \\
CT-028 & Theo dõi trạng thái giao & \ok & ch.5 & Hai bước cuối của thanh tiến trình \textbf{chưa có giờ thật}. \\
CT-029 & Tạo yêu cầu đổi trả & \ok & ch.6 & \\
CT-030 & DS đơn được đổi trả & \ok & ch.6 & \textbf{UAT hiện 0 đơn lọt cửa sổ 10 ngày} --- không ai tạo được yêu cầu. Mốc tính từ \emph{ngày đặt}, không phải ngày nhận hàng. \\
CT-031 & Lịch sử yêu cầu đổi trả & \ok & ch.6 & Gộp mọi cửa hàng, \textbf{không có cột và bộ lọc cửa hàng}. \\
CT-032 & Chi tiết yêu cầu đổi trả & \ok & ch.6 & Xác nhận lại: một sản phẩm mỗi yêu cầu. Thêm: \textbf{video tải lên không bao giờ hiện lại}. \\
CT-033 & Danh sách lịch thi & \ok & ch.8 & \textbf{UAT không kỳ thi nào đăng ký được} (hai kỳ mở đều đã qua ngày thi). \\
CT-034 & Đăng ký thi & \ok & ch.8 & Phiếu \textbf{chờ duyệt đã chiếm ghế}; người dự thi gõ tay, không nối danh sách nhân sự. \\
CT-035 & Lịch sử đăng ký & \ok & ch.8 & Chỉ thấy \textbf{cửa hàng đang chọn}. \\
CT-036 & Kết quả thi & \ok & ch.8 & \\
CT-037 & Danh sách bài viết & \ok & ch.7 & Kiến thức \textbf{không phân biệt cửa hàng và vai trò}. \\
CT-038 & Danh mục / thẻ bài viết & \ok & ch.7 & UAT có \textbf{danh mục trùng tên} (Vận hành, Sản phẩm mỗi cái hai bản ghi). \\
CT-039 & Chi tiết bài viết & \ok & ch.7 & Bài hết hạn \textbf{chỉ rụng khi tác vụ nền chạy} (hằng ngày). \\
CT-040 & Ghi lượt xem / đã đọc & \pt & ch.7 & Vẫn đúng: có đếm lượt xem, \textbf{không} ghi ai đã đọc. Thêm: lượt xem \textbf{không chống bấm lặp}. \\
CT-041 & Danh sách thông báo & \ok & ch.7 & UAT \textbf{chưa có thông báo gửi đích danh cửa hàng} nào đang phát hành. \\
CT-042 & Thông báo trên chuông & \ok & ch.7 & Chuông \textbf{cố định 5 dòng}; bỏ thông báo hết hiệu lực, khác màn danh sách. \\
CT-043 & Đánh dấu đã đọc & \ok & ch.7 & Ghi theo \textbf{người + thông báo + cửa hàng}: đọc ở cửa hàng A rồi sang B lại thành chưa đọc. UAT có \uat{6} dấu cũ chưa gắn cửa hàng. \\
CT-044 & Chi tiết thông báo & \ok & ch.7 & Mở khi \textbf{chưa chọn cửa hàng} thì không được tính là đã đọc. \\
CT-045 & Tạo ticket hỗ trợ & \ok & ch.7 & Gửi xong \textbf{không ai được báo}. \\
CT-046 & Danh sách ticket & \ok & ch.7 & Xác nhận lại: lọc theo \textbf{người tạo}. Thêm số thật: UAT có \uat{15} phiếu, \uat{0} do tài khoản cửa hàng tạo. \\
CT-047 & Chi tiết ticket & \ok & ch.7 & Phiếu \textbf{Đã huỷ} bị giấu ở danh sách nhưng chi tiết \textbf{vẫn mở được}. \\
CT-048 & Gửi phản hồi trong ticket & \ok & ch.7 & Trả lời \textbf{không đổi trạng thái} phiếu; phiếu đã đóng vẫn nhận lời nhắn. \\
CT-049 & Lịch sử trao đổi ticket & \ok & ch.7 & Portal hiện \textbf{mọi lời nhắn} --- nhân sự Ngô Gia phải dùng đúng ô ghi chú nội bộ nếu muốn giấu. \\
CT-050 & Tổng quan công nợ & \ok & ch.6 & \textbf{Đính chính} như CT-014. Thêm: vai trò xét \textbf{tại cửa hàng đang chọn}; chưa chọn cửa hàng thì cổng vai trò không chạy. \\
CT-051 & Danh sách hoá đơn công nợ & \ok & ch.6 & Trạng thái \emph{Dư có} \textbf{đè lên} \emph{Có quá hạn}, làm mất nút thanh toán. \\
CT-052 & Lọc theo kỳ công nợ & \ok & ch.6 & Dropdown \textbf{chỉ 6 tuần} --- hoá đơn cũ hơn vĩnh viễn không chọn được. Hạn thanh toán do portal tự tính (thứ Hai + 10 ngày), \textbf{không} đọc hạn của hoá đơn. \\
CT-053 & Tải PDF công nợ & \no & ch.6 & Vẫn chưa có. \\
CT-054 & Lịch sử thanh toán & \ok & ch.6 & Tổng \textbf{tách theo loại tiền, không quy đổi}; kỳ mở sẵn là tháng hiện tại, trên UAT tháng đó rỗng. \\
CT-055 & QR thanh toán & \pt & ch.6 & Vẫn treo ảnh QR. Thêm: nội dung chuyển khoản \textbf{làm tròn xuống số nguyên}. \\
CT-056 & Danh sách báo cáo portal & \no & ch.5 & Vẫn chỉ có một báo cáo. \\
CT-057 & Dữ liệu báo cáo theo bộ lọc & \ok & ch.5 & Thẻ \emph{Đơn hoàn thành} \textbf{luôn bằng 0}; bộ lọc ngày \textbf{không quy đổi múi giờ}; biểu đồ luôn nới ra 12 tháng. \\
CT-058 & Export báo cáo & \ok & ch.5 & Tệp Excel \textbf{chứa cả đơn đã huỷ}, màn hình thì không --- hai bên không khớp. \\
CT-059 & Danh sách nhân viên cửa hàng & \ok & ch.4 & Sắp xếp theo vai trò đang là thứ tự chữ cái của mã trạng thái. \\
CT-060 & Tạo/cập nhật nhân viên & \pt & ch.7 & Vẫn đi qua yêu cầu cập nhật thông tin. Thêm hai điều: \textbf{duyệt yêu cầu KHÔNG tự sửa hồ sơ cửa hàng}, và UAT có \uat{0} yêu cầu --- phân hệ chưa từng dùng. \\
CT-061 & Danh sách ca làm & \no & --- & Chưa có dòng mã nào. \\
CT-062 & Đăng ký lịch / ca làm & \no & --- & Chưa có dòng mã nào. \\
CT-063 & Lấy lịch làm việc & \no & --- & Chưa có dòng mã nào. \\
CT-064 & Ghi nhận chấm công & \no & ch.9 & Chưa có trên portal. Có điểm danh \textbf{trong phiếu khảo sát} --- khác nghiệp vụ, và thao tác đó \textbf{sửa được hồ sơ nhân sự cửa hàng}. \\
CT-065 & Tạo yêu cầu nghỉ phép & \no & --- & Chưa có dòng mã nào. \\
CT-066 & Ghi nhận chi phí phát sinh & \no & --- & Chưa có dòng mã nào. \\
CT-067 & Lịch sử chi phí phát sinh & \no & --- & Chưa có dòng mã nào. \\
CT-068 & Upload file / ảnh & \pt & ch.6/7/9 & \textbf{Đính chính}: hiện là \textbf{ba mức} --- Đổi trả đọc nội dung thật của tệp; Hỗ trợ và Yêu cầu thông tin dùng chung một hàm, chỉ tin lời khai của trình duyệt; Khắc phục khảo sát \textbf{không kiểm gì}. \\
CT-069 & Lấy cấu hình portal & \pt & ch.4 & Vẫn vậy. Portal đang bật \uat{3} ngôn ngữ. \\
CT-070 & Kiểm tra quyền menu & \pt & ch.4 & Vẫn vậy. Thêm: \textbf{ba phân hệ chặn vai trò theo ba cách khác nhau} --- xem mục \ref{sec:vai-tro}. \\
CT-071 & Trả thông báo lỗi chuẩn & \pt & mọi ch. & Vẫn chưa thống nhất. Cụ thể: \textbf{sáu màn} chặn ngày ngược và báo tại thanh lọc, \textbf{một màn} lặng lẽ đảo ngày, \textbf{một màn} trả về rỗng không nói gì. \\
\end{longtable}}

\textbf{Tổng kết: \ok{} 56 mục · \pt{} 6 mục · \no{} 9 mục} --- giữ nguyên như bản 19/09; phần
soi sâu không làm mục nào đổi trạng thái, chỉ nói rõ hơn nội dung.

\section{Code đã có nhưng tab CT chưa mô tả}

Phần ngược lại: mã nguồn có, spec chưa có mục nào. BA cân nhắc bổ sung vào tab
\textit{3. Controller} để sau này không ai tưởng là thiếu.

{\footnotesize
\begin{longtable}{@{}L{5.6cm}L{9.4cm}@{}}
\rowcolor{hdrbg}\textbf{Đường dẫn / nhóm} & \textbf{Vì sao nên có mục CT}\\ \midrule \endfirsthead
\rowcolor{hdrbg}\textbf{Đường dẫn / nhóm} & \textbf{Vì sao nên có mục CT}\\ \midrule \endhead
\rt{/portal/forgot-pass} · \rt{/portal/reset_password} & Quên và đặt lại mật khẩu --- nghiệp vụ thật, 1.500 tài khoản chắc chắn cần. \\
\rt{/portal/set-lang/<mã>} & Đổi ngôn ngữ; portal đang bật 3 ngôn ngữ. \\
\rt{/portal/purchase-history/<id>.pdf} & Tải PDF \textbf{đơn hàng} --- khác CT-053 là PDF công nợ. Đang dùng mẫu in mặc định của Odoo. \\
\rt{/portal/delivery/<id>.ics} & Tải lịch giao hàng vào ứng dụng lịch. \\
6 nhóm \rt{.../results} & Lọc không tải lại trang --- quyết định kỹ thuật nhưng đổi cách viết kịch bản kiểm thử. \\
\rt{/portal/order/cart/fragment} & Đồng bộ giỏ khi hai người cùng cửa hàng sửa cùng lúc. \\
\rt{/portal/exam/line/<id>/photo} · \rt{/portal/profile/avatar/<id>} & Hai đường dẫn ảnh có kiểm quyền riêng. \\
\rt{/portal/knowledge/search} & Gợi ý tìm kiếm khi đang gõ. \\
\rt{/portal/info-request/franchise/<id>/values} & Lấy giá trị hiện tại của cửa hàng khi soạn yêu cầu --- xem điểm cần chốt số \textbf{Q58}. \\
9 chuyển hướng 301 & Đường dẫn cũ thời v14; tài liệu cũ của BA ghi URL cũ vẫn vào được. \\
15 đường dẫn \textbf{Khảo sát cửa hàng} & Cả phân hệ khảo sát (8 phía cửa hàng + 7 phía người đi khảo sát) không có mục CT nào. \\
2 đường dẫn \textbf{Metabase} & Nhúng bảng số liệu BI. \\
\end{longtable}}

% =============================================================================
\section{Bốn cách hiểu ``cửa hàng đang xem''}
\label{sec:pham-vi}

Đây là quyết định lớn nhất của tài liệu này. Cùng một người, cùng một lúc đăng nhập, nhưng bảy
phân hệ hiểu ``cửa hàng của tôi'' theo \textbf{bốn cách khác nhau}:

{\footnotesize
\begin{longtable}{@{}L{4.6cm}L{4.2cm}L{6.2cm}@{}}
\rowcolor{hdrbg}\textbf{Cách hiểu} & \textbf{Phân hệ dùng} & \textbf{Hệ quả người dùng thấy}\\
\midrule \endfirsthead
\rowcolor{hdrbg}\textbf{Cách hiểu} & \textbf{Phân hệ dùng} & \textbf{Hệ quả người dùng thấy}\\
\midrule \endhead
\textbf{Bắt buộc chọn một cửa hàng} --- chưa chọn thì màn rỗng &
  Lịch sử mua hàng · Công nợ · Đăng ký thi &
  Người phụ trách nhiều cửa hàng phải bấm đổi cửa hàng để xem từng nơi. Đăng ký thi còn báo nhầm nguyên nhân (``Chưa có đăng ký thi''). \\ \addlinespace
\textbf{Cộng mọi cửa hàng khi chưa chọn} &
  Trang chủ · Giao hàng · Báo cáo &
  Chưa chọn cửa hàng thì các con số là \textbf{tổng toàn bộ} --- dễ tưởng là số của một cửa hàng. \\ \addlinespace
\textbf{Luôn gộp mọi cửa hàng truy cập được} &
  Đổi trả · Yêu cầu thông tin · Khảo sát &
  Danh sách trộn nhiều cửa hàng mà \textbf{không có cột cửa hàng, không có bộ lọc cửa hàng}. \\ \addlinespace
\textbf{Không phân biệt cửa hàng} &
  Kiến thức · Thông báo toàn hệ · Metabase &
  Mọi tài khoản thấy cùng một nội dung. \\
\end{longtable}}

\begin{ask}
\textbf{Đây là một quyết định, không phải bốn.} Anh chốt giúp: portal nên theo cách nào làm
chuẩn, và những màn nào được phép làm khác chuẩn (kèm lý do). Chốt xong dev sẽ đưa tất cả về
một luật --- \textbf{không} cần BA lên task cho từng màn.
\end{ask}

% =============================================================================
\section{Ba cách chặn vai trò khác nhau}
\label{sec:vai-tro}

{\footnotesize
\begin{longtable}{@{}L{4.0cm}L{4.4cm}L{6.6cm}@{}}
\rowcolor{hdrbg}\textbf{Phân hệ} & \textbf{Xét vai trò ở đâu} & \textbf{Không đạt thì thấy gì}\\
\midrule \endfirsthead
\rowcolor{hdrbg}\textbf{Phân hệ} & \textbf{Xét vai trò ở đâu} & \textbf{Không đạt thì thấy gì}\\
\midrule \endhead
Công nợ & Vai trò \textbf{tại cửa hàng đang chọn} & Trang thông báo thân thiện. Chưa chọn cửa hàng thì \textbf{cổng không chạy}. \\
Báo cáo đặt hàng & Vai trò \textbf{cao nhất trên mọi cửa hàng} & Trang thông báo. Dữ liệu vẫn chỉ của cửa hàng đang chọn --- hai phạm vi lệch nhau. \\
Yêu cầu thông tin (màn tạo) & Vai trò \textbf{cao nhất trên mọi cửa hàng} & \textbf{Trang lỗi kỹ thuật 403} của máy chủ. \\
Mọi phân hệ còn lại & \textbf{Không chặn vai trò} & Nhân viên đặt hàng, đăng ký thi, gửi phiếu hỗ trợ, xem danh sách đồng nghiệp kèm điện thoại. \\
\end{longtable}}

\begin{ask}
Cũng là \textbf{một} quyết định: chốt bảng ``vai trò nào làm được gì'' cho cả portal, dev áp
một lần. Riêng \textbf{trang lỗi 403 kỹ thuật} thì dev sửa luôn, BA không cần task.
\end{ask}

% =============================================================================
\section{Toàn bộ điểm cần BA xác nhận}

Bảng dưới gom \textbf{mọi} điểm ghi ở mục 7 của từng màn. Cột \emph{Nhóm}: \textbf{PV} phạm vi
cửa hàng · \textbf{VT} vai trò · \textbf{NC} ngõ cụt quy trình · \textbf{SL} số liệu lệch ·
\textbf{DL} thiếu dữ liệu UAT để nghiệm thu · \textbf{TB} không ai được báo ·
\textbf{CH} chuẩn hoá (dev tự làm, \textbf{BA không cần task}).

{\footnotesize
\begin{longtable}{@{}L{0.8cm}L{1.1cm}L{3.4cm}L{9.4cm}@{}}
\rowcolor{hdrbg}\textbf{\#} & \textbf{Nhóm} & \textbf{Màn} & \textbf{Điểm cần chốt}\\
\midrule \endfirsthead
\rowcolor{hdrbg}\textbf{\#} & \textbf{Nhóm} & \textbf{Màn} & \textbf{Điểm cần chốt}\\
\midrule \endhead
\multicolumn{4}{@{}l}{\textbf{Chương 3 --- Đặt hàng}}\\
Q1 & SL & Danh sách sản phẩm & Sản phẩm chưa gán danh mục vẫn hiện ở ``Tất cả'' (UAT 2/5). \\
Q2 & PV & Danh sách sản phẩm & Chưa chọn cửa hàng thì giá hiện là giá niêm yết, không phải giá cửa hàng. \\
Q3 & CH & Danh sách sản phẩm & Ô tìm kiếm chưa tìm theo tên tiếng Trung dù màn có hiện. \\
Q4 & DL & Danh sách sản phẩm & UAT chưa sản phẩm nào đặt trần tối đa --- chưa nghiệm thu được luật trần. \\
Q5 & SL & Chi tiết sản phẩm & ``Sản phẩm liên quan'' dựa hoàn toàn vào danh mục portal. \\
Q6 & NC & Giỏ hàng & Ghi chú là của cửa hàng, không của người gửi. \\
Q7 & NC & Giỏ hàng & Giỏ không có lịch sử --- hai người sửa chồng nhau thì người sau đè người trước. \\
Q8 & VT & Gửi đơn hàng & Nhân viên gửi được đơn --- có cần chặn không. \\
Q9 & NC & Gửi đơn hàng & Gửi đơn lần hai trong ngày huỷ đơn nháp cũ một cách im lặng. \\
Q10 & --- & Gửi đơn hàng & Đơn không tự xác nhận --- xác nhận đúng thiết kế. \\
Q11 & CH & Đặt hàng thành công & Máy tính và điện thoại kết thúc ở hai màn khác nhau. \\
Q12 & DL & Khung giờ đặt hàng & Tên khung giờ trên UAT không khớp giờ thật. \\
Q13 & PV & Khung giờ đặt hàng & Giờ tính theo múi giờ tài khoản, không theo khu vực cửa hàng. \\
Q14 & DL & Khung giờ đặt hàng & Khu vực Hà Nội (3/4 cửa hàng) chưa có khung riêng. \\
\midrule
\multicolumn{4}{@{}l}{\textbf{Chương 4 --- Khung portal · Trang chủ}}\\
Q15 & CH & Đăng nhập & Chưa chặn thử mật khẩu hàng loạt ở màn đăng nhập. \\
Q16 & NC & Đăng nhập & Đăng xuất không xoá cửa hàng đang chọn --- máy dùng chung giữ lựa chọn cho người sau. \\
Q17 & CH & Quên mật khẩu & Luôn báo thành công nên người gõ nhầm email ngồi chờ email không tới. \\
Q18 & --- & Đặt lại mật khẩu & Luật mật khẩu chỉ ``từ 8 ký tự'' --- có cần thêm yêu cầu không. \\
Q19 & SL & Trang chủ & Số ``Thông báo chưa đọc'' ở Trang chủ lệch với màn Thông báo. \\
Q20 & SL & Trang chủ & Thẻ ``Đổi trả đang mở'' bỏ sót trạng thái \emph{Đang xét}. \\
Q21 & SL & Trang chủ & Thẻ số và danh sách bên dưới dùng hai bộ trạng thái khác nhau. \\
Q22 & CH & Trang chủ & Mọi khối danh sách cố định 2 dòng. \\
Q23 & PV & Trang chủ & Chưa chọn cửa hàng thì các con số là tổng mọi cửa hàng. \\
Q24 & PV & Trang chủ & Khối ``Bài kiến thức'' không lọc theo cửa hàng. \\
Q25 & CH & Chọn cửa hàng & Chọn cửa hàng không được phép thì im lặng, không báo. \\
Q26 & NC & Chọn cửa hàng & Cửa hàng đang chọn lưu trên trình duyệt --- máy khác phải chọn lại. \\
Q27 & NC & Hồ sơ cá nhân & Sửa email ở Hồ sơ không đổi email đăng nhập. \\
Q28 & VT & Hồ sơ cá nhân & Ai cùng cửa hàng đều xem được ảnh đại diện của nhau. \\
Q29 & CH & Đổi mật khẩu & Đổi mật khẩu không đăng xuất thiết bị khác. \\
Q30 & CH & Đổi mật khẩu & Luật mật khẩu lệch nhẹ với màn Đặt lại mật khẩu. \\
Q31 & DL & Đổi ngôn ngữ & Portal đang bật 3 ngôn ngữ --- xác nhận có đủ bản dịch. \\
Q32 & CH & Danh sách cửa hàng & Hai bộ màn song song cho cùng một nội dung. \\
Q33 & VT & Chi tiết cửa hàng & Nhân viên xem được đủ danh sách đồng nghiệp kèm điện thoại. \\
Q34 & CH & Chi tiết cửa hàng & Bốn đường dẫn cho ba màn gần giống nhau. \\
Q35 & CH & Thông tin cửa hàng & Sắp xếp thành viên theo thứ tự chữ cái của mã trạng thái. \\
Q36 & DL & Thông tin cửa hàng & Nhánh ``cửa hàng bị khoá'' chưa có dữ liệu để nghiệm thu. \\
\midrule
\multicolumn{4}{@{}l}{\textbf{Chương 5 --- Lịch sử mua · Giao hàng · Báo cáo}}\\
Q37 & PV & Lịch sử mua hàng & Ba phân hệ hiểu ``cửa hàng đang xem'' khác nhau (xem mục \ref{sec:pham-vi}). \\
Q38 & NC & Lịch sử mua hàng & Đơn Đã huỷ bị giấu hoàn toàn (UAT 16/47). \\
Q39 & CH & Lịch sử mua hàng & Tìm kiếm chỉ theo mã đơn, không theo tên sản phẩm. \\
Q40 & VT & Lịch sử mua hàng & Nhân viên xem được toàn bộ số tiền các đơn. \\
Q41 & NC & Chi tiết đơn hàng & Không hiện lịch sử thay đổi của đơn. \\
Q42 & CH & PDF đơn hàng & Dùng mẫu in mặc định của Odoo, chưa có nhận diện WujiaTea. \\
Q43 & DL & Theo dõi giao hàng & Cửa hàng đông đơn nhất (Cầu Giấy, 8 đơn) không có chuyến nào. \\
Q44 & SL & Theo dõi giao hàng & Chuyến chưa đặt lịch biến mất khi lọc theo ngày. \\
Q45 & CH & Theo dõi giao hàng & Không lọc riêng được chuyến bị huỷ. \\
Q46 & PV & Chi tiết chuyến giao & Điều kiện xem chi tiết hẹp hơn điều kiện hiện ở danh sách. \\
Q47 & SL & Chi tiết chuyến giao & Hai bước cuối thanh tiến trình mượn ``giờ sửa bản ghi'', chưa có giờ thật. \\
Q48 & SL & Tệp lịch giao hàng & Giờ trong tệp lịch lấy từ trường khác với giờ hiển thị trên màn. \\
Q49 & SL & Báo cáo đặt hàng & Thẻ ``Đơn hoàn thành'' luôn bằng 0 (Odoo 19 không còn trạng thái đó). \\
Q50 & SL & Báo cáo đặt hàng & Số tiền dán sai ký hiệu tiền tệ. \\
Q51 & SL & Báo cáo đặt hàng & Bộ lọc ngày không quy đổi múi giờ, khác Lịch sử mua hàng. \\
Q52 & VT & Báo cáo đặt hàng & Vai trò xét trên mọi cửa hàng nhưng dữ liệu chỉ của cửa hàng đang chọn. \\
Q53 & SL & Báo cáo đặt hàng & Biểu đồ luôn nới ra 12 tháng kể cả khi lọc một tuần. \\
Q54 & SL & Excel báo cáo & Tệp Excel chứa cả đơn đã huỷ, màn hình thì không. \\
Q55 & SL & Excel báo cáo & Cột \emph{Số dòng} đếm cả dòng tiêu đề, màn hình thì không. \\
Q56 & SL & Excel báo cáo & Cột \emph{Tổng tiền} là số trần, không kèm loại tiền. \\
Q57 & VT & Excel báo cáo & Điều kiện vai trò của nút Xuất khác của màn. \\
\midrule
\multicolumn{4}{@{}l}{\textbf{Chương 6 --- Công nợ · Đổi trả}}\\
Q58 & SL & Công nợ theo tuần & Số ở Trang chủ và số ở màn Công nợ là hai phép tính khác nhau. \\
Q59 & NC & Công nợ theo tuần & Dropdown chỉ 6 tuần --- hoá đơn cũ hơn vĩnh viễn không chọn được. \\
Q60 & SL & Công nợ theo tuần & ``Dư có'' đè lên ``Có quá hạn'', làm mất nút thanh toán. \\
Q61 & SL & Công nợ theo tuần & Trạng thái rỗng nói ``không còn công nợ'' dù có hoá đơn quá hạn. \\
Q62 & SL & Công nợ theo tuần & Hạn thanh toán do portal tự tính (thứ Hai + 10), không đọc hạn của hoá đơn. \\
Q63 & --- & Công nợ theo tuần & Tuần mở sẵn = tuần quá hạn cũ nhất còn dư --- xác nhận cách mở. \\
Q64 & VT & Công nợ theo tuần & Nhân viên bị chặn khỏi toàn bộ phân hệ Công nợ. \\
Q65 & DL & Lịch sử thanh toán & Kỳ mở sẵn là tháng hiện tại, trên UAT tháng đó rỗng. \\
Q66 & CH & Lịch sử thanh toán & Ngày ngược thì tự hoán đổi thay vì báo lỗi --- lệch 4 màn khác. \\
Q67 & SL & Lịch sử thanh toán & Tổng tách theo loại tiền, không quy đổi. \\
Q68 & DL & Lịch sử thanh toán & 1/6 phiếu thanh toán chưa gắn cửa hàng nên không ai thấy. \\
Q69 & --- & Chuyển khoản & Ảnh mã QR: tĩnh hay động theo số tiền. \\
Q70 & SL & Chuyển khoản & Nội dung chuyển khoản làm tròn xuống số nguyên. \\
Q71 & TB & Chuyển khoản & Chuyển khoản xong portal không biết gì --- đợi kế toán ghi nhận. \\
Q72 & PV & DS đổi trả & Gộp mọi cửa hàng, không có cột và bộ lọc cửa hàng. \\
Q73 & NC & DS đổi trả & Nhãn ``Nháp'' có trong bộ lọc dù phiếu nháp không gửi được. \\
Q74 & DL & DS đổi trả & 10/14 phiếu trên UAT là dữ liệu mẫu, thiếu đơn gốc và sản phẩm. \\
Q75 & DL & Tạo yêu cầu đổi trả & Cửa sổ 10 ngày làm màn này \textbf{không dùng được} --- UAT 0 đơn hợp lệ. \\
Q76 & CH & Tạo yêu cầu đổi trả & Ô chọn đơn không lọc theo cửa hàng ở ô bên cạnh --- báo lỗi gây hiểu nhầm. \\
Q77 & NC & Tạo yêu cầu đổi trả & Phiếu nháp là ngõ cụt --- lưu được nhưng không đường dẫn nào gửi được. \\
Q78 & CH & Tạo yêu cầu đổi trả & Điện thoại không có nút Lưu nháp --- hai nền tảng khác luật. \\
Q79 & NC & Tạo yêu cầu đổi trả & Sản phẩm chưa cấu hình bù thì chặn cả đổi và trả. \\
Q80 & --- & Tạo yêu cầu đổi trả & Mỗi phiếu đúng một sản phẩm. \\
Q81 & NC & Tạo yêu cầu đổi trả & Video tải lên không bao giờ hiện lại cho cửa hàng. \\
Q82 & NC & Tạo yêu cầu đổi trả & Không sửa, không huỷ được phiếu sau khi gửi. \\
Q83 & NC & Chi tiết yêu cầu & Màn chi tiết không có nút nào cho cửa hàng. \\
Q84 & DL & Chi tiết yêu cầu & Thanh tiến độ bù chưa chạy được trên UAT --- chưa có bản ghi phân bổ. \\
Q85 & SL & Chi tiết yêu cầu & Thanh phần trăm bị chặn trần 100\% --- bù dư vẫn hiện 100\%. \\
\midrule
\multicolumn{4}{@{}l}{\textbf{Chương 7 --- Hỗ trợ · Yêu cầu thông tin · Kiến thức · Thông báo}}\\
Q86 & PV & DS phiếu hỗ trợ & Phiếu chỉ người tạo mới thấy --- Chủ tiệm không thấy phiếu nhân viên. \\
Q87 & CH & DS phiếu hỗ trợ & Không có bộ lọc ngày, cỡ trang cố định 20. \\
Q88 & NC & DS phiếu hỗ trợ & Phiếu Đã huỷ bị giấu ở danh sách nhưng chi tiết vẫn mở được. \\
Q89 & CH & Tạo phiếu hỗ trợ & Định dạng tệp chỉ kiểm theo lời khai của trình duyệt. \\
Q90 & TB & Tạo phiếu hỗ trợ & Gửi phiếu xong không ai được báo. \\
Q91 & CH & Tạo phiếu hỗ trợ & Chưa giới hạn số phiếu gửi trong ngày. \\
Q92 & --- & Chi tiết phiếu hỗ trợ & Ranh giới nội bộ / công khai nằm ở ô nhân sự Ngô Gia chọn. \\
Q93 & NC & Chi tiết phiếu hỗ trợ & Phiếu đã đóng / đã huỷ vẫn nhận lời nhắn mới. \\
Q94 & NC & Gửi lời nhắn & Trả lời không tự đổi trạng thái phiếu. \\
Q95 & NC & Gửi lời nhắn & Không đính kèm được ảnh trong lời nhắn. \\
Q96 & VT & DS yêu cầu thông tin & Nhân viên xem được mọi yêu cầu của cửa hàng nhưng không tạo được. \\
Q97 & NC & DS yêu cầu thông tin & Không có nhãn ``Đã huỷ'' --- tự huỷ hiện thành ``Từ chối''. \\
Q98 & --- & DS yêu cầu thông tin & Bảy loại thông tin cố định trong mã nguồn. \\
Q99 & VT & Tạo yêu cầu thông tin & Quyền xét trên mọi cửa hàng nhưng gửi được cho từng cửa hàng. \\
Q100 & CH & Tạo yêu cầu thông tin & Nhân viên bị chặn bằng trang lỗi kỹ thuật 403. \\
Q101 & CH & Tạo yêu cầu thông tin & Loại ``Khác'' bắt người dùng gõ tên kỹ thuật của trường. \\
Q102 & NC & Tạo yêu cầu thông tin & Duyệt yêu cầu \textbf{không} tự cập nhật hồ sơ cửa hàng. \\
Q103 & NC & Tạo yêu cầu thông tin & Phiếu nháp ở đây gửi được, khác Đổi trả --- cần thống nhất. \\
Q104 & SL & Chi tiết yêu cầu & ``Giá trị cũ'' không phải ảnh chụp --- đọc lại mỗi lần mở màn. \\
Q105 & CH & Huỷ yêu cầu & Huỷ và bị từ chối trông giống nhau. \\
Q106 & CH & Lấy giá trị cửa hàng & Đường dẫn đọc được bất kỳ trường nào của hồ sơ cửa hàng, không kiểm vai trò. \\
Q107 & VT & DS kiến thức & Kiến thức không phân biệt cửa hàng và vai trò. \\
Q108 & DL & DS kiến thức & Danh mục trùng tên trên UAT (Vận hành ×2, Sản phẩm ×2). \\
Q109 & SL & DS kiến thức & Bài hết hạn chỉ rụng khi tác vụ nền chạy hằng ngày. \\
Q110 & SL & DS kiến thức & Lượt xem đếm mỗi lần mở, không chống bấm lặp. \\
Q111 & NC & Chi tiết bài viết & Không biết ai đã đọc bài nào. \\
Q112 & --- & Chi tiết bài viết & Đường dẫn dùng tên rút gọn; đổi tiêu đề không đổi tên rút gọn. \\
Q113 & CH & Gợi ý tìm kiếm & Chưa giới hạn số lần gọi --- cần theo dõi với 1.500 tài khoản. \\
Q114 & PV & DS thông báo & ``Đã đọc'' tính theo từng cửa hàng --- người 3 cửa hàng phải đọc 3 lần. \\
Q115 & DL & DS thông báo & 6 dấu đã đọc cũ chưa gắn cửa hàng. \\
Q116 & SL & DS thông báo & Thẻ đếm chưa đọc ở Trang chủ tính khác màn này (trùng Q19). \\
Q117 & DL & DS thông báo & UAT chưa có thông báo gửi đích danh cửa hàng nào đang phát hành. \\
Q118 & NC & Chi tiết thông báo & Mở khi chưa chọn cửa hàng thì không được tính là đã đọc. \\
Q119 & CH & Hộp chuông & Cố định 5 dòng, không có ``xem thêm''. \\
Q120 & CH & Đánh dấu đã đọc & ``Đánh dấu tất cả'' chỉ là tất cả tại cửa hàng đang chọn, còn hiệu lực. \\
\midrule
\multicolumn{4}{@{}l}{\textbf{Chương 8 --- Đăng ký thi}}\\
Q121 & CH & DS đăng ký thi & Chưa chọn cửa hàng nhưng màn báo ``Chưa có đăng ký thi''. \\
Q122 & CH & DS đăng ký thi & Cùng trạng thái, hai chữ khác nhau giữa máy tính và điện thoại. \\
Q123 & SL & DS đăng ký thi & Huy hiệu ``Có kết quả'' đè lên trạng thái phiếu trên bản điện thoại. \\
Q124 & PV & DS đăng ký thi & Chỉ xem được cửa hàng đang chọn dù luật dữ liệu cho phép xem mọi cửa hàng. \\
Q125 & DL & Màn đăng ký thi & UAT không có kỳ thi nào đăng ký được. \\
Q126 & CH & Màn đăng ký thi & Chưa chọn cửa hàng thì bị đá về danh sách không lời giải thích. \\
Q127 & --- & Màn đăng ký thi & ``Cửa sổ 60 ngày'' đặt riêng theo từng khóa thi. \\
Q128 & --- & Khung giờ thi & Kỳ thi đặt giới hạn người riêng, đè lên giới hạn của khóa. \\
Q129 & NC & Gửi phiếu đăng ký & Phiếu chờ duyệt đã chiếm ghế của cửa hàng khác. \\
Q130 & NC & Gửi phiếu đăng ký & Người dự thi gõ tay, không nối danh sách nhân sự cửa hàng. \\
Q131 & TB & Gửi phiếu đăng ký & Gửi phiếu xong không ai được báo. \\
Q132 & CH & Gửi phiếu đăng ký & Bản điện thoại không có bước xem lại trước khi gửi. \\
Q133 & NC & Chi tiết phiếu thi & Portal không huỷ được phiếu. \\
Q134 & NC & Chi tiết phiếu thi & Phiếu bị từ chối là ngõ cụt --- phải tạo lại từ đầu. \\
Q135 & PV & Chi tiết phiếu thi & Chỉ mở được phiếu của cửa hàng đang chọn. \\
\midrule
\multicolumn{4}{@{}l}{\textbf{Chương 9 --- Khảo sát · Metabase} \note{(mã anh Thái --- dev trao đổi trực tiếp)}}\\
Q136 & CH & DS khảo sát & Ngày ngược thì hiện rỗng, không báo gì. \\
Q137 & CH & DS khảo sát & Cỡ trang cố định 5. \\
Q138 & --- & DS khảo sát & Phiếu còn nháp portal không thấy. \\
Q139 & SL & Chi tiết khảo sát & Điểm trên portal và điểm trong ERP có thể lệch nhau. \\
Q140 & SL & Chi tiết khảo sát & Tổng điểm tối đa ghi cứng 100. \\
Q141 & CH & Chi tiết khảo sát & ``Vi phạm nghiêm trọng'' nhận biết bằng cách dò chữ. \\
Q142 & CH & Chi tiết khảo sát & Còn vài giá trị mẫu sót lại khi dữ liệu thiếu. \\
Q143 & SL & Màn gửi khắc phục & Mốc thời gian là giờ tạo phiếu, không phải ngày khảo sát. \\
Q144 & CH & Gửi khắc phục & Đường dẫn gửi khắc phục không kiểm cửa hàng. \\
Q145 & CH & Gửi khắc phục & Ảnh khắc phục không kiểm định dạng, không kiểm dung lượng. \\
Q146 & TB & Gửi khắc phục & Gửi xong phiếu vẫn ``Chờ phản hồi'', không ai được báo. \\
Q147 & NC & Gửi khắc phục & Gửi lại đè lên lần trước, không lưu lịch sử. \\
Q148 & CH & Màn làm khảo sát & Phiếu chưa giao cho ai thì tài khoản nào cũng mở được. \\
Q149 & CH & Lưu bài khảo sát & Đường dẫn lưu không kiểm người gửi. \\
Q150 & NC & Lưu bài khảo sát & Điểm danh tự tạo thành viên cửa hàng. \\
Q151 & NC & Lưu bài khảo sát & Phiếu tự chuyển sang ``Chờ phản hồi'' ngay khi lưu tạm. \\
Q152 & NC & Doanh thu PosApp & Bấm đồng bộ là xoá sạch số đã gõ tay, không hỏi lại. \\
Q153 & CH & Doanh thu PosApp & Lỗi kết nối hiện nguyên văn câu lỗi kỹ thuật. \\
Q154 & VT & Bảng điểm danh & Màn khảo sát sửa được hồ sơ nhân sự của cửa hàng. \\
Q155 & NC & Bảng điểm danh & Xoá dòng điểm danh không gỡ thành viên đã tạo. \\
Q156 & VT & Metabase & Bảng không khai nhóm quyền thì ai đăng nhập cũng xem được. \\
Q157 & PV & Metabase & Không truyền cửa hàng sang Metabase --- người xem thấy số liệu toàn hệ. \\
\end{longtable}}

\textbf{Tổng: 157 điểm.} Trong đó \textbf{42 điểm nhóm CH} là việc dev tự chuẩn hoá ---
\textbf{BA không cần lên task}. Còn lại 115 điểm cần anh chốt, nhưng \textbf{không phải 115
quyết định}: nhóm \textbf{PV} (14 điểm) gộp về \textbf{một} quyết định ở mục \ref{sec:pham-vi},
nhóm \textbf{VT} (11 điểm) gộp về \textbf{một} quyết định ở mục \ref{sec:vai-tro}.

% =============================================================================
\section{BA đừng lên task viết lại --- ba nhóm}

\subsection*{1. Nhóm CH --- dev tự chuẩn hoá}

42 điểm đánh nhóm \textbf{CH} ở bảng trên đã nằm trong kế hoạch chuẩn hoá nội bộ của dev (cụm
F): thống nhất thông báo lỗi, thống nhất bộ lọc ngày, thống nhất cỡ trang, vá các chỗ kiểm
thiếu khi tải tệp và khi ghi dữ liệu. BA \textbf{không cần} lên task cho từng cái --- lên task
chỉ khiến hai bên làm trùng.

\subsection*{2. Controller sắp đổi chỗ --- đường dẫn giữ nguyên}

Kế hoạch F6--F13 sẽ \textbf{dời mã nguồn} của bảy phân hệ sang mô-đun khác cho đúng kiến trúc
tầng (ADR-027). \textbf{Đường dẫn không đổi một ký tự nào}, màn không đổi, dữ liệu không đổi.
Nếu BA thấy tên mô-đun trong tài liệu này khác với tên trong bản ghi chép sau đó, đó là lý do
--- \textbf{không phải lỗi, không cần task}.

\subsection*{3. Mã của anh Thái}

Chương 9 (15 đường dẫn khảo sát + 2 đường dẫn Metabase) do anh Thái viết. Những điểm dev thấy
cần bàn --- nhất là Q144, Q145, Q148, Q149, Q156 --- dev sẽ trao đổi trực tiếp với anh Thái.
BA chỉ cần biết để trả lời khi cửa hàng hỏi.

% =============================================================================
\section{Phân hệ chưa có gì: vận hành nội bộ cửa hàng}

Bảy mục \textbf{CT-061 … CT-067} (ca làm, lịch làm việc, chấm công, nghỉ phép, chi phí phát
sinh) \textbf{chưa có một dòng mã nào}. BA đánh ưu tiên Thấp cho cả bảy. Nếu muốn làm thì đây
là \textbf{một phân hệ mới}, không phải sửa cái có sẵn --- ước lượng và kế hoạch phải tính
riêng.

\section{Dữ liệu UAT cần dựng trước khi nghiệm thu}

Mười lăm điểm nhóm \textbf{DL} ở bảng trên có chung một nguyên nhân: \textbf{UAT thiếu dữ liệu
mẫu}, nên có những luật đã viết xong mà chưa ai bấm thử được. Ba chỗ nặng nhất:

\begin{enumerate}[leftmargin=*]
  \item \textbf{Đổi trả · Bù hàng} --- \uat{0} đơn hàng lọt cửa sổ 10 ngày. Không tài khoản nào
        tạo được yêu cầu. Cần một đơn đặt trong vòng 10 ngày gần nhất.
  \item \textbf{Đăng ký thi} --- \uat{0} kỳ thi đăng ký được (hai kỳ đang mở đều đã qua ngày
        thi 16/08 và 25/08). Cần một kỳ thi ngày tương lai, còn ghế.
  \item \textbf{Hỗ trợ} --- \uat{0} phiếu do tài khoản cửa hàng thật tạo; cả \uat{15} phiếu đều
        của tài khoản nội bộ, nên mọi tài khoản cửa hàng mở ra màn rỗng.
\end{enumerate}

Ngoài ra: \uat{0} yêu cầu cập nhật thông tin từng được tạo; \uat{0} thông báo gửi đích danh cửa
hàng đang phát hành; \uat{0} bản ghi phân bổ bù hàng; chỉ \uat{1}/4 cửa hàng có dữ liệu kế
toán để xem Công nợ.


\end{document}
